\documentclass{article}
% >>> COMPILE WITH XeLaTeX (Overleaf: Menu -> Settings -> Compiler -> XeLaTeX). <<<
\usepackage{graphicx} % Required for inserting images
\usepackage{xcolor} % Must be loaded before soul
\usepackage{soul} % For highlighting text
% inputenc/fontenc intentionally omitted: XeTeX is natively UTF-8 and fontspec provides the encoding.
\usepackage{amssymb} % \therefore used in the metaphysics premise blocks
\usepackage{newunicodechar}
\usepackage{tikz}
\usetikzlibrary{positioning}
\usepackage{fontspec}           % For Greek text support with XeLaTeX/LuaLaTeX
\usepackage{polyglossia}        % Multilingual support
\setmainlanguage{english}
\setotherlanguage{greek}

\usepackage{geometry}
\geometry{margin=1in}

\usepackage{csquotes}           % Block quotes
\usepackage{enumitem}           % Better lists
\usetikzlibrary{arrows.meta, positioning, shapes.geometric}

\usepackage{fancyvrb}           % Verbatim environments for ASCII diagrams
\usepackage{setspace}           % Line spacing

% --- Greek support (XeLaTeX) ---
\newfontfamily\greekfont{Gentium Plus} % good polytonic Greek; try Times New Roman if needed
\newcommand{\gk}[1]{{\greekfont #1}}

% RODA chain glyphs: subscripts A_0..A_4, plus the arrow / logic / therefore symbols used in the text
\newunicodechar{₀}{\ensuremath{_0}}
\newunicodechar{₁}{\ensuremath{_1}}
\newunicodechar{₂}{\ensuremath{_2}}
\newunicodechar{₃}{\ensuremath{_3}}
\newunicodechar{₄}{\ensuremath{_4}}
\newunicodechar{→}{\ensuremath{\rightarrow}}
\newunicodechar{∧}{\ensuremath{\land}}
\newunicodechar{∴}{\ensuremath{\therefore}}
\newunicodechar{′}{\ensuremath{{}^{\prime}}} % prime, e.g. A_0' (next-cycle deposit)

\usepackage[backend=biber,style=authoryear]{biblatex}
\addbibresource{references.bib} %Add References page for in text citation
\usepackage[colorlinks=true, linkcolor=blue, urlcolor=blue, citecolor=blue]{hyperref} %enable hyperlinks to be colored and interactable
\usepackage{changepage} %create individual blocks of text that can be altered apart from general document.

\usepackage{comment}

\newcommand{\ra}{\ensuremath{\rightarrow}\ }

\newcommand{\inlinenote}[1]{\par\noindent
  \fcolorbox{orange}{orange!20}{\parbox{0.97\linewidth}{#1}}\par\noindent}

\begin{document}

\begin{center}
{\LARGE\bfseries Part II — Virtual Worlds as Rhetorical Ecologies:}\\[4pt]
{\large An Ontology, a Method, and Two Designed Worlds}
\end{center}
\vspace{1em}


\begin{quote}\small\itshape This document assembles the five Part II sections worked through to date — the metaphysics of virtual and phantasmic things, the justification of the method (RODA), the introduction to the applications, and the tutorial analyses of two designed worlds at opposite poles of the medium. The metaphysics keeps its own, more periodic register; the remaining sections hold the figured-but-plain voice of Part I. The §2 disruption analysis of Red Dead Redemption 2 --- the Sonny event --- is now included below as Section 6.\end{quote}

\bigskip

\bigskip

\section*{1. What Are Virtual and Phantasmic Things?}

\subsection*{What is A Virtual Thing?}

The conception of the \textit{virtual} and \textit{virtual reality}, as Michael Heim argues, can be traced as far back as Plato and Diotima (Heim, \textit{Virtual Realism}, 18). Philosophically speaking, as Heim observes, ``virtual'' is a term meaning ``not actually, but as if'' (Heim, \textit{Virtual Realism}, 220). Etymologically, ``virtual'' evolved from post-classical Latin’s virtualis: ``of or relating to power or potency,'' that which has the power to produce an effect (``virtual''). Contemporarily speaking, it is used to refer to the virtual environments, objects, persons, and events produced by digital technologies. Considering the wide range of things to which \textit{virtual} can be—and has been—applied, my use of \textit{virtual reality} refers to computer-generated—and mediated—scenes, objects, persons, and events that manifest an interactive virtual environment into which the user is immersed through devices such as Head Mounted Displays (HMD) and handheld controllers.

It is, likely, due to this history and understanding of the virtual and virtual reality technologies that people commonly compare virtual objects to the non-virtual objects they supposedly represent, suggesting that they are not ``real'' or not as ``real'' as the things of objective reality.\footnote{I want to note that I use the term ``objective reality'' to refer to the everyday reality outside of or external to that produced by VR technologies. I use ``non-virtual'' to refer to the objects experienced in everyday, objective reality.} However, I believe this to be a mistaken assumption. Our, often instinctual, responses to virtual objects are an indication of their reality: insofar as a virtual object or thing produces or provokes a response from the user experiencing it is grounds for thinking that a real object is present.

Thus, I occupy a position similar to that taken by Heim and David Chalmers: evinced by their causal powers that are able to affect users, I argue that virtual objects and virtual reality are very much real things, they are just virtual things. However, while I agree with Chalmers and accept that virtual objects are real and not merely fictional or illusory objects, I disagree with the identity claim he makes and which serves as the foundation for his argument: that virtual objects are digital objects. Just as we need to make a distinction between a non-virtual object and a virtual object, so too do we need to make a distinction between a virtual object and the digital object on which it depends.\footnote{I use ``depend'' here and throughout to refer to ontological dependence.} Accordingly, I will argue against Chalmers and suggest that digital objects and virtual objects are two different \textit{kinds} of objects: a digital object is a \textit{functional} object in that it functions to produce a specific virtual object, whereas a virtual object is a \textit{perceptual} object that is the rendered output of the specific digital object on which it depends. Our work, then, is to pin down the nature of a \textit{digital object}, \textit{digital property}, and \textit{digital event} versus a \textit{virtual object}, \textit{virtual property}, and \textit{virtual event}.\footnote{Here I am taking Jaegwon Kim’s notion of an event as an object and a property that is temporally situated.}

Before pinning down these distinctions, I want to fix the ontological register in which they operate. To call a thing \textit{virtual} is, etymologically and ontologically, to ascribe to it the character of a \textit{potency}: it is not a finished, freestanding actuality but something \textit{rendered}—brought from potency into act—by an agency that actualizes it. This is precisely Aristotle's distinction between a first and a second actuality. One who knows English but sits silent possesses that knowledge as a \textit{first actuality}: a developed capacity, fully real, yet unexercised; when one speaks, that knowledge passes into \textit{second actuality}, its active exercise (Aristotle, \textit{De Anima}, 417a21–b2). A digital object, I will argue, stands to the virtual object it generates as first actuality stands to second: the code persists as a structured and entirely real, but unexercised, potency, while the computer that reads and renders it is the agency that actualizes it into the perceptible virtual object. The non-virtual object, by contrast, is already fully actual; it depends on nothing further in order to be what it is. This actualization of a potency is the genus to which both the virtual object and, as I will later argue, the phantasmic object belong: each is a potency rendered into act by its respective renderer, and each therefore exists only in the act of being so rendered.

I believe we can start to get at the distinction between the two kinds of objects by thinking through the following questions: (1) On what does a virtual object and a digital object ontologically depend? (2) What is the difference between a digital property of a digital object versus a virtual property of a virtual object? (3) What is the difference between a digital event and a virtual event?

\subsubsection*{On What Does A Virtual Thing Ontologically Depend?}

A virtual object, like the lightsaber one is able to wield in the VR program \textit{Trials on Tatooine}, consists of various lines of code which determine its visual properties—such as size, shape, color—as well as what kinds of actions it can and cannot perform or be a part of. This segment of code, and the program overall, is then located in the hard drive of a computer that reads the code and renders the virtual object that is mediated through the HMD. The fact that a specific segment of code or data structure constitutes a certain virtual object leads Chalmers to argue that virtual objects are ``digital objects constituted by computational processes on a computer'' (Chalmers, ``The Virtual and the Real,'' 317).

Thus, he understands virtual objects as digital objects that ``can be regarded as data structures, which are grounded in computational processes which are themselves grounded in physical processes on one or more computers'' (Chalmers, ``The Virtual and the Real,'' 317). The perceptible properties experienced when we see a virtual lightsaber ``directly reflects the properties of this data structure'' (Chalmers, ``The Virtual and the Real,'' 317). Moreover, the interaction of two virtual objects, say the lightsaber and the storm trooper’s blaster fire, is generated ``by causal interaction among these data structures'' which are themselves ``causally active on real computers in the real world'' (Chalmers, ``The Virtual and the Real,'' 317). Therefore, Chalmers’ identity claim derives from and is supported by his argument from causal power: ``(1) Virtual objects have certain causal powers (to affect other virtual objects, to affect users, and so on). (2) Digital objects really have those causal powers (and nothing else does). ∴ (3) Virtual objects are digital objects'' (Chalmers, ``The Virtual and the Real,'' 318).

Arguing that virtual objects are digital objects that are \textit{grounded} in computational processes which are \textit{grounded} in physical computers, Chalmers attempts to explain the existence and ontological status of virtual objects in terms of the code, data structures, or digital objects that constitute them which in turn are explained by the computational process of the physical computer. However, while I agree with Chalmers’ characterization of a digital object as a data structure which determines a virtual object’s perceptible properties, there seems to be complications with Chalmers’ identity claim. Firstly, his definition of a digital object could easily be applied to more traditional objects such as the icon for \textit{Microsoft Word} and the program itself would also be qualified as such. Thus, his definition fails to capture the distinction between a virtual object experienced in VR versus other types of digital objects. Secondly, and more important, it seems that a digital object can exist without the virtual object, but a virtual object cannot exist without the digital object on which it depends. This state of affairs is our first clue that virtual objects and digital objects are not identical.

The digital object only depends on the computer which stores and reads it, whereas the virtual object ontologically depends on a particular digital object, the computer that renders it, and the HMD which mediates our experience of it. Even if the computer is turned off the digital object still exists in the computer’s memory, but the rendered virtual object does not. While we might be able to open the program and observe the actual code or digital object as it is stored in the computer, the virtual object which arises from it cannot be perceived without the VR technology that mediates the final rendered output. Virtual objects are real objects insofar as they are images composed of light—a massless object—projected and mediated by the HMD. A specific virtual object, therefore, can only exist when the computer reads the digital object and actively renders it as a virtual object which is projected and mediated through the HMD; the existence of a digital object, on the other hand, needs only the computer in which it is stored. Furthermore, Chalmers’ claim does not seem to adequately capture and explain our experience of virtual objects, and it is in regard to this aspect that his argument does not seem to hold up. When I dodge, say, a storm trooper’s blaster fire headed my way, I am not responding to the data structure or digital object but the perceptual visual and spatial properties of the virtual object the digital object manifests.

\subsubsection*{Virtual Blueness vs. Digital Blueness}

An adept and experienced programmer fluent in the coding language may be able to look at the segment of code and understand that the object will be a lightsaber with such and such properties, but the experience of the virtually rendered object is in no way identical with the experience of observing the digital object on which it depends. Just because a virtual object covaries with\footnote{For you cannot make a change in the virtual object without making a change in the digital object on which it depends.} and depends on the digital object does not necessitate that the two objects—and, especially, our experience of them—are identical. Even though the properties of the virtual object are related to the properties of the digital object, the virtual object has properties that the digital object does not: namely, the perceptible properties the user experiences when encountering a virtual thing. You cannot hear the crackling sound of the lightsaber as it moves through the air or see the slight halo of light that radiates from it like heat waves coming off a hot sidewalk by merely examining the digital object on which it depends. You must experience the virtual object in and of itself.

The fact that my virtual lightsaber is virtually blue is dependent on a certain value input in the underlying data structure which generates it, but that data structure or digital object is not blue in any perceptual sense; for the only thing perceived when observing the actual digital object (code) is a series of numbers, letters, and other various signs. At most one might, as Chalmers does, argue that the data structure of the blue lightsaber is \textit{functionally} blue in that it produces or causes a visual experience of the virtual lightsaber as being blue (Chalmers, ``The Virtual and the Real,'' 321). Thus, we can think about the specific input value or segment of code which \textit{functionally} produces the experience of blueness as a property of the digital object: digital blueness. Because Chalmers’ argument from causal power leads him to identify the virtual object as the digital object which constitutes it, he seems to suggest that the digital object’s property of digital blueness functions as a causal explanation for our experience of the virtual lightsaber’s virtual blueness. But to do so is to conflate two different kinds of properties of two different kinds of objects: the perceptible visual property of the lightsaber’s virtual blueness with the specific value input or digital property (digital blueness) of a digital object that generates or causes that perceptible property.

In regard to our perceptual experience of a virtual object, Chalmers invokes an argument from perception to support his claim.

\begin{quote}
(1) When using virtual reality, we perceive (only) virtual objects. \\
(2) The objects we perceive are the causal basis of our perceptual experiences. \\
(3) When using virtual reality, the causal bases of our perceptual experiences are digital objects. \\
∴ \\
(4) Virtual objects are digital objects. (Chalmers, ``The Virtual and the Real,'' 318)
\end{quote}

While I agree with his first and second premise, his third premise, though, is problematic. As he continues, a virtual object is a digital object because ``many people can see the same digital object by experiencing virtual reality with different displays'' (Chalmers, ``The Virtual and the Real,'' 319), since on his view that digital object is the common causal basis of all of their images.

However, to use his own words, ``a moment's reflection suggests that this cannot be right'' (Chalmers, ``The Virtual and the Real,'' 318-19).

Let us consider if the same digital object (\textit{d}), the data structure of our virtual lightsaber, is mediated through two different HMDs. One of them is working normally and projects the virtual lightsaber in a particular shade of virtual blue (\textit{b1}), while the other has a slight defect or perhaps is a different brand, causing the virtual blueness of the lightsaber to be projected in a different shade of virtual blue (\textit{b2}). If the blueness of the lightsaber is a functional property of a digital object, this implies that there is a relation between inputs and outputs: the same input will always produce the same output. However, in this example the same input produces two different outputs: \textit{d → b₁ ∧ b₂}. So, while we may have the same kind of digital object we have two different instances: \textit{b1} and \textit{b2}; and, if the objects we perceive are the causal basis of our perceptual experiences, then, in this case, there are two distinct perceptual experiences, similar though they may be. Because we have, ostensibly, two distinct perceptual experiences of the same digital object suggests that the digital object is not the causal basis of our perceptual experience. The only thing left, therefore, that can occupy such a causal role is the virtual object.

Because, as I mentioned above, our actual experience of the digital object’s digital blueness (code or data structure) and the virtual object’s virtual blueness (our virtual blue lightsaber) is so drastically different, I believe it is more useful to take a naïve realist stance which suggests that the blueness of the virtual lightsaber is an intrinsic feature or property of the virtual object while the digital object’s digital blueness is a functional property. Therefore:

\begin{quote}
(1) When using virtual reality, we perceive (only) virtual objects. \\
(2) The object we perceive is the causal basis of our perceptual experiences. \\
(3) When using virtual reality, the causal bases of our perceptual experiences are virtual objects \\
(4) Digital objects are the causal basis of the virtual objects we perceive \\
∴ \\
(5) Virtual objects are not digital objects.
\end{quote}

If this is the case, then it follows that we have two entirely different kinds of properties of two different kinds of objects both of which have their own causal power: the digital object’s property of digital blueness functions to cause the veridical blueness of the virtual object that causally explains our experience of the blueness of the virtual lightsaber. Therefore, the relation between the virtual object and our perceptual experience, and the relation between the virtual object and the digital object on which it ontologically depends is \textit{causal} or \textit{explanatory}.

\subsubsection*{What Is The Difference Between A Digital Event And A Virtual Event?}

In his attempt to explain why virtual objects are as real as non-virtual objects and how virtual objects exert causal powers which affect users, Chalmers grounds his explanation in the identity claim that virtual objects are digital objects. In so doing, Chalmers seems to be working to maintain what Kim calls ``\textit{the principle of explanatory exclusion}'': there can be at most one complete explanation of any given event (Kim, ``Explanatory Exclusion,'' 233). In a sense, when first thinking about a virtual event we are presented with two potential causes for the perceptual event, say, when we first activate the virtual lightsaber and experience its virtual blueness: either the virtual object and its perceptible properties or the digital object on which it depends and generates such.

As Kim notes, there are six possibilities when we seemingly have two causal explanations—the digital object and the virtual object—for the occurrence of one event—my perceptual experience of the lightsaber’s virtual blueness.\footnote{For more on the six possibilities, see Kim, ``Explanatory Exclusion,'' 233-35.} Chalmers adopts the stance that the two causal explanations are identical; I, on the other hand, see two different events with two different causal explanations. One of the upsides to Kim’s explanatory realism is that it helps individuate explanations: ``explanations are individuated in terms of the events related by the explanatory relation (the causal relation, for explanations of events)'' (Kim, ``Explanatory Exclusion,'' 233). Thus, in the moment I experience the event, the virtual blueness of the virtual lightsaber is explained by—or stands in causal relation to—the veridical, perceptible property of the visual object; and, the virtual blueness of the virtual lightsaber is explained by—or stands in causal relation to—the property of digital blueness of the digital object. In other words, when I turn on the virtual lightsaber and experience its virtual blueness two events are taking place almost simultaneously. (1) A digital event: the computer reads and activates the digital object (code or data structure) and digital property (digital blueness) which causes the corresponding virtual object (virtual lightsaber) and virtual property (virtual blueness) to perceptibly appear mediated through the HMD; and (2) the virtual object’s veridical blueness causally explains my perceptual experience of the lightsaber’s virtual blueness.

These two events are best understood as the two articulations of a single \textit{occurrence}: the actualization of the digital object's standing potency into the perceptible actuality of the virtual object. Kim's principle is what individuates that occurrence into its distinct events—it is what licenses the claim that there are \textit{two}, and, therefore, that the virtual object is not identical to the digital object on which it depends. The Aristotelian register names what binds them: the one is the potency, the other its actualization by the rendering computer. The more basic unit, the event as an object exemplifying a property at a time, is in this way nested within the occurrence rather than a replacement for it.

\subsubsection*{Virtual Objects, Properties, and Events}

My distinction between a virtual object and the digital object on which it depends does not invalidate Chalmers’ argument that virtual objects are real, but it does, I believe, offer us a more useful philosophical conception of such by allowing us to more fully pin down a virtual object’s causal powers which affect the user. Accordingly, I would define a virtual object as the virtually rendered output of a digital object; a virtual property, then, is a perceptible property that is the virtually rendered output of a digital property; and, lastly, a virtual event is the perceptual experience of a virtually rendered digital event. Furthermore, this distinction is useful because it allows us to clarify and individuate between two different explanations for two different events. This individuation, therefore, more wholly captures how virtual objects, virtual properties, and virtual events ontologically depend on and are causally explained by their underlying digital objects, digital properties, and digital events, as well as causally explain our perceptual experience of a virtual event.

\subsection*{What is a Phantasmic Thing}

We must now consider other types of virtual artifacts that are revealed in a particular capacity of the human animal: namely, Aristotle's conception of \textit{phantasia}. However, before proceeding, a terminological clarification is in order. The phantasmic object is itself a \textit{virtual} object—not a non-virtual thing of objective reality, but the rendered actualization of a real ground on which it ontologically depends; and, it differs from the virtual-reality object only in \textit{kind}: the virtual-reality object is rendered by a computer and mediated by a head-mounted display, whereas the phantasmic object is rendered by the soul through \textit{phantasia}. If anything, the phantasmic object is the more originary of the two: the imagination renders virtual objects natively, and virtual-reality technology may be understood as a late, external prosthesis for a power the soul has always exercised.

\subsubsection*{On What Does a Phantasmic Thing Ontologically Depend?}

Aristotle states that \textit{phantasia} ``does not exist apart from sense perception'' (Aristotle, \textit{De Anima}, 428b, 431a). It is initiated when perception occurs, and it continues to function even when direct sensory input is absent. This is why \textit{phantasia} enables one to recall, reconfigure, or anticipate things not currently being perceived. Aristotle further describes \textit{phantasia} as ``some kind of movement'' produced by sensory activity (Aristotle, \textit{De Anima}, 428b, 431a). This implies that sensory input leaves a residue — a \textit{resonant kinēsis} (\gk{κίνησις}), the after-motion that lingers in the soul once the object that stirred it has withdrawn — and on which \textit{phantasmata} depend. In this sense, \textit{phantasia} serves as a bridge between perception and cognition. He distinguishes \textit{phantasia} from both thought (\textit{nous}) and perception, emphasizing that it is fallible: ``\textit{Phantasia} is an affection [\textit{pathos}] which lies in our power whenever we choose'' (Aristotle, \textit{De Anima}, 427b). This suggests that—unlike direct perception, which is always true in its apprehension of its proper objects—\textit{phantasia} can produce false or distorted images based on prior sensory experiences. Logically, then, \textit{phantasia} does not function independently of perception, but instead relies on perception as its ontological ground. In other words, in order to recall an apple we must have \textit{already} experienced an apple either directly or conceptually and have that experience embedded in the soul as a discrete chunk of data. Without this chunk of data, it would not be possible to recall an apple. This residual trace of a perceptible object, which I will refer to as a \textit{resonant thing} or \textit{object},\footnote{I use ``object'' here in the metaphysical sense at issue throughout this section—the relatum of an event, paired with a property and temporally situated—rather than in the sense of a discrete substance. The usage is licensed by Aristotle's own treatment of the residual sense-impressions (\textit{aisthēmata}) as ``themselves objects of perception'' (\textit{De Insomniis} 2, 460b2–3). This residual trace is what, in Part I, I term the \textit{resonant kinēsis}: the persisting eidetic and orectic impression—the \textit{resonant aisthēma} and \textit{resonant epithymia}—deposited by a completed perception, upon which any later \textit{phantasma} of the perceived object ontologically depends.} is the discrete residual trace of the apple that exists latently, or \textit{potentially} within the soul and serves as the ontological ground for the recollected apple. If that resonant object did not exist, it would be impossible for the recollected apple to exist. Or, if a person suffers traumatic brain injury and falls into a persistent vegetative state, despite being alive and still possibly containing the necessary resonant object in their mind, they will be unable to exercise the cognitive processes needed for active recollection. Lastly, the phantasmic apple only exists when the person is actively recalling or imagining it, whereas the resonant object on which it depends exists as a potentiality regardless of whether the person is actualizing it. In other words, the resonant object only depends on the soul of the person where it is stored and which actualizes and renders it as a perceptible object, whereas the phantasmic object ontologically depends on a particular resonant object and the soul of the being that stores, renders, and mediates the perceptible experience. In Aristotle's idiom, the resonant object is a first actuality—a developed yet unexercised retention of past perception—and the phantasmic object is its second actuality, that retention actively rendered by \textit{phantasia}.

In fact, I believe there are clear similarities between virtual and phantasmic things.\footnote{The parallel drawn here between the computational case (the digital object and its rendered virtual object) and the mental case (the resonant object and its rendered phantasmic object) is structural and heuristic—a shared logic of a real ground rendered into a perceptible appearance—not an endorsement of any computational theory of mind. The disanalogies are as telling as the analogy: most pointedly, the perceptual trace carries an orectic charge—it \textit{matters} (recalled grief is felt as grief)—that no data structure possesses, so the soul's rendering is desire-laden in a way the computer's is not.} To reiterate, within the context of a digitally rendered object, entity, or event, \textit{virtual} refers to the final rendered output of the digital object. Thus, ontologically, the virtual object depends on the digital object which is grounded in computational processes which are themselves grounded in physical processes on a computer. Analogously, phantasmic things are mental objects (the ``perceptible'' image) that depend on prior perceptual traces (resonant things) which are grounded in the functions of \textit{phantasia} which are themselves grounded in powers exercised by the ensouled being. There are, however, important disanalogies. The resonant thing is subject to degradation, whereas the digital object—so long as its hardware and code remain unchanged—reproduces the same result identically; the mental trace carries no such guarantee. More telling still is a difference of charge. The perceptual trace carries an orectic weight—it \textit{matters}, as recalled grief is felt as grief—that no data structure possesses; the soul's rendering is desire-laden in a way the computer's is not, and it is this charge, latent in the resonant trace and rendered into the phantasma, that the analysis to come will track as the orectic spring of action.

Just as our virtual lightsaber in \textit{Trials of Tatooine} depends on lines of code that determine its perceptible properties, a phantasmic thing depends on a mental trace that derives from direct perception (\textit{aisthēsis}) which I will call a resonant thing. The mental equivalent of a digital object is the stored \textit{resonant kinēsis}, the residual imprint of past sensory experience. These traces are not perceptible as discrete objects in and of themselves, just as a digital object is not perceptibly identical to the virtual object rendered from it. Now, while one might be able to open the file and view the digital code on which a specific virtual object depends, our current technology does not let us actively view residual perceptual traces stored in the soul. Even if one day we are able to see exactly where a specific perceptual trace or resonant thing is stored, it will be perceptibly distinct from the phantasmic thing that we perceive in our mind. Conceptually, however, I believe we can understand the functioning of resonant objects as similar to that of digital objects. Much like how a virtual object only exists when computational processes actively read and project the digital object as a rendered image, a phantasmic thing is only experienced when the mind actively \textit{renders} it through \textit{phantasia}. The mediation of a phantasmic thing occurs through the soul's faculties rather than an external technological device like an HMD. However, similar to VR, consciousness serves as a mediating apparatus that enables the perception of what is imagined, recalled, or projected. Thus, just as a virtual object requires a functional technological infrastructure to appear, a phantasmic thing requires a functional perceptual and cognitive infrastructure. The virtual object \textit{depends} on the digital object, and the phantasmic thing \textit{depends} on the residual perceptual trace, the resonant thing. It is crucial to note that the perceptual trace can exist without being actively rendered into a phantasmic thing: even when one is not consciously recalling a memory, the perceptual trace that allows for recall remains embedded in the mind. This is akin to a digital object being stored in a computer's memory even when the virtual object it generates is not being actively rendered. Unlike a digital object, which is stored externally and remains relatively stable, perceptual traces are stored in an embodied, mutable form within the soul's lived experience. Phantasmic things are, therefore, ontologically dependent entities: they require a ground (perceptual traces or resonant things), an active rendering process and a mediating apparatus (\textit{phantasia}) in order to be and be consciously experienced (the perceptible object that appears within the mind).

Accordingly, a memory is not a direct imprint of the past but, in the moment it is recalled, a rendered reconstruction (phantasmic thing) of past sensory data (resonant things). It does not exist independently of the soul's storage and retrieval processes in the same way that a digital object exists even when not being rendered as a virtual object. An imagined object (phantasmic thing) could then be seen as a reconfiguration of resonant things: a combination of past perceptual experiences into a perceptible composition. Like a virtual object rendered from code, it is an active manifestation that requires \textit{phantasia} for mediation. Similarly, projected hopes or fears into the future are future-oriented renders constructed from previous experiences (perceptual traces) and future expectations. They exist within the mind's capacity for simulation, extending perception beyond the immediate moment into an imagined future.

Lastly, just as virtual objects possess causal powers that serve as an indication of their reality (they interact with other virtual objects, affect people, and produce real-world outcomes), phantasmic things, similarly, are causally operative. For instance, a traumatic memory can affect our emotional responses and decision-making in the future. One can act based on imagined outcomes: a fear of a future wrought by climate change can move someone to act more environmentally responsible in the present. Thus, phantasmic things exert real causal force on the subject, making them ontologically significant beyond mere subjective illusions. Just as is the case with virtual objects where we do not experience the data structure itself (digital object) but the rendered virtual object, so too with our experience of phantasmic things. We do not engage with the resonant things or perceptual traces themselves, but with their reconstructed, rendered forms as phantasmic things. The phantasmic thing may appear directly within the mind's awareness, but it is mediated by cognitive structures in the same way that a virtual object appears within a VR headset but is mediated by computational structures. Phantasmic things, therefore, are fundamentally and experientially distinct from their perceptual traces or resonant things on which they ontologically depend.

\subsubsection*{Resonant and Phantasmic Properties}

Just as a virtual property (e.g., the virtual blueness of our lightsaber) is causally dependent on the corresponding digital property (the code specifying the object's color), a phantasmic property (e.g., the recalled warmth of a deceased parent's embrace or the vivid redness of an imagined apple) is causally dependent on the corresponding resonant property (the stored perceptual trace that enables its reconstruction). This distinction highlights a key ontological relation: the resonant property functions as the causal basis for the phantasmic property; and the phantasmic property is perceptually intrinsic to the phantasmic thing, just as the virtual property is intrinsic to the virtual object. A phantasmic property, then, is not identical to its resonant basis but is instead rendered as an experiential phenomenon in consciousness.

Having laid the ontological grounds, we must now discuss the difference between the functional and perceptual properties of phantasmic things. Resonant properties (functional cause) are the neural or mental structures that store, process, and retrieve perceptual traces and are analogous to the \textit{digital property} of a digital object. These exist even when not being actively experienced. Phantasmic properties are the experienced qualities of a phantasmic thing, such as the imagined color of an object and are analogous to the \textit{virtual property} of a virtual object. They only exist when \textit{phantasia} actively renders them in consciousness. So, a resonant property might be the encoded neural trace of the redness of an apple. Whereas the phantasmic property would be the recalled \textit{vivid red} of that apple when one consciously remembers it. Just as the digital object's \textit{digital blueness} is not blue in a perceptual sense, a stored resonant trace of an apple's redness is not \textit{red} in and of itself—the phantasmic redness is only generated when actively recalled.

\begin{quote}
(1) When actively contemplating, remembering, or anticipating/projecting, we perceive (only) phantasmic objects\footnote{The inclusion of contemplation is licensed by Aristotle's thesis that the intellect ``never thinks without an image'' (Aristotle, \textit{De Anima}, 431a16–17): even when \textit{nous} grasps a universal, a \textit{phantasma} is rendered as its vehicle, so a phantasmic object is present in every exercise of \textit{phantasia} — contemplative no less than recollective or anticipatory — and the premise holds across all three modes.} \\
(2) The object we perceive is the causal basis of our perceptual experiences. \\
(3) When actively contemplating, remembering, or anticipating/projecting, the causal bases of our perceptual experiences are phantasmic objects \\
(4) Resonant objects are the causal basis of the phantasmic objects we perceive \\
∴ \\
(5) Phantasmic objects are not resonant objects.
\end{quote}

Accordingly, we have two different kinds of properties of two different kinds of things, both of which have their own causal power: the resonant property (the latent perceptual trace) functions to cause the experiential blueness of the phantasmic object, which causally explains our experience of the blueness of the remembered lightsaber. Therefore, the relation between the phantasmic object and our perceptual experience, and the relation between the phantasmic object and the resonant object on which it ontologically depends is \textit{causal} or \textit{explanatory}. To conclude, therefore, I will define ``resonant property'' (causal input) as ``the stored perceptual trace existing latently in the brain'' (Perhaps existing ``potentially'' in the brain and must be ``actualized''). And, I will define phantasmic properties (perceptible effect) as the perceptible qualities intrinsic to phantasmic things, causally dependent on resonant properties but irreducible to them, and responsible for the phenomenological reality of mental experience. Just as virtual properties determine how we experience a virtual world, phantasmic properties determine how we experience memories, imaginations, and projections—making them ontologically significant and causally powerful.

\subsubsection*{Resonant and Phantasmic Events}

A phantasmic event is the dynamic process by which a phantasmic thing (a memory, an imagined object, or a projected hope/fear) is generated, rendered, and experienced in consciousness. This event unfolds in an ontological hierarchy wherein different layers of reality interact. So, for example, when I am actively recalling turning on the virtual lightsaber for the first time two events are taking place almost simultaneously. (1) A resonant event: \textit{phantasia} reads and actualizes the resonant object (residual perceptual traces) and resonant property (residual resonant blueness) which causes the corresponding phantasmic object (rendered recollection of lightsaber) and phantasmic property (recollected blueness) to perceptibly appear in our consciousness mediated through \textit{phantasia}; and, (2), the phantasmic thing's phantasmic blueness causally explains my perceptual experience of the recollected lightsaber's blueness.

A phantasmic event, then, is a temporally situated moment in which a latent resonant thing (residual perceptual trace) and its resonant properties is actualized by \textit{phantasia}, generating a phantasmic thing, which is experienced through its phantasmic properties, resulting in a cognitive-affective response. In other words, a phantasmic event is the perceptible experience of a cognitively rendered resonant event.

\begin{quote}
(1) Resonant Thing (Ground): The perceptual trace existing latently in the mind that serves as the causal basis for the phantasmic thing. \\
(2) Phantasmic Thing (Rendered Object): The constructed mental representation. \\
(3) Resonant Properties (Ground): The specific property of a resonant thing that serves as the causal basis for a phantasmic property. \\
(4) Phantasmic Properties (Perceptible Quality of Rendered Object): The rendered experience of the stored trace, the resonant property, actively mediated by \textit{phantasia}. \\
(5) Resonant Event (Ground): Actualization of resonant thing and resonant properties (latent perceptual traces) that serves as the causal basis for the phantasmic event. \\
(6) Phantasmic Event (Perceptual Experience): The perceptual experience of the rendered resonant event.
\end{quote}

Temporally speaking, a phantasmic event unfolds in four distinct phases, corresponding to the actualization, rendering, and experience of the phantasmic thing. Phase 1 - Actualization: begins when the subject's capacity for \textit{phantasia} is engaged and a perceptual trace (resonant thing and \textit{properties}) is retrieved from the mind and is analogous to a computer reading the digital object before rendering a virtual object. This process may be intentional (e.g., deliberately recalling a past event) or spontaneous (e.g., a memory being triggered by an external sensory stimuli). Phase 2 - Rendering: \textit{phantasia} renders the perceptual trace (resonant thing and \textit{properties}) into a perceptible thing (phantasmic thing and \textit{properties}). This process may involve recombination, distortion, or augmentation (e.g., imagining a new scenario based on past memories), and the phantasmic thing is a rendered mental virtuality. Phase 3 - Phenomenological Presence: the phantasmic thing enters consciousness, where it is perceived with phantasmic properties (e.g., color, shape, emotional tone). The experience is veridical within consciousness even if it does not correspond to present external reality. Phase 4 - Causal Effects: the phantasmic event influences cognition, emotion, and behavior. For instance, the subject may react emotionally after recalling a nostalgic memory of a deceased loved one, perhaps leading the subject to contacting family members who they may not have spoken to in some time.

A phantasmic event, therefore, is not merely an internal illusion; rather, it possesses real causal power, structuring how an individual engages with the world. Like virtual events, phantasmic events can alter perception, shape emotional states, and influence future actions and/or how we interpret past events. Thus, a phantasmic event operates as mental simulations that actively shape one's understanding of the past and experience of the present and future.

\subsection*{The Three Kinds of Thing and the Data Discrepancy}

Let us now consider the differences and similarities between a non-virtual thing, a virtual thing, and a phantasmic thing. A non-virtual apple exists as a physical object in the ``objective'' world. Its continued existence does not depend on anything once it has come to be; it simply \textit{is}—even when one is not perceiving it. A virtual apple, on the other hand, depends on (i) specific lines of code (the digital object), (ii) the physical hardware that store and render that code (a computational device and computational processes), and (iii) the user's mediating device (e.g. a display or HMD) that renders the virtual apple perceptible. The virtual apple only exists for a subject while the computer program is running and actively rendering the apple in a virtual environment. If the computer is turned off, or the program is closed, the virtual apple ceases to exist as a perceptible object (even if the underlying data structure—the digital object—remains in the computer's memory).

Similar to the virtual apple, the phantasmic apple ontologically depends on (i) a prior experience with a real apple in at least some capacity (directly experienced or perceived through other representations) to create the necessary perceptual trace (resonant object), (ii) the perceptual trace in and of itself, and (iii) the soul of the person that stores and capacity that renders that perceptual trace (a rationally ensouled being and \textit{phantasia} as the capacity or process which reconstructs stored traces into phantasmic things that are perceptually experienced.) The phantasmic apple only exists for the subject who is actively imagining it. If that person were to become suddenly distracted by external stimuli or suffer amnesia and lose the necessary perceptual trace, then the phantasmic apple will cease to exist as a perceptible object (even if the underlying perceptual trace of resonant object still exist within the mind of the subject). Metaphysically speaking, then, the difference between the virtual, phantasmic, and non-virtual apple comes down to the difference between causal and ontological dependence. A non-virtual apple and the tree it comes from have a causal relationship, but once the apple is formed it remains a physically independent entity. The apple does not need the continuing, real-time existence of the tree to remain an apple. In this sense, the non-virtual apple is not ontologically dependent on its tree, but rather is \textit{causally} derived from them. Whereas the virtual apple and its digital object or the phantasmic apple and its resonant object have an ontological relationship: neither the virtual apple nor phantasmic apple \textit{can exist} at any point in time without the active presence (and operation) of the digital code \textit{and} hardware or the resonant object and the soul of the subject which actively render the perceptual events. Therefore, even if the phenomenological human experience of the virtual or phantasmic apple were identical (which they are not) to the experience of the non-virtual apple, the difference in their hierarchy of being dictates that, while they are all very much \textit{real} entities, they are distinct kinds of entities. Whereas the non-virtual and virtual apples are externally experienced, the phantasmic apple exists \textit{internally} within the mind of the \textit{phantasizing} subject that perceives it. Another major difference between, specifically, the virtual and phantasmic apple is that the phantasmic apple requires a perceptual trace or resonant object which in turn requires the person to have actually experienced or learned about apples (through direct perception, testimony, or conceptual learning). In other words, the subject would not have that particular apple-memory or apple-image \textit{unless} they had encountered an apple (or the concept of such) at some point in their life. Once their mind encodes the apple as a resonant object or memory, the existence of that memory no longer requires the \textit{continuing} existence of the actual apple in the real world. You can remember apples that have long since been eaten or decayed. So, the memory or image is \textit{causally} dependent on the real apple having existed in the past (you needed an encounter or some conceptual input), but it is \textit{not} ontologically dependent on the apple still existing now.

To be perceived or recognized as an apple, both the virtual and the phantasmic apple presuppose that appleness has been experienced before—by the designer who models the virtual apple, and by the subject who recognizes either as an apple. In this sense each gives presence to what is absent: to an apple that was not present before one entered the virtual environment or began to remember. Now let us turn to a consideration of the experience of all the apples. A virtual apple can activate \textit{phantasia} to fill in the gaps—missing smell, taste, weight—by drawing on previous experiences with non-virtual apples. Herein lies a crucial connection between all three types of apples.

As the experiential virtual object ontologically depends on the digital object, and since a digital object like our virtual red apple is made up of a discrete data structure that takes up physical space on our hard drive, we can claim that the amount of data that makes up both the digital and virtual apple are identical. However, even if the amount of data is the same between the digital and virtual apple, our experience of that data is nowhere near the same. The virtual apple primarily consists of mathematical descriptions which create the 3D model: like vertices, polygons, textures, and other attributes needed for rendering a visually similar representation on a computer screen or HMD. Additionally, it can contain information regarding how the apple might move in a simulation, like when I move to pick it up from a virtual fruit bowl. A non-virtual apple includes details like precise atomic composition, surface texture at the microscopic level, smell, tactility, and mass distribution. Accordingly, the virtual apple contains a much smaller data footprint compared to the non-virtual apple's full complexity. So while both the virtual and non-virtual apple are very much real things, they are two different kinds of things and the primary difference between the two is what I will call, the data discrepancy. Even if every potentiality the non-virtual apple contains could be duplicated in the virtual apple, we still won't be able to taste it if we take a bite, feel the juice spreading across our tongue, or experience the pleasing scent that wafts up to our nose. However, this data discrepancy only holds true in our current moment. If technology reaches a point where we can inject the human sensorium into a virtual world with the same level of data and potentiality as the non-virtual world, then what is the difference between the non-virtual and virtual apple?

While the data discrepancy is made quite clear when comparing a virtual object to a non-virtual object, we can also see this discrepancy at work in different types of media. For instance, here we have three different depictions of storming Omaha Beach in Normandy France on D-Day. Though profound, the photo fails to capture the same level of intensity that the war documentary does, and that pales in comparison to the clip from \textit{Saving Private Ryan}. Here, we see a gradation in the amount of data presented by each respective media format. Each progressive increase in sensorial data amplifies the viscerally felt response to the content. Thus, simply viewing the photo will not be felt as intensely real as say watching the movie clip. Why? At least one reason and perhaps the defining differentiation is the additional sensorial data. This suggests that media which incorporates a wider gamut of sensorial data will be experienced as more ``real.'' This is especially true of virtual environments. The more sensorial data it contains as well as how closely it can simulate the total possibilities for action present in a non-virtual environment (lighting, sound, interactivity, realistic entities, etc.), then, potentially, the more ``real'' it will feel for the person who experiences it. This has profound implications regarding the nature of reality. It suggests that if we can create a virtual apple with the same amount of data and interactive potentialities as a non-virtual apple, and if we are able to wholly immerse the human sensorium into a virtual environment, then our discrete phenomenological experience of the virtual apple will be identical to that of the non-virtual apple. However, this convergence would be merely experiential, and it would leave the deeper difference wholly untouched. Even were our experience of the two apples identical in every sensory respect, they would remain different in kind—for the difference between them was never experiential to begin with, but ontological. The virtual apple is, and remains, a potency rendered into act by a renderer and ontologically dependent upon the digital object that grounds it; the non-virtual apple is already fully actual and depends on nothing further to be what it is. No closing of the data discrepancy can close that gap, for the data discrepancy is a difference of sensory bandwidth, whereas the difference in kind is a difference in mode of being. The two apples may become phenomenologically indistinguishable; they cannot thereby become the same kind of thing. This is also why the limit case—a total verisimilitude that closed the discrepancy entirely—would be no triumph of the real but a foreclosure: a world grown too seamless ever to break is a world that can no longer disclose its own renderedness, and what such a world would conceal is precisely the difference in kind just established.

Each of the three apples is tied to something physical, though not in the same manner. The non-virtual apple is itself a physical thing in objective reality, depending on nothing further for its continued existence; the virtual apple depends on the digital object, which occupies physical space on the hard drive that stores it. The phantasmic apple, however, is not lodged at some site in the brain, as though each memory were an item filed at an address. Its ground, the resonant trace, is an enmattered modification of the ensouled body itself—a persisting alteration of the living, en-formed body through whose capacity, \textit{phantasia}, the apple is rendered. Its existence therefore depends not on a cerebral coordinate but on the physical integrity of that ensouled body, its matter and its form together; so a grave injury to the brain may leave one simply unable to recall what one once could—not because a stored item has been erased from its place, but because the body that bore the capacity has lost the integrity on which that capacity depended.

The head-mounted display is only the most enclosing case of a more general arrangement. What actualizes a digital object into a perceptible virtual one is a renderer and a display; the headset is one such display, the monitor another, and the structure—a digital object rendered into act and mediated to a perceiver—holds across them. What changes from device to device is not the kind of being produced but the \textit{bandwidth} of its production: how much of the non-virtual object's full sensorial field the medium supplies, and how much it leaves the perceiver to furnish. This is the data discrepancy named from the side of the world. Named from the side of the perceiver it is its converse, which I will call the \textit{phantasia-load}: the smaller the data a medium delivers, the more the beholder's own image-making must contribute to make a world of it. A head-mounted, room-tracked world carries a small discrepancy and asks little of \textit{phantasia}; a world rendered on a flat screen and worked through a controller carries a large one and asks a great deal. Both produce virtual objects in the single sense established here, and both enclose a perceiver in a designed world; they differ only in the route by which that enclosure is achieved—and it is that difference of route, not any difference in kind, that the two studies to come will trace.

The data discrepancy, then, is the axis along which these three kinds of thing are ranged: the non-virtual stands at the limit of full sensorial actuality, while the virtual and the phantasmic each present a narrower bandwidth that their respective renderer must actualize from potency. And it is here that the present inquiry reaches its threshold. For the gap between the data a medium presents and the full sensorial field of its non-virtual referent is not left empty: it is \textit{phantasia} that fills it, conscripting the backlog of prior perception to supply what the medium withholds. How that labor is apportioned—how much the medium supplies and how much \textit{phantasia} must contribute—is what sets one medium's route to the perceiver apart from another's; but that question, and the larger one of how such rendered worlds come to work upon us and even to reshape us, belongs to the analysis that follows. What the present account has sought to establish is only its proper burden: that virtual and phantasmic things are real, that they are distinct in kind both from the non-virtual and from the grounds on which they depend, and that they share one mode of being—a potency rendered into act.

Where the non-virtual apple is already actualized, the virtual and phantasmic apple exist as a potentiality until actualized by the computational processes or \textit{phantasia} of the person who remembers. Unlike the non-virtual apple, the virtual and phantasmic apple \textit{only} exists for the subject while the VR program is running or actively remembering. If the subject were to lose power or become suddenly distracted while remembering, then the virtual and phantasmic apple cease to exist as a perceptible object (even if the underlying digital or resonant object still exists within the memory of the computer or person).

\bigskip

\section*{2. The Justification of the Method}

\subsection*{Rhetorical-Ontological Diachronic Analysis: The Mechanism of a Reshaping}

The previous section argued that the virtual and the phantasmic are real — distinct in kind both from the non-virtual and from the grounds on which they hang — and that they share one mode of being: a potency rendered into act. What was left unaddressed was a question of quite another order: not what the rendered world \textit{is}, but how it comes to work upon, and, in time, to remake, the one who enters it. The first question was metaphysical. The second is rhetorical — and it calls for a method, not a further ontology. That method is the one I develop here, and I call it Rhetorical-Ontological Diachronic Analysis (RODA).

What RODA does is best said through what it reads. It takes a stretch of play — a bounded run of engagement — and reads it as the \textit{actualization of desire}: the passage of a desiring being from the first stir of perception, through a charged image and a committed judgment, to a completed act. That passage is the object the method traces. It then asks how a designed world drives that passage toward a substantial union with itself — the union I call \textit{rhetorical incorporation}. And it asks what that union becomes once it is won: it does not lapse when the screen goes dark, but settles into the player's \textit{hexis}, the standing disposition from which the next encounter sets out. The wager is that engaging and reshaping are not two events but one: the second is the sediment of the first, so that to trace the actualization is already to trace the remaking.

Three words in the method's name mark the three things it must do, and each governs a movement below. The analysis is \textit{diachronic}: it does not photograph a state of involvement but follows a motion through time. It is read forward, then — actualization to actualization, each motion carrying it into the next — even as it is understood backward from its end. It is \textit{ontological} before it is psychological: the union it traces is a convergence in being, not a swell of feeling, and for that it must borrow an ontology rather than a vocabulary of affect. Lastly, it is \textit{rhetorical}: a world does not merely stand before the player but moves the player, and the grammar of such moving — how an act is given its motive, and how that motive's true source is concealed — belongs to rhetoric as the ancients conceived it. Where the metaphysics carried the reality-burden, the method carries the mechanism.

\subsection*{The Producing Art and the Using Art: What the World Supplies and What the Player Does}

Two features of the change the method traces carry the argument. The first is its unity. It is tempting to think that, because a played world has both a maker and a player, there are two changes, each with its own full complement of causes; but the four causes belong to a single change, and they cannot be parcelled out between two agents. Designing and playing are indeed two acts — yet they are joined. They are joined at the relation Aristotle draws between the producing art and the using art (\textit{Phys.} II.2, 194a34–b5): ``the helmsman knows and prescribes what sort of form a helm should have, the other from what wood it should be made and by means of what operations'' (194b5–7) — the user knows the form, the maker the matter. Yet the maker is not confined to matter: Aristotle's builder has knowledge ``both of the form of the house and of the matter, namely that it is bricks and beams'' (194a23–26). What divides the two is not what each knows but what each does. What the world \textit{supplies} is the maker's finished product — its matter and the form worked into it — handed over as a given; what the player \textit{does} — the bringing-about and the for-the-sake-of-which — is the using art worked upon that given. The crossing is exact: the designer's end-product is the player's given means. To say that a world \textit{engages} the player, then, names neither a bare meeting nor a metaphor but this precise handing-over, in which the player's own efficient and final causation takes up the world's material and formal supply as the stuff of its act. The designer cannot supply the player's \textit{desire}, and so is not the efficient cause of a \textit{freely} undertaken act — its source is the player's own appetite, which the designer can only give a world to meet. At the near end the world merely sets the occasion — why the act comes \textit{now} — and leaves \textit{which way} it runs to the player. Further along it \textit{requires} an act as the price of passage, leaning on the player's course; at the limit it \textit{compels} outright, moving the body without or against desire. Such compulsion is forced motion — the one place this partition breaks, where the efficient source has passed to the world and the player is for that moment made patient.

The second feature is the shape of the actuality toward which the change tends. Part I took this shape from Aristotle's account of perceiver and perceived: ``the activity of the sensible object and that of the sense is one and the same activity, and yet the distinction between their being remains'' (\textit{DA} III.2, 425b26–27). This is the structure the method inherits — one in substrate, different in being — and it must be held with care, for the whole third movement turns on it. It does not say the two collapse into one thing. It says that one actuality is at once the actuality of both, while what it is to be the one differs from what it is to be the other. Agent and world stand so, joined in act and separate in being. This formula, and only this, lets the method affirm a \textit{substantial} oneness of player and world without the absurd consequence that they become numerically one.

\subsection*{Rhetorical Incorporation: True Being-in-a-Virtual-World}

The spine of the method is the actualization of desire chain. It is read forward — actualization to actualization, each motion carrying it into the next — yet understood backward from the \textit{orekton} (\gk{ὀρεκτόν}), the desired object that draws the whole toward itself and lends each earlier node its sense, as an end lends sense to a means. It opens at \textbf{A₀}, the ontological horizon: a sensible object and the sense-faculty given together within motion and time, the designed \textit{Umwelt} into which the player is always already thrown. The perceptual motion M₀→A₁ reaches \textbf{A₁}, completed perception, which lays down a dual residue — the \textit{resonant kinēsis}, the after-motion that lingers when the object withdraws, in its two strands: the \textit{resonant aisthēma}, the persisting eidetic trace, and the \textit{resonant epithymia}, the persisting hedonic-orectic tone. The \textit{phantastic} motion M₁→A₂ renders, from that residue, \textbf{A₂}, the \textit{phantasma} proper — the determinate, available, charged image on which every later operation works, the pivot of the chain. The cognitive motion M₂→A₃ carries the image to \textbf{A₃}, completed cognition: the \textit{phantasma} taken up by the three orientational modes — intellection, memory, the discursive-deliberative — and ratified by the orthogonal committal \textit{doxa}, the \textit{taking-as-true}, in which the emotion concretizes. The \textit{orektikon} motion M₃→A₄ issues in \textbf{A₄}, the completed act, which sediments into \textit{hexis}. That cross-pass settling is \textit{temporal sedimentation}, and it must not be confused with the within-pass residue of A₁; the two are different deposits at different scales. And the sediment is of two strands. A completed act of skill settles as a \textit{technē-hexis}, the craftsman's standing readiness; a completed act of morally weighted choice settles as an \textit{ethical hexis}, the character's. These are not two chains but two deposits of one — laid down together, each conditioning what the other can become — so that a single stretch of play commonly thickens both at once.

One motion in the chain can also be inverted, and the method must mark the inversion as carefully as it marks the completion. Ordinarily the \textit{orektikon} motion M₃→A₄ issues from the agent's own desire into the agent's own act; but a designed world can \textit{oppose} that desire with a force of its own, and where it does, the agent is made patient — the motion suffered rather than enacted. This is \textit{bia} (\gk{βία}), forced motion: the standing exception to the chain's agency.\footnote{\textit{Bia} (\gk{βία}) is commonly rendered ``by force'' or ``under compulsion.'' Aristotle fixes the sense in his account of the involuntary: ``Those things, then, are thought involuntary, which take place under compulsion or owing to ignorance; and that is compulsory of which the moving principle is outside, being a principle in which nothing is contributed by the person who acts or is acted upon, e.g. if he were to be carried somewhere by a wind, or by men who had him in their power'' (\textit{NE} III.1, 1110a1–4). The Greek turns on the \textit{archē} lying \textit{outside}: ``\gk{δοκεῖ δ᾽ ἀκούσια εἶναι τὰ βίᾳ ἢ δι᾽ ἄγνοιαν γινόμενα· βίαιον δέ ἐστιν ᾧ ἡ ἀρχὴ ἔξωθεν, τοιαύτη οὖσα ᾗ μηδὲν συμβάλλεται ὁ πράττων}'' (\textit{NE} III.1, 1110a1–2, Bywater, OCT, 1894). Forced motion in the chain is just this: a moving principle outside the agent, to which the agent contributes nothing.} It takes several forms the analyses will need — \textit{cinematic} (the cutscene, total patiency, the body's control withdrawn entire); \textit{soft} (the rebuff, the invisible wall, the leash that turns the player back without seizing the body); \textit{suffered} (the grapple, the bear's charge, control stripped and won back only by counter-input); \textit{inflicted} (the player as the agent of force upon another); and \textit{de-inverted} (the act that restores to another the agency \textit{bia} had taken).\footnote{Each form of \textit{bia} commandeers one of the channels of involvement that Gordon Calleja's Player Involvement Model supplies below (and in the Calleja appendix). \textit{Cinematic bia} is what Calleja calls a \textit{push} narrative — the designer driving the story by cutscene, Act-constraining, the player swept along — as against the \textit{pull} narrative the player solicits (the audio diary, the in-world book) (\textit{In-Game}, ch. 7). \textit{Soft bia} seizes the \textit{spatial} channel: the invisible wall, the world's edge turning the body back. \textit{Suffered bia} strips the \textit{kinesthetic} and exacts a \textit{ludic} counter-input to win it back. \textit{Inflicted} and \textit{de-inverted bia} both work the \textit{shared} channel — the one player forcing, or freeing, another agent. The \textit{affective} channel alone answers to no single form, being the felt register of them all. So construed, \textit{bia} is not something outside involvement but involvement's own channels turned against the player's agency.} Where the chain runs unbroken to its \textit{orekton}, desire is consummated; where \textit{bia} intervenes, desire is overridden, and the analyst reads an agent momentarily unmade.

The method's central event falls at the \textbf{A₃→A₄} transition, where committed cognition passes into act, and it is in just this passage that rhetorical incorporation comes to actuality. It is one in substance, dual in being, held not as two events but as one: \textit{world-disclosedness} (\textit{Welterschlossenheit}) and \textit{co-disclosedness} (\textit{Miterschlossenheit}). The parallel in the names is the claim itself — one disclosedness, opened at once toward a world and toward others, as one actuality is at once the actuality of both: one in substrate, different in being. It is the same formula Part I traced through the chain: as perception and appetite — the \textit{aisthētikon} and the \textit{orektikon} — are one in substrate, different in being (\textit{DA} III.7, 431a13–14), so world-disclosedness and co-disclosedness are one disclosedness in two beings. World and others are disclosed not in succession but equiprimordially — Being-with and Dasein-with being, in Heidegger's terms, ``equiprimordial with Being-in-the-world'' (Heidegger, \textit{Being and Time}, 149). In world-disclosedness the player is taken into the world acted in, and the world into the player's awareness — the designed \textit{Umwelt} disclosed as a place dwelt in, not a scene surveyed. In co-disclosedness the same act is an acting-together, a being-with those the player acts alongside — a ``we'' the act at once presupposes and ratifies. The two are not successive and not separable: the player drawn into the world is, in the very same motion, drawn into a being-with-others within it.\footnote{Both names are Heidegger's. \textit{Welterschlossenheit} is the disclosedness of world — its worldhood; \textit{Miterschlossenheit}, which Macquarrie and Robinson render ``co-disclosedness,'' he states of space, ``the essential co-disclosedness of space'' (\textit{BT} 145), where space is ``disclosed with the worldhood of the world,'' and I extend it to the co-disclosedness of \textit{others}. \textit{BT} 160 holds the two together: with Dasein's Being-with, the others' ``disclosedness has been constituted beforehand,'' and this ``goes to make up significance — that is to say, worldhood.'' The constituted worldhood then ``lets us encounter what is environmentally ready-to-hand … in such a way that together with it we encounter the Dasein-with of Others,'' who ``show themselves in the world in their special environmental Being'' (\textit{BT} 160). World-disclosedness and co-disclosedness are there co-constituted — disclosed in one stroke, not in succession.}

The convergence is experiential and ontological, an oneness \textit{in act} — and it is emphatically not a change of kind. The player wholly taken into a virtual world has not thereby made that world non-virtual, no more than the perceiver who is one-in-substrate with the perceived has turned the seeing into the thing seen. The reality-burden secured the virtual as a distinct kind of being; rhetorical incorporation does not undo that distinctness but presupposes it.

\subsection*{Incorporation, Re-Grounded: Calleja and the Channels of Involvement}

The chain shows that a designed world engages the player's perception, cognition, and desire — \textit{desire} taken in its full span, as the \textit{orexis} that avoids no less than it pursues. A world found tedious and boring engages that faculty too: it is an \textit{orekton} to be fled and avoided, and its chain runs to the leaving-off of play as a craved world's runs to the grasp. What it cannot say, from its own resources, is \textit{where} in the world's design that engaging happens — nor \textit{whether}, in a given case, the player has truly been taken into the world rather than merely set before it. For both the dissertation turns to Gordon Calleja, and it turns to him because his is a vocabulary built from games and tested against players, not deduced from a philosophy of mind. The turning is a sublation, in three beats — a preserving, a transforming, an elevating.

What is \textit{preserved} is incorporation itself, kept as world-disclosedness, together with the six dimensions of Calleja's Player Involvement Model, taken up as so many designed channels through which a world reaches the player along the chain.\footnote{The six dimensions of the Player Involvement Model (Calleja, \textit{In-Game}, 43–44; full verbatim articulation in the Calleja appendix): \textit{kinesthetic}, ``all modes of avatar or game piece control… ranging from learning controls to the fluency of internalized movement''; \textit{spatial}, ``engagement with the spatial qualities of a virtual environment… the sense that they are inhabiting a place, rather than merely perceiving a representation of space''; \textit{shared}, ``engagement derived from players' awareness of and interaction with other agents… cohabitation, cooperation, and competition''; \textit{narrative}, ``the narrative that is scripted into the game and the narrative that is generated from the ongoing interaction with the game world, its embedded objects and inhabitants, and the events that occur there''; \textit{affective}, ``various forms of emotional engagement… either purposefully designed into the game or precipitated by an individual player's interpretation of in-game events and interactions with other players''; \textit{ludic}, ``players' engagement with the choices made in the game and the repercussions of those choices.''} Calleja's own definition is kept whole: incorporation is ``the absorption of a virtual environment into consciousness, yielding a sense of habitation, which is supported by the systemically upheld embodiment of the player in a single location, as represented by the avatar'' (Calleja 169); and, its double axis with it: the player ``incorporates… the game environment into consciousness while simultaneously being incorporated through the avatar into that environment,'' where ``the simultaneous occurrence of these two processes is a necessary condition for the experience of incorporation'' (169). No one, moreover, is taken into a virtual world and sees the reality of it: a virtual thing is, at base, bytes of data, rendered into a screen's pixels or a headset's paired displays. The question is what we see, and what we see. Virtual leaves, grass, trees — we instantly identify them and see them as such; so then what we see (pixels resolved through a display) is a far cry from what we see (leaves, grass, and trees). Thus the power of a designed world to take the player in shows, first and foremost, its capacity to change her mode of perception: it must move her from seeing the rendered world as an artifact she stands apart from, to a mode of perception in which she perceives herself as implicated within it — a participant who acts and is acted on. Each channel names a point where the designed world's supply bears on the chain — where its matter and form actually engage some motion or node of the player's actualization. For the world to \textit{press} there is for its supply to be doing work at that point: a fallen fruit set in the path presses along the kinesthetic and spatial channels, drawing the eye and soliciting the hand; a goblin's greeting-by-name presses along the shared and narrative. The channels enter an analysis only where the world is in fact pressing — only where such work is being done — never run down as a checklist to fill.

What is \textit{transformed} is the standing of the concept. In Calleja's hands incorporation is a psychological intensification of involvement, set on ``their experientialist ontology'' after Lakoff and Johnson — for whom meaning is grounded in the body's traffic with its surroundings, ``our conceptual system emerges from our constant successful functioning in our physical and cultural environment'' (Lakoff and Johnson, \textit{Metaphors We Live By}, 180; quoted by Calleja, \textit{In-Game}, 168) — and on ``a view of consciousness as an internally generated construct'' after Dennett and Damasio (168–169). The reading I develop lifts it from that ground onto another: incorporation is the player's \textit{dwelling} in the designed \textit{Umwelt} — what Part I, after Heidegger, names ``the characteristic absorption of concern in its equipmental world'' (BT 405) — the \textit{circumspective concern} in which the tool withdraws into its use, the hammer disappearing into the hammering itself (BT 98). Dwelling I take in Thomas Rickert's sense as much as Heidegger's — not location but \textit{attunement}: ``Dwelling is an attunement that can generate various kinds of knowledge, in particular a knowledge of how the world gives back'' (Rickert, \textit{Ambient Rhetoric}, 27). To dwell in a designed world is to be tuned to how it answers — its supply giving back to the act it draws. World-disclosedness simply \textit{is} this absorption of concern, redescribed for the virtual world. I must say plainly that this re-grounding is mine and not Calleja's: where Calleja reaches for cognitive science I am attempting a humanistic return through rhetorical ontology. Yet his own text invites the move twice over. He grants that ``incorporation is ultimately a metaphor like presence or immersion'' (3) — and a metaphor leaves its ontological debt unpaid; and there is a gap, plain on the page, between the experientialist ground he claims and the phenomenological words he in fact reaches for, the ``sense of habitation'' (169), the ``inhabiting a place'' (43–44). The re-grounding does not import a foreign vocabulary so much as pay for the one he had already spent.

What is \textit{elevated} is incorporation's reach: no longer the whole but one of the two disclosednesses — world-disclosedness — joined to co-disclosedness within a single event. The joining is prefigured in Calleja's own list — for one of the six dimensions, the shared, already names involvement with one's fellows, and so points past the world toward co-disclosedness. Here Calleja earns his place as an instrument and not an ornament, for his own writing yields a test of incorporation against mere involvement, which the method numbers and extends as the criteria R1–R7 and makes the gauge of world-disclosedness.\footnote{The incorporation criteria (Calleja, \textit{In-Game}, 169–173, 178; full articulation in the Calleja appendix): R1 assimilation; R2 embodiment; R3 simultaneity (the strictest gatekeeper, Calleja's own necessary condition); R4 the spatial-kinesthetic cornerstone; R5 blending; R6 temporal apex; R7 intrusion-integration. R3–R7 are Calleja's own incorporation-versus-involvement test; R1–R2 are drawn from the double axis.} The strictest gatekeeper is R3, simultaneity, which is Calleja's own necessary condition; the cornerstone is R4, the spatial-and-kinesthetic base ``without [which] incorporation cannot take place'' (170). The claim to rhetorical incorporation is thereby kept from mere assertion: it is read off conditions as a threshold met or unmet, and it returns a verdict — achieved, merely-involved, or ruptured.

Why Calleja, and not another account of immersion or presence? The answer runs through his own correcting moves, not through a contest with rivals. He excises the magic circle (\textit{In-Game}, 46–53)\footnote{The \textit{magic circle} is Johan Huizinga's figure for the boundary that sets play apart from ordinary life, taken up as a near-axiom of game studies (Huizinga; Salen and Zimmerman; Juul). Calleja gives a large portion of his model-chapter to dismantling it. His central objection: one cannot keep Huizinga's circle while dropping Huizinga's thesis that play is set apart, for the circle just \textit{is} the instrument of that apartness (46–48); and for digital worlds Juul's spatial version — the circle located in screen and controller — is redundant, since the only on-screen space one can act in is the navigable environment itself (48–49). With Ehrmann he denies that ``reality'' can be bracketed against play at all. The excision matters here because a world real in kind cannot be sealed behind such a membrane: incorporation presupposes the very traffic the magic circle would forbid.}, that supposed clean wall between the played world and the lived one — and the excision is exactly what a world real in kind demands, since what is real in kind cannot be sealed behind such a membrane. And he corrects the one-way metaphors of immersion and presence, which figure the player as merely lowered into a world; his double axis restores the two-way traffic — the world taken into awareness as much as the player given over into the world — that ``one in substrate, different in being'' requires.

\subsection*{Consubstantiation and the Grammar of Motives: Burke at the Other Pole}

Incorporation joins the player to the world acted in. It does not join the player to those acted alongside; and the chain, which tells what moved the player, does not tell how that motion is afterward rendered, nor how its true source comes to be concealed. Two things are left unaccounted for, then — and neither is a \textit{residue} in the sense the term carries here, but each is a real gap: first, a consubstantial bond among several agents, a ``we'' that can itself \textit{source} an act, so that the deed springs from the bond and not from any single locus; second, a grammar by which a finished act is furnished with its motive, and by which that motive's source may be displaced. For both the dissertation turns to Kenneth Burke, and the turning is again a sublation.

What is \textit{preserved} is consubstantiation. ``In being identified with B,'' Burke writes, ``A is `substantially one' with a person other than himself. Yet at the same time he remains unique, an individual locus of motives. Thus he is both joined and separate'' (Burke, \textit{Rhetoric}, 21). That ``joined and separate'' is, word for word, the grammar of ``one in substrate, different in being,'' and its ground is act: ``substance… was an act; and a way of life is an acting-together'' (21). The bond, moreover, is born of division and answers to it — ``Identification is affirmed with earnestness precisely because there is division. Identification is compensatory to division'' (22) — and the player arrives divided in the most literal way, an actual body set over against the avatar acted as. What is \textit{transformed} is the engine of the bond. For Burke the ``we'' is worked by symbolic persuasion; in the RODA method it is worked by shared desire — the ``we'' ratified at A₃, where the committed cognition of agents acting together fixes a common \textit{orekton}. What is \textit{elevated} is its place: consubstantiation becomes co-disclosedness, correlative to world-disclosedness, and folded into what Part I called the affective architecture of bodily-being-in-the-world-with-others.

The grammar of motives Burke supplies on his own warrant. He ``refers human actions to seven causes (\textit{aitia}): chance, nature, compulsion, habit, reason, anger, and desire'' (Burke, \textit{Grammar}, 292) — and that warrants reading the seven causes of \textit{Rhetorica} I.10 (1369a5–6) as a table of motive. The crosswalk from those seven to the pentad is mine, built on Burke's own move. His explicit map runs from the \textit{four} causes; I mark the difference rather than blur it. The pentad itself — Act, Scene, Agent, Agency, Purpose, with Attitude as the sixth term\footnote{The dramatistic pentad (Burke, \textit{Grammar}, xv; full articulation in the Burke appendix) is Burke's answer to the question of motive: any complete statement of what was done names five terms — Act (``what was done''), Scene (``when or where it was done''), Agent (``who did it''), Agency (``how he did it''), and Purpose (``why''). Their analytic force lies less in the terms than in their \textit{ratios} — the proportions an account strikes between them (scene–act, agent–act, and the rest), since it is by weighting one term against another that a motive is assigned; and dramatism reserves the pentad for \textit{action}, the domain of the symbol-using animal, as against sheer \textit{motion}. Attitude — ``the preparation for an act,'' an ``incipient act'' (Burke, \textit{Grammar}, 20) — Burke subordinates to Agent in the 1945 text and promotes to a sixth term only in the 1969 edition; I take it here as that sixth term.} — the method takes as a grammar of motive-\textit{attribution}, ``a concern with the terms alone'' (xvi), not a psychology of the drives beneath an act: it names what a finished deed gets ascribed to, and lets one and the same deed be relocated from term to term. To it attaches a concealment diagnostic, which I frame narrowly and not as a blanket suspicion: ``in banishing the term, far from banishing its functions one merely conceals them,'' so that the analyst may ``detect its covert influence even in cases where it is overtly absent'' (21). The diagnostic catches a displaced motive — a source smuggled in under another term's name, as when an act the agent in fact sourced is re-credited to the Scene that merely housed it — and claims nothing more sweeping than that.

One last piece of evidence binds the two sources into one. Calleja, grounding incorporation in cognitive science, and Burke, grounding consubstantiation in symbolic action, arrive — each blind to the other, neither ever citing the other — at a single figure, the joined-yet-separate. Calleja reaches it as a \textit{two-way axis} — the player takes the game-world into consciousness and is, through the avatar, taken into it, in one act — and makes their simultaneity the necessary condition of incorporation (169). Burke reaches it as consubstantiation: ``both joined and separate,'' identified with another yet still ``an individual locus of motives'' (\textit{Rhetoric}, 21). The two are one structure — one actuality, two beings — reached from grounds that share no vocabulary, no method, no common source. That structure is what the method names rhetorical incorporation; world-disclosedness and co-disclosedness are its two poles — not two lenses laid over a scene, but the two sides a single convergence brings to light.

\bigskip

\bigskip

\section*{3. Introduction to the Applications}

\subsection*{The Virtual Environment as Rhetorical Ecology}

Over the course of his career, Jakob von Uexküll would conceive and propose a radical new perspective for zoological studies: a new way to perceive the relationship between the organism and its environment. Dissatisfied with the traditional scientific perspective which saw a single world within which all living organisms are hierarchically situated, he argued that every living creature exists within a species-specific, perhaps even within an individual-specific, perceptual world: ``for everything a subject perceives belongs to its perception world [Merkwelt], and everything it produces, to its effect world [Wirkwelt]'' (42). The receptor and effect systems constitute the organism's ``functional circle'' (49)—contemporarily rephrased as a ``feedback loop''—and ``these two worlds, of perception and production of effects, form one closed unit, the environment'' (42). Accordingly, there are an infinite variety of perceptual worlds as each individual's subjective environment [Umwelt] —metaphorically figured as a ``soap bubble'' which encloses the subject on all sides—is determined by their respective capacities for perception and production of effects. This perspective not only helps to blur the hierarchical distinction and prioritization between humans and animals, in that all organisms are now understood to perceive and react to environmental sense-data as signs, but also lays the groundwork for Heidegger's deep ontological adaptation of Umwelt to articulate human being-in-the-world. Indeed, as Heidegger argues, ``[t]hat world of everyday Dasein which is closest to it, is the environment'' (BT 94). Thus, it seems as if for both scholars the environment or Umwelt constitutes the grounds through which the perceiving subject comes to be, it ``forms a self-enclosed unit, which is governed in all its parts by its meaning for the subject'' (Uexküll 144).

While Uexküll primarily focused on examining the animal environment, Heidegger's method of analyzing human being by situating it within an environmental context reinvigorates the concept of environment within a human perspective: it is no longer something ``out there,'' as the natural world against the world of human creation, but rather serves as the overarching orientation through which things are related to other things. Heidegger's reoriented understanding of environment works to break down the dichotomy between natural and artificial or built environments by figuring environment as the network or ``totality of involvements'' within which all understanding and possibilities for action originate. Indeed, as Katherine Hayles comments, this is a conception of environment as the ``crucial scaffolding and support for human cognition'' (92). In Gibsonian terms, the environment contains the total sum of possible affordances for thought and action for the respective subject: ``[i]t implies the complementarity of the animal and the environment'' (119).

According to Timothy Morton, ``[i]n order to have an environment, you have to have a space for it; in order to have an idea of an environment, you need ideas of space (and place)'' (11). The environment's spatial character is itself relational, answering to Heidegger's structure of ``worldhood'' as a context of relationships (\textit{Bewandtniszusammenhang}) — a structure reflected in Michael Heim's analysis of virtual worlds, which reads Heidegger's definition as suggesting that ``the entities in a world are constituted by their interrelationships … a continual linking of one thing with another'' (Heim, \textit{Virtual Realism}, 91). Lefebvre seizes on the interconnected nature of spatial environments and poses a new way to understand the lived environment of the human subject. He was one of the first to highlight the fact that our environment not only shapes who we are and how we understand our place in the world, but also how ``real'' lived space is constructed through social practices: social relations ``project themselves into a space, becoming inscribed there, and in the process producing that space itself'' (129). The key element to Lefebvre's thinking, according to Hayles, ``is conceptualizing place not as a fixed site with stable boundaries but rather as a dynamic set of interrelations in constant interaction with the wider world, which nevertheless take their flavor and energy from the locality they help to define'' (185).

Lefebvre's work helps us understand how a built environment is inflected with various ideologies and political implications and demonstrates how those environments work, in turn, to inflect our lived experience of them. Sarah Ahmed illustrates this in her critique of institutional spaces as ``orientation devices'' which are shaped by ``what resides within them'' (157). In other words, these spaces are shaped by the bodies that built and occupy them and thereby orient our encounter with other bodies' experience within them. To return to the natural environment/built environment binary, we can see that fundamentally they function in the same way: they orient our understanding, actions, and experiences in the world. Just as the institutional space orients those who occupy it, so too does the natural world. If, on a hike, we hear the telltale sound of a rattlesnake, our orientation is directed to a potential threat and that plays a dramatic role in our future actions. Therefore, in essence, all environments, artificial or natural, exert fundamental rhetorical force that directs our attention and moves us to do something. However, while all environments are fundamentally rhetorical, there are undoubtedly different kinds of environments. The university or academic environment is a very different environment than, say, Yosemite in regard to the objects and entities that populate each respectively. As such, I believe we need to make a distinction between environment and ecology. Accordingly, I define ``environment,'' as a specific, bounded location wherein circulation between various interrelations take place; and although circulation can occur throughout multiple environments, the rhetorical effects, analytically speaking, are best examined within a particular localized system. That said, I believe it is better to consider the rhetorical force exerted by environments along ecological lines. For me, ``ecology'' represents the broader, all-encompassing scale of multiple environments. As such, I use ``ecology'' to describe the theoretical lens through which we can understand the rhetorical effects of environments in general; in other words, ``rhetorical ecologies.''

Even if how a specific individual makes sense of their position within a specific environment and how they conceive the possibilities for action, for the production of effects, is linguistically circumscribed, both space and language itself is affected by and embedded with various environmental, cultural, institutional, and political perspectives. Our practical consciousness, therefore, is not an isolated product. It does not, as Jenny Edbauer argues, exist ``outside the prior and ongoing structures of feeling that shape the social field,'' the interconnection of forces, rhetorics, moods, and environmental, social, and political histories that cohere around material spaces and physical locations (10); rather, our practical consciousness is formed in relation to the dynamic rhetorical ecology into which we are born and live our lives. To this end, Thomas Rickert recasts rhetoric's work as ambient and ontological: ``persuasion gains its bearings from an affectability that emerges with our material environments both prior to and alongside the human'' (254). This way of being forces us to recognize that our engagement with environments and the things within them not as just simple objects, but rather as active constituents of the world that shape us and are shaped by us. Indeed, as Heidegger concludes, rhetoric is the self-interpretation ``of concrete being-there, the hermeneutic of being-there itself'' (Concepts 75); rhetoric is the means by which we understand our being-in-the-world and how our environmental context shapes this understanding. An understanding of the rhetorical force exerted by environments as ambiance dissolves the dichotomy between ``passive'' material environments and human action in favor of, what Byron Hawk calls, ``emergent entanglements that tacitly coproduce dwelling'' (20).

Human understanding and discourse does not simply take place within a specific environment, it takes shape within dynamic relations between human nature, environment, and social infrastructures. Indeed, a change in, for example, our material information environment brings with it a change in human being. Social media in particular is a great example to understand how a change in our environment can lead to fundamental changes in our everyday being-in-the-world. It exemplifies how new media environments fostered through new technologies of communication transform how we understand these new techno-social matrices. In short, it demonstrates how an environment isn't merely shaped by us, but how it plays a role in who we are, how we know, as well as what we say and do. Our everyday environment serves as the backdrop through which things and objects show up as what they are, as we conceive and understand them. However, as Rickert emphasizes, our everyday environment is already ``choreographed in advance'' (55) and thereby rhetorically shapes our primordial understanding of it. Therefore, the range of our thoughts, expressions, and discourses is shaped by our dynamic rhetorical ecology: the concatenation of and interaction with other objects, entities, peoples, discourses, material environments, social norms, and economic and governmental institutions and policies. And it is within such rhetorical ecologies that individual and collective self-determination is realized. What an understanding of rhetorical ecologies discloses or opens to us is a complex, dynamic, and relational world consisting of mutually interacting elements which do not outright control or predetermine our experiences, but rather rhetorically attune us to the affective and symbolic ground from which consciousness, lived experience, and willed activity is derived.

If every environment is a rhetorical ecology — orienting its inhabitants, attuning them, and over time re-making the being of those who dwell within it — then the virtual environment is that proposition carried to its limit. A natural or built environment is choreographed only in part: shaped by histories, materials, and forces no single hand has authored, much of it withdraws from any design. The virtual environment is choreographed entire. Nothing in it is found; everything is rendered. Each perception-mark and effect-mark, each affordance and functional tone, each law of the place is laid down in advance by a designer, so that the choreography Rickert discerns in all dwelling becomes, in the virtual, a nearly total authorship of the world.

Such a world has no objective referent. It does not depict a reality lying behind it; it simulates one outright — a simulacrum in a deliberately minimal, Baudrillardian sense, a situation rendered rather than the representation of a prior one, though without Baudrillard's stronger claim that no real lies behind it at all, since the foregoing ontology has located that real ground precisely in the digital object — ``the generation by models of a real without origin or reality'' (Baudrillard, \textit{Simulacra and Simulation}, 1). The virtual environment is therefore situational in the strongest sense: not a scene the user surveys from without but a world the user is within, a designed Umwelt whose every relation has been set so as to orient and move the one it encloses. This is what may be called a persuasive virtual ecology — the affordance-structured horizon within which the user comes to perceive and to act; named in the present terms, it is simply a rhetorical ecology every element of which is authored.

The designer is thus the hidden agent of the whole ecology, the efficient cause for which the rendered world stands in. And yet — this is the crux — a wholly authored environment does not on that account determine the one who enters it. It attunes. The player who moves through such a world acts from a desire that is genuinely the player's own, even where the ecology has been arranged so that the player's desire and the designer's end fall into alignment. In just this sense the virtual environment is the purest instrument of environmental rhetoric: it curates without coercing, moving its inhabitant to act not by instruction but by the very disposition of a world. What in the natural environment is diffuse and unauthored — the rattlesnake's rattle bending a hiker's next step — is here total and deliberate, and so can be studied in isolation.

But to say that a virtual environment attunes those it encloses already presumes that they are, in some genuine sense, enclosed — taken up into the world rather than set before it as a picture. What that being-taken-up amounts to, and why the older language of ``presence'' must give way to a language of incorporation, is the next question.

\bigskip

\subsection*{From Presence to Incorporation}

To be enclosed by a rhetorical ecology and moved by it is, in the language of virtual-reality research, to be \textit{present}. The dominant accounts define presence as the sense of inhabiting a mediated world ``as if … only mediated by human senses and perceptual processes'' (ISPR par. 2) — a co-location felt below cognition, registering as the involuntary flinch or the stumble at an edge that is not there. Presence in this sense is, as Perelman and Olbrechts-Tyteca said of rhetorical presence generally, something that ``acts directly on our sensibility'' and is operative ``at the level of perception'' (116). It is, moreover, a matter of degree: worlds are more or less convincingly inhabited, and presence is to be distinguished from mere \textit{immersion} — the engrossment in an activity better described, with Heidegger, as a ``concernful absorption'' (BT 146) in a task than as a feeling of being-there. Such presence, though, marks only the floor. It registers even against the player's certain knowledge that nothing is there: the hand moves to steady itself against a virtual tree and passes clean through it. The reaction comes so readily because a world familiar enough to our ordinary expectations draws a familiar response with little work of the design's own — and precisely because it asks so little, this reflexive co-location must not be mistaken for the achievement the analyses track. To flinch at a world is not yet to dwell in it.

In my earlier work I pressed this distinction toward a limit I called \textit{true presence}: the condition reached when the apparatus falls transparent and the division between a subject and a world of objects dissolves, so that the user is no longer set over against the virtual world but absorbed into it — present as a craftsman is present in the work, not as a spectator before a picture. That criterion — equipment-transparency and the dissolved subject/object — I keep. The name I do not.

For ``presence,'' even qualified, carries a psychologistic and faintly binary cast: a state of mind one is in or out of, registered on a felt scale. It sits poorly with what the analyses to come will show — an achievement built up act by act, across distinct channels, over time. I therefore set the term aside and take up, from Gordon Calleja, the more exact one: \textit{incorporation} (Calleja, \textit{In-Game}) — the very union the foregoing method anatomized into its two correlative poles, world- and co-disclosedness. Incorporation names a double movement: the player assimilates the virtual environment into the field of habit and attention and, in the very same stroke, comes to be incorporated \textit{into} it — present within the world as one who dwells there. It is graded, embodied, and processual, gathered across several dimensions of involvement at once; and it preserves what ``true presence'' rightly saw — the transparent equipment, the dissolved subject and object — now as its achieved condition rather than its definition. To be incorporated into a designed environment is to have come to dwell in it. The difference from mere presence is not academic. A spectator before a vivid but non-interactive panorama is \textit{present} to it as to a picture, set over against a world that does not answer; the player who lifts a fallen fruit and carries it to a waiting goblin, or who pours a morning coffee in a gang's camp and is greeted by name, is \textit{incorporated} — taken up into the world as one who dwells and acts there, the world taken up into the player in the same stroke. It is that achieved, two-way dwelling, and not a felt intensity registered on a scale, that the analyses to come will track.

\bigskip

\subsection*{Veri(dis)similitude: The Aim, Function, and Motive of the Virtual}

If incorporation is what a designed environment accomplishes, the next question is \textit{how} — by what aim, function, and motive a rendered world draws the player in. The triad is Heideggerian: a phenomenon is delimited by asking what it is for, what it does, and what moves it (BT 24, 29).

The motive is legible in the name. ``Virtual,'' in Burke's sense, is an \textit{epithet} — a word that assigns a substance and thereby carries an implicit program of action (\textit{GoM} 57); and etymologically it means ``not actually, but as if'' (Heim 220), from the Latin \textit{virtualis}, power or potency (OED). The motive of the virtual, then, is to be taken \textit{as} the real — to pass, for the duration of the experience, as indistinguishable from the actual. From this motive follows a double function. \textit{Technically}, virtual-reality technology functions to \textbf{supplant} objective reality: to evacuate the real-world horizon and render another in its place. \textit{Rhetorically}, it functions as an \textbf{actant} (Latour 373; Gries 74–75) — the rendered world is not an inert backdrop but a granted source of action that works upon the user by symbolic and material means (Burke, \textit{RoM} 43). The rhetorical ecology of the foregoing pages, seen now from the side of the maker, is just this: an environment authored to act.

To supplant reality and be taken as real, however, a rendered world must satisfy a double and paradoxical condition, which I name \textbf{veri(dis)similitude}. It must be \textit{verisimilar} enough — ``familiar enough to our normal everyday mode of interacting and understanding objective reality'' — for the player to be thrown into it at all; and it must be \textit{dissimilar} enough — holding possibilities ``dissimilar to what is familiarly experienced in everyday objective reality'' — that what it then does to the player can carry ontological force (\textit{Veri(dis)similitude} 9). The credibility a virtual world must win merely to be entered — its \textit{ethos} as a world — is the verisimilar pole of this condition; its power to re-make those who enter is the dissimilar pole. A world too familiar cannot move; a world too strange cannot be entered. Only the veri(dis)similar world both throws the player in and, having done so, can rupture and re-ground. The verisimilar pole has a technical name already given in the ontology — the \textit{data-discrepancy}, the gap between what a medium renders and the full sensorial field of what it renders — and an experiential one, the \textit{phantasia-load}, the labor of image-making by which the player fills that gap. Veri(dis)similitude, data-discrepancy, and phantasia-load are thus one condition seen in turn from three sides: from rhetoric, from the medium, and from the player who must make a world of what the medium supplies.

\bigskip

\subsection*{Two Worlds at the Poles of the Medium}

Two designed worlds anchor what follows, chosen because they stand at opposite poles of the medium. \textit{Gnomes \& Goblins} (Wevr and Reality One, 2016) is a virtual-reality world at the near edge of the virtual: room-scale, headset-supplanted, its data-discrepancy small and its demand on the player's image-making correspondingly light. \textit{Red Dead Redemption 2} (Rockstar, 2018) is a screen-based world at the far edge: rendered on a flat display and worked through a controller, its data-discrepancy large and its demand on the player's image-making correspondingly heavy. Each is a rhetorical ecology built to incorporate; set side by side, they let us watch how the medium's data-discrepancy sets the \textit{route} to incorporation — and to the ontological change it can carry — without either pole holding a monopoly on the end. That comparison is drawn out on its own terms in the differential analysis that closes these readings; here it is enough that the two cases are not rival demonstrations of one thing but a controlled variation, the medium changed and the rest held as steady as two such worlds allow.

Each world is read in the same two movements. First the \textbf{tutorial} — how the design, without instruction, naturalizes its mechanics and draws the player from a being-in-the-world toward a being-in-a-virtual-world. Then the \textbf{disruption} — the engineered rupture at which the incorporated world breaks and, in breaking, discloses what it had quietly made of the one who dwelt in it. The first case is \textit{Gnomes \& Goblins}.

\bigskip

\section*{4. \textit{Gnomes \& Goblins}}

\subsection*{The Tutorial: Coming to Dwell}

\subsubsection*{The Candle: A Threshold Crossed by Discovery}

In the entire span of the demo, across every scene and encounter it contains, the program offers the player exactly one line of instruction, and it spends that line at the very threshold. The experience opens on a night scene and a signpost, each of its directions hung with an unlit candle, with a single lit candle resting on a stump nearby. A player who happens to look down at the tracked hands finds a line of text there: ``[l]ight a [c]andle to make a selection'' (00:33). Nothing requires that the player look, and the instruction is easily missed. Even when the instruction has been found, the task it names resists: the lighting is ``more difficult than one might imagine; it takes some experimenting'' (Salvo, \textit{Veri(dis)similitude}, 31), and the player must fumble with grip, angle, and flame before the wick catches and a way opens. What I want to draw attention to is not the difficulty itself but what the difficulty teaches, for the player who works through it learns something no instruction states. At first the signpost reads as what a lifetime of signage has settled it to be — a trail-marker, a ``direction tone'', pointing the ways one might go — and the candles read as things that merely burn. In the handling, and only in the handling, they take on another sense entirely: an ``interface tone'' (Salvo, \textit{Veri(dis)similitude}, 31), an apparatus the player must operate in order to begin the experience at all. The opinion the player holds about what a signpost is, sedimented from years of roads and trailheads, is not corrected by any caption; it is overridden by a new judgment that the player's own trying earns, ratified in the moment the flame finally catches and the world responds. In the terms this project has developed, a standing \textit{doxa} is displaced by an occurrent one, and the displacement is authored by the player rather than announced by the design.

Because this is the world's single sentence, everything after it must be learned another way, and the threshold is where the program declares — by withholding rather than by saying — what that way will be. The maker has supplied the matter and the form of the scene: the signpost, the candles, the flame, and the rule that binds them; what the maker has not supplied, and structurally cannot, is the working that turns them into a beginning. That labor, and the meaning it produces, belong to the player. Here, in miniature, is the method the whole world will run on, and the method this analysis will track through each of its beats: a sense withheld and then conscripted from the player's own \textit{phantasia} rather than declared. Nothing in this is yet a dwelling. The player stands at the edge of the world and not within it, the first skill of grip and reach only beginning to take hold; but the disposition the candle installs is precisely the one the rest will build upon and the rupture will, in the end, turn — a settled readiness to find the world's sense by working at it, cued by almost nothing. When the flame catches and the candle on the direction marked ``Goblins'' takes light, the menu gives way to a forest, and the world's one sentence has been spent.

\subsubsection*{The Forest: Dwelling, and What Withholds}

Lighting the candle on the direction marked ``Goblins'' transports the player to the main world: a fairy-like forest, complete with tiny homes in hollowed-out trees and rope suspension bridges strung between the dwellings, fireflies drifting in the dark. The player is not before this world but within it. This is the designed surround the reading will keep returning to, a virtual \textit{Umwelt} in the sense Uexküll gives the word when he bids us ``imagine all the animals that animate Nature around us [\ldots] as having a soap bubble around them, closed on all sides, which closes off their visual space and in which everything visible for the subject is also enclosed'' (Uexküll, 69) — a surrounding world furnished to be inhabited and not surveyed. The headset draws something like that bubble about the player, closing the visual field on all sides so that all the player sees is enclosed within it. In virtual reality the inhabiting is nearly literal. Because the program renders physical locomotion one to one — the player navigates the virtual space with directional steps rather than a controller, a single stride forward being just that, a single stride — a real reach becomes a reach among the trees, and the tracked hands that move there are the player's own, with no avatar interposed (Salvo, \textit{Veri(dis)similitude}, 28). Read through Calleja's PIM, whose six dimensions this analysis will use throughout to name the structural elements at work, this is the \textit{kinesthetic} and \textit{spatial} channels fused at the medium's floor. The involvement model's cornerstone criterion — space internalized as inhabited place — is here laid down in the body itself, since the controls to be mastered just are the player's ordinary powers--of stepping, turning, and reaching--along with the controller's trigger to grasp and hold objects. The felt sense of being somewhere, able to move and to handle, is not conscripted from the player's imagination but supplied by the hardware, and the cornerstone of the dwelling is set before the player has formed a single intention about it.

With the world's one instruction already spent, everything the player now learns is learned by exploration, and for a time the world simply obliges the learning. Tempted to investigate the tiny treehouses, the player is invited to push open their wooden shutters and peek inside; sticks and bridges sway when pushed; a fallen fruit can be lifted, an acorn taken up. Uexküll's account gives the fine structure of these dealings. Each animal subject, he writes, meets its objects in a pincer of two arms, one perceptive and one effective: ``With the first, it imparts each object a perception mark [Merkmal] and with the second an effect mark'' (Uexküll, 49). Each object thus carries the perception-marks, the features by which the player's attention finds it, and receives in turn the effect-marks the player's handling imprints — the two closing into the functional circle of perceiving-and-doing that his biology made the elemental unit of any creature's world. Each experimentation--grasping twigs, opening shutters, picking up acorns and handing them to goblins--builds a fluency: each succesful interaction sedimenting a (\textit{technē}), a naturalization of the mechanics for interacting in this world. This is a kind of kinesthetic learning: the player interiorizes the world's mechanics until they are second nature, a skill thickening into the settled ease of being at home somewhere. Because the dwelling is laid down in use, its criterion is transparency: ``[w]hile constructing a table in a skilled and trouble-free way, the craftsman has no conscious experience of the hammer, nails, or wood as they would were they to sit and merely contemplate those items'' (Salvo, \textit{Veri(dis)similitude}, 41). What holds for the carpenter's hammer holds, in the rendered world, for the shutters and the grip itself. When the player's dealings run smoothly, the controls and the handled things withdraw from notice, and there arises that ``concernful absorption in whatever work-world lies closest to us'' (Heidegger, \textit{Being and Time}, 101) in which no awareness remains of a subject set against a world of objects. By the time the forest has been explored for a few minutes, the player is no longer operating an interface. The player is somewhere, doing things, and that unremarked being-somewhere is what the whole tutorial has been for; they have started undergoing the process of rhetorical incorporation. The degree of psychological separation between the player and the virtual world has lessened.

One thing, though, refuses. A curiously shaped rock lies on the ground, and the hand that reaches to take it closes on nothing it can use: the rock is, in the earlier study's words, ``resistant or entirely unserviceable'' (Salvo, \textit{Veri(dis)similitude}, 40), a thing that cannot be lifted or moved or turned to any task. In refusing, it shows what the obliging world had kept hidden. The compliant objects had been ready-to-hand, so transparent in use that they drew no notice; the rock, balking, goes un-ready-to-hand in Heidegger's sense, the thing that announces itself precisely by breaking down under the hand that reaches for it. What balks is then set aside. It falls out of the player's dealings into mere perception, a thing that, proving unusable, ceases to matter. What refuses in this way is, in Rickert's sense, a thing and not an object — for ``a thing is not simply an object `out there,' picked out within a dualistic paradigm in which a subject perceives and distinguishes objects'' (Rickert, \textit{Ambient Rhetoric}, 225); its being is not used up in what can be done with it but exceeds the use-relation, withdrawing into a depth no dealing reaches, for ``[i]t is the nature of tool-being to recede from every view'' (Harman, \textit{Tool-Being}, 4) — a hidden surplus that never becomes fully present to the hand that reaches for it. Almost everything else in the forest hides that withdrawal by complying, and so the rock is the one point at which the world's material resistance becomes palpable. Through this exception the player learns, again by exploration and by nothing else, that this world has a ``cannot'' as well as a ``can.''

Yet open refusal is not this world's main way with the player. Because it is the engine of everything that follows, the design's economy deserves precise statement. The forest rhetorically encourages the player to dwell less by pushing back than by holding back: withholding cues, withholding instruction, withholding the meaning of its candles and the consequence, soon, of its bell. Because the design limits its interactive objects and says almost nothing — the old adage ``less is more'' governing the whole construction — the player's dealings are turned onto the world itself, and the player is left to supply what is kept from view. The rock withdraws from use; the design withholds the world's sense; and both throw the player onto the same conscripted engagement of \textit{phantasia}. What looks like a generous world is in this one respect austere, handing over its objects readily and its meanings almost never, and the austerity is the pedagogy: a world that never tells can only be learned by trying, and a player who can only try becomes, minute by minute, a being disposed to try. Into this dwelling a first relation is then seeded. A goblin appears---shy, glimpsed at the edges of vision and gone again, then nearer---and the affective tone of the exploration, a childish and unguarded delight, begins to gather around the creature. So the dwelling stands: a body at home in a tracked space, a fluency that needs no thought, a world of things that mostly answer and one that would not, and a creature met but not yet engaged.

\subsubsection*{The Fruit: How a World Moves the Player Without a Word}

Of all the moments in the opening, the fruit exchange is the one in which the design's method is most nakedly itself. It moves the player to do exactly what it wants — pick up a fruit and give it to a goblin — without a single instruction, by curating a world rather than issuing a command. I will therefore walk this episode through the actualization chain step by step, from the opening perception to the completed action, naming at each motion the involvement dimension through which the design applies its pressure. The two instruments answer different questions, and it is in their joint application that the method earns its keep. Calleja's model identifies \textit{which} structural elements of the game are at work upon the player at a given beat — \textit{kinesthetic, spatial, shared, narrative, affective, ludic} — while the chain says \textit{what} those elements do to the actualization of desire as it runs from perception through image and judgment into act. The channel names the designed pressure-point; the chain names the motion being pressed.

Before the walkthrough begins, however, let us put the episode's end in its proper place, for each motion in the chain takes its name from its terminus and is understood backward from the actuality it reaches. In the skilled engagements of ordinary life this poses no puzzle: the outfielder who tracks a fly ball knows the catch from the crack of the bat, and agent and world share one \textit{for-the-sake-of-which} from first perception to closing glove. Here the case is stranger, and the strangeness is the analysis. At the moment the chain opens, the player does not know that the episode will end in a gift; only the designer holds that terminus, and has held it since the scene was built. The sequence is read forward by the player, discovery by discovery, and was understood backward by the maker, who arranged every prior motion from the end it was to serve. What the walkthrough will show is those two readings converging on a single act — the player freely performing the very deed the design projected — and that convergence, rather than any command, is what it means to say the world curates rather than coerces. A found world exerts no final cause on the one who acts in it; a designed world is nothing but final causes, laid down in advance and waiting for a desire to run through them.

Before the fruit ever falls, the chain has already begun, at the horizon the tutorial has spent its minutes preparing. By this point the player has had unguided play enough that a disposition has quietly settled: in this world, things are to be picked up. The forest has, over the same minutes, produced its goblin, the shy creature making brief appearances at the periphery, glimpsed and gone, until the player has begun to attend to it. The horizon (A₀) thus holds a field of graspable things and a creature of growing interest, the \textit{spatial} and \textit{shared} dimensions of involvement already primed; and neither element of that readiness was taught, both having been cultivated by exploration alone.

The opening motion fires at a sound. High in a tree the goblin turns to a piece of fruit on one of the higher branches; grunting with effort and unable to reach it, it goes briefly inside to sharpen its axe, then returns and sets to cutting. The player hears the chopping before seeing it. Sound, in Hawk's account, ``is produced by an encounter, a disturbance, a clash with noise, and is diffracted as it resonates throughout space'' (Hawk, \textit{Resounding the Rhetorical}, 26--27); the axe biting the branch is such an encounter, and the chop ``resonates in space as a nonsemantic expression of a material environment'' so exactly anchored by the headset that those who hear it ``locate themselves in an environment via reverberation'' (Hawk, \textit{Resounding the Rhetorical}, 29) — the head turning toward the source before the eyes have found it. This authored, spatially anchored auditory cue is the designed stimulus by which perception itself is steered. The \textit{spatial} channel directs the eyes; the \textit{affective} channel, the helping tone already gathered around the creature, colors what they find; and the player looks up to watch the goblin chop at the fruit stem until it comes loose and falls to the floor (06:52.8). What completed perception (A₁) deposits is not yet a plan but a charge: the persisting image of a thing that has just come free.

From that charged residue the image of a possible act takes form (A₂): \textit{a fruit I can pick up}. Unlike the outfielder's projection, which composes trajectory-residues into a determinate landing-point, this \textit{phantasma} is under-determined by design — nothing has said what the fruit is for — and what fills it out is the kinesthetic disposition the free play laid down. The settled sense that in this world the graspable thing is to be grasped, and the experimental play that has demonstrated you can give the goblins the acorns that are strewn about. The design has withheld all instruction, and into that withholding the player's own \textit{phantasia} supplies the act the design declines to name. This is the conscription the candle first practiced, now bearing weight: the world authors the occasion, the affordance, the potentiality; and, the player authors its sense and actualization.

What turns \textit{pick it up} into \textit{give it to the goblin}? The game's design takes this into careful consideration. If the player has not yet discovered that the acorns strewn about can be handed to the goblins, the creature itself supplies the impetus, through a pressure that runs along the \textit{shared} channel alone. Once the fruit is down, the goblin waits; it merely looks at the fruit and then at the player, no word exchanged, and it is left entirely to the player to read what it wants.

On that bare expectant nearness the player ratifies a judgment (A₃), and the judgment has a fine structure the method must slow down to see. Its cognitive orientation is the \textit{deliberative-discursive} mode — not a memory, and not an intellection of anything universal, but a prospective reading-forward of the scene: what does the creature want, and what would meet it? That orientation is one axis; the \textit{doxa} that rides it is another, and the two must be held apart, for the mode is only \textit{how} the player is turned toward the scene while the \textit{doxa} is \textit{whether} its appearance is taken as true. And the committal \textit{doxa} of the pass — \textit{this just fell, and the creature wants it} — is itself layered. It rests on a \textit{resonant doxa}, the standing reading a lifetime has settled: that a creature which turns to a thing, then to you, and waits, is a creature that wants. And it consists in an \textit{occurrent} ratification, the involuntary taking-as-true of \textit{this} creature, now, wanting \textit{this} — which does not wait upon deliberation but fires on the appearance itself. No line of text proposes the gift; the goblin's posture, read through the player's own \textit{phantasia} and against the player's own standing \textit{doxai}, does.

But what inclines that judgment is not cognition alone. Emotion, on this project's account, is the form desire takes under an evaluative appearance, and it enters the chain \textit{transversely} — coloring the tone at which perception arrives (A₁) and concretizing into a determinate feeling at the moment of judgment (A₃). The mood was set long before the fruit fell: the childish, unguarded delight the forest cultivates is the \textit{Stimmung} under which every perception in it is given, so that the chopping is heard, and the goblin seen, already within a warmth the design laid down. A mood does not itself concretize; what concretizes, formed within it and at A₃, is an empathy — a wanting to help the goblin — a \textit{philia} built by degrees from the curiosity of first sight into something like tenderness for a small creature that cannot reach its fruit. The emotion is not the player's spontaneous overflow; it is precipitated, and precisely. The design made the goblin endearing — shy, child-scaled, glimpsed and gone — and made its very approachability answer to the player's conduct: an abrupt movement reads as threat and sends it scampering, a gentle one draws it near, so that over the preceding minutes the player's own gentleness has been the instrument by which curiosity tipped into care.

The emotion does cognitive work, and the method can say what. It rests, as the \textit{doxa} does, on a resonant layer — that a small creature in want is a fitting object of care — and consists in the occurrent feeling of \textit{this} creature's plight; and once aroused it bends the judgment toward itself. An aroused emotion can impair the corrective faculty — not inevitably, but readily — and here, with nothing in the world to check it, it does, so the appearance it inflects is taken as true unopposed: the reading \textit{it wants this, and I would help} meets no skeptical counter — no intruding thought that the creature is a timed script — and the feeling, unopposed, deepens as the exchange proceeds. This is the pathetic loop, intra-event, distinct from any disposition the pass will sediment. Its magnitude must be read on three independent axes, and the reading is the point: this is an emotion of low hedonic intensity, low existential weight, and only moderate cognitive articulation — a quiet, warm, low-stakes feeling, nothing like the anxiety a rupture will later work. That the design reaches for so mild an affect is itself the design: a helping made to feel not momentous but natural. And it is this prepared feeling that sources the wish the next motion completes — the \textit{boulēsis} for the goblin's good is just this care given the form of a desire. The judgment that completes the chain's cognitive station is thus the player's own, ratified on a cue no stronger than a waiting body and moved by a feeling no stronger than tenderness. But the cue was written, the tone was trained, and the feeling was prepared.

So the player acts (A₄): carries the fruit to the goblin and sets it on the ground between them (07:13.3). Two channels carry the act, and they carry different things. The execution runs through the \textit{kinesthetic} channel, and runs there transparently: the carry is a continuous perceptual-motor adjustment, grip and footing revised step by step, each revision re-entering the chain at its image-station — perceived, imaged, judged, made — so that the short walk is itself a rapid series of miniature actualizations. The choice runs through the \textit{ludic} channel: to carry the fruit rather than drop it or turn away is a decision taken under repercussion, and refusal stayed open the whole way. Naming the cause fixes the freedom. The act is voluntary — its source lies in the player, not outside — so it is in no way \textit{bia}, the forced motion that moves a body against its desire; and among the seven causes of action it is agent-caused and desire-sourced, two of them working together: the settled disposition the free play had built (the habituated readiness, here, to take a graspable thing and offer it) and the wish that disposition serves — not an appetite for the fruit but a \textit{boulēsis}, a wish for the goblin's good, carried through to \textit{prohairesis} in the deliberate choice of a means, the offering. The freedom, then, has a precise location: it lies in the engaging, not in the outcome. The design has so arranged the scene — the chop, the fall to the ground, the patient wait — that of everything the player might do, this act is made salient, very nearly inevitable; what the design cannot supply, and need not, is the desire that moves the player through it. This is the producing-art and using-art division the method draws: the maker furnishes the act's \textit{matter} (the fruit, the goblin, the ground between them) and its \textit{form} (that a held thing may be given), while its \textit{efficient} and \textit{final} causes — the desire that moves and the end it serves — stay the player's own. Curation as against coercion names just that division: a world that hands over every cause of the act but the two an agent must supply, and so moves the player to its end without forcing the player to it.

The goblin's part seals the reading. It does not take the fruit at once; it waits some eight seconds, then bends, lifts it, and carries it off (07:21.7). That delay is the tell. The creature is not a fellow agent answering a gift but a scripted trigger firing on a timer — its actions ``are nothing more than ... programmed responses triggered when the user does something the program recognizes'' (Salvo, \textit{Veri(dis)similitude}, 36) — so that even as it appears to act on its own, it acts as a superficial agent whose motive is the programmer's, the hidden maker for whose efficient causation the goblin's ``wanting'' stands in. The player reads none of this. \textit{Phantasia} fills the pause as expectancy and hesitancy, and the script as intention: a timed retrieval lived as a creature's natural response. Out of that lived misreading an emergent story writes itself — \textit{I helped the little one who could not reach its fruit} — a narrative no author scripted, composed by the player's own play — what Calleja calls the \textit{alterbiography}, ``the ongoing narrative generated during interaction with a game environment'' (Calleja, \textit{In-Game}, 124), set against the scripted narrative the designers write in: the \textit{narrative} channel operating in its generated rather than its scripted register. What the exchange leaves behind is correspondingly real. Another helping tie further sediments an ethical disposition (\textit{hexis}) toward the goblin: the player is now one who has helped, and the world one in which help is taken up and answered. The \textit{shared} dimension now opens, and the union this project calls rhetorical incorporation has gained the beginnings of its second pole — an accepted invitation. At this point the player should have internalized the mechanics, and the design cultivated it deliberately through a tutorial whose only method was to let exploration teach. Read whole, then, the fruit exchange is the persuasive virtual ecology in miniature. Every cause but two is the designer's — the fruit, the goblin, the sound, the fall, the wait, and the rule that binds them — and the two that remain the player's, the desire that moves the act and the end it serves, are exactly what the design cannot author and has no need to, because it has built a world in which the player's desire and the designer's end fall into a single line. To curate an environment, on this evidence, is to arrange the conditions under which a free act becomes the likely one: to move someone to do what you would have them do by handing them a world, and never a word. This is more than a technique of games. Ordinary persuasion works on what a person believes; curation of this kind works beneath belief — on what the player wants, and beneath that, on the disposition from which wanting springs — so arranging the world that the desire the design intends arises as the player's own and issues, unforced, in the act the design foresaw. A rhetoric that operates on the springs of action rather than on the propositions of assent reaches the agent where argument cannot: it is, in the most exact sense, one way of reaching into and reshaping the soul.

\bigskip

\subsection*{The Disruption: Rupture and Carry-Over}

\subsubsection*{The Miniaturization: A Rupture That Re-Incorporates}

By the time the gold bell lies on the moss at the foot of the tree, the player has already come to dwell, and the ringing that follows draws its whole force from that fact. The opening minutes have done their quiet work — a fluent grip, a path learned by walking it, a goblin met and helped — a skill (\textit{technē}) settled into a disposition and a relationship begun, sedimented act by act, until the forest is no longer a scene surveyed but a place inhabited. The act that triggers everything is the tutorial's own disposition come home. A world that teaches only by exploration produces a player who probes; the bell is rung because ringing intriguing things to see what they do is, by now, who the player has become. Curiosity moves the act — what does this bell do? — a desire sourced in the player and in no goal the design has set; and that desire completes itself as a wish, to learn what ringing brings, whose means, the bell raised and sounded, is deliberately chosen. This is \textit{prohairesis} doing its quiet work — the deliberative completion of a wish, not a fourth kind of motive beside the three desires — and the end it serves is the player's own, proposed by nothing in the world. What the player cannot author is the consequence. Nothing on the bell and nothing in the wordless world declares what ringing will bring, so the charged image the player acts from — the \textit{phantasma} of a bell to be rung — is under-determined as to its outcome, filled out from the settled judgment that bells are meant to be rung. For a thing is never met as a bare object but as a function: Uexküll's stone, taken up on a path, ``becomes a carrier of meaning as soon as it enters into a relationship with a subject'' and ``had, we could say, a `path tone''' (Uexküll, 140), and Heidegger says the same of equipment, which ``is essentially `something in-order-to''' (Heidegger, \textit{Being and Time}, 97). The bell reaches the player already tuned to its use, a thing-to-be-rung, and the player rings it because that — here as in every world the player has known — is what a bell is for. The freedom is real but orchestrated and anticipated, exactly as it was for the fruit: it lies in the \textit{whether}, the choosing to ring, and not in the \textit{what-it-does}, which is foreclosed. If the episode were nothing more than a set-piece the player cannot refuse once triggered, like a cutscene, that foreclosure would be damning. What saves it is the character of what the ringing then sets off — an overwriting of the whole perceptual field, the forest at human scale replaced in an instant by another — and, decisively, what that overwriting does not touch: the player's standing as an agent within the world it transforms.

For the world then overwrites the image. Fireflies wake and swirl, light saturates to white, and the point of view drops as the world swells — grass risen into trees, the human scale of a moment before collapsed to the scale of an insect (07:48--08:14). For this the player has no fore-structure; the anticipatory \textit{phantasmata} derived from prior experience of bell-ringing is fulfilled (the bell makes a noise when rung) and then exceeded. Yet the disruption does not expel. Movement remains, gaze remains, a second ringing remains available. The player is transformed, not dragged from the agent's place to the patient's. The distinction turns on the desire that rang the bell. The player rang it wishing to learn what ringing would bring; the transformation is that learning delivered in full — not an outcome imposed against what the player wanted, but the very thing wanted, answered beyond all asking. Because the change fulfils the desire that occasioned it rather than overriding it, it is no forced motion, no cutscene seizing the body, but a designed and productive rupture. What it ruptures is the scale at which the player was incorporated; what it leaves intact is the incorporation itself, and the being-affected it works is of the preserving and fulfilling kind, perception stretched by wonder rather than damaged by force.

Set against the other breakdowns this world can stage, the miniaturization shows its true structure, for it is the deepest of them and the only one that re-incorporates. Each of the others discloses by breaking, and each breaks in its own way. The rock breaks by refusing the hand: reaching to take it and finding it unusable, the player meets a thing gone un-ready-to-hand, and in that refusal the world discloses that it need not yield. The apparatus breaks by ceasing: when the headset slips or the power cuts, the virtual world — which exists only while it is rendered — simply passes out of being, and in that interruption the rendered discloses its dependence on the renderer that holds it. Each such break discloses something true, and each pays for its disclosure with the dwelling itself — the rock falls out of use, the world winks out. The miniaturization breaks more deeply than either, for it remakes the very scale of being-in-the-world; and yet, alone among them, it breaks without expelling, its disclosure undergone from within a dwelling it does not end. Before going on, it is worth pausing over how easily the design could have declined this moment, for the avoidance it did not choose is the counterfactual that gives the bell its stakes. A world committed to seamlessness — to a verisimilitude so total that nothing in it ever exceeded the player's fore-structure — would by that very commitment foreclose the disclosure the bell delivers, leaving its inhabitant absorbed, fluent, and forever unshown the rendered character of the rendering. The miniaturization is the design choosing to break its own seamlessness at the very point where the illusion was complete. Everything the rupture reveals depends on that choice.

This is also why the shock lands as wonder and not as dread. Because the forest had set its mood in advance — that childish, unguarded delight — the perceptual moment into which the transformation arrives is already colored, and what the world's giving-way attunes is anxiety in its disclosive register rather than its threatening one: the vertigo of a ground withdrawn, recovered almost at once into the poise of an agent still standing. Anxiety (\textit{Angst}) here individualizes the player without leaving it, in the end, not at home in the world — the ground drops away and is regained. Rendered as motive, the episode then runs on a single ratio. The player attributes the change not to any agent within the world — the goblins, agents in the fruit exchange, do nothing here — but to the scene itself, the bell and the world it transforms: \textit{the scene did it}. And the scene, in transforming, transforms what the agent is and can do: a Scene--Act relation, the container reshaping the contained (Burke, \textit{Grammar}, 443). Two terms go concealed in that attribution, as a banished term goes on doing its work unseen (Burke, \textit{Grammar}, 21): the player's own settled readiness to engage with virtual objects---which the tutorial built---and to ring a bell---which previous experience with and knowledge of bells has established---and the maker who wrote the rule---the programmer, the hidden Agent for whose efficient causation the magic stands in.

Read as union, the beat is the height of the whole arc, and Calleja's criteria for incorporation — the R1--R7 test that tells it apart from mere involvement — give the rubric by which the height can be named rather than merely felt. Everything the forest had built — the fluent body, the inhabited space, the wonder, the dawning story — is co-active at once, and rhetorical incorporation reaches here what R6, the temporal apex, names: the culmination of its climb from off-line anticipation to absorbed, moment-to-moment play (Calleja, \textit{In-Game}, 178). The rupture then puts R7, intrusion-integration, to its limit case — an intrusion that, rather than breaking the absorption, is taken up into it. Calleja's own instance is a call for help from the room, which, far from dispelling incorporation, ``will most likely intensify the sense of incorporation because it fosters the absorption of the game environment into the player's consciousness more decisively by blending stimuli from the game world and from the immediate surroundings'' (Calleja, \textit{In-Game}, 172); the miniaturization is that blending driven to its limit, the intrusion authored from within. By that criterion the miniaturization earns a description that would otherwise be lyric: the world's most violent move deepens rather than dispels belonging, the dissimilar shock re-incorporating the player at a new scale instead of casting the player out. World-disclosedness (\textit{Welterschlossenheit}) — disclosed, for Heidegger, equiprimordially with the others who share the world and with one's own existence (Heidegger, \textit{Being and Time}, 176) — is not undone but re-pitched: the designed surround is disclosed again, smaller and stranger, and again as a place to be dwelt in. This is \textit{veri(dis)similitude} made event. A world verisimilar enough to be entered through fluent dwelling is broken open by a possibility dissimilar enough that the breaking carries ontological force, and the force is disclosive: what the miniaturization unconceals is the contingency of a constitution that ordinarily withdraws from notice — that the player's scale, orientation, and standing in a world are not given but rendered, and can be rendered otherwise. The bell that worked the rupture becomes, now, a way through. Rung again, it teleports the player, who moves among the small world's places once only observed and, in so doing, moves toward the others who dwell there.

\subsubsection*{The Tavern: Co-Disclosedness at Its Threshold}

Rung from within the miniaturized world, the bell no longer transforms; it carries, setting the player down with each ringing in a new place among the tiny world's locations. One of these is a tavern — a drinking-hall built for a crowd, its long tables and benches and scattered mugs sized for company, lamplit and warm — and two goblins are at home in it, one asleep, one bent over a drink (08:56--09:16). Nothing here is staged for the player's coming. The room and its dwellers go on with a life of their own, and what the place discloses, in disclosing that life, is the goblins' \textit{being-with-one-another}: a way of living together already under way before the player arrived among it, and not pausing for the arrival.

To be set down in that hall is to undergo one of the founding disclosures of \textit{being-in-a-world} at all. Because its dwellers come forward not as objects met but as others already living together, the world the player now inhabits shows itself as a with-world, and to enter their place is to receive the ``disclosedness of the Dasein-with of Others'' (Heidegger, \textit{Being and Time}, 160) that belongs to every being-with. The disclosure runs both ways. At some of the bell's destinations a goblin lifts its head and looks — a small thing, and the whole of it — for in the returned look the other is given as an other and the player as one who is there to be seen, co-present in a shared place. What this project names co-disclosedness (\textit{Miterschlossenheit}) opens here not as a doctrine but as that mutual look across a lamplit room, and the \textit{shared} dimension of involvement, which the fruit exchange first opened as a single seeded tie, now stands at its structural ceiling: a whole social world, present and alive, with the player invited and admitted into its midst.

Read as motive, the tavern works by what Burke names consubstantiation, a shared substance built out of shared act. ``A way of life is an acting-together'' (Burke, \textit{Rhetoric}, 21), and in acting together the members come to hold the sensations and attitudes that make them of one stuff. The goblins' communal drinking is just such an acting-together, their consubstantial way of life made visible; and the warmth of the inhabited scene works on the player as the sympathetic imagination works. Identification, in the account Burke adopts from Hazlitt, has us ``deriving from imagination both our ideas of self-interest and our sympathy for others'' (Burke, \textit{Rhetoric}, 83) — it is the liveliness of a present image that makes its object matter as one's own. So the lamplit hall, rendered vivid, comes to matter to the player and to solicit belonging, before any deed is done among them. What the player cannot do is enter the substance. Admitted to the hall and co-present in it, the player remains ``both joined and separate'' (Burke, \textit{Rhetoric}, 21) in the thinnest register the phrase will bear — joined as a guest is joined, present and unturned-away; separate as a stranger is separate, unaddressed, no member of the ``we'' that is present in the tavern.

For the joining stops at the look. The stopping has a precise grammar in the chain's own terms. Though the goblins drink and sleep and raise their heads, they never act with the player — never hand a cup, never make room at the table, never take up a common task — and so the one thing that would consummate the bond is never put before the player at all. A committal judgment requires something to ratify: where the fruit exchange offered a cue on which the player could \textit{take-as-true} that \textit{the creature wants this}, the tavern offers no cue on which anyone could \textit{take-as-true} that \textit{we are doing this together}. The ``we'' is not proposed and declined. It is never proposed. Their motion is programmed responsiveness, scripted by the hidden maker whose efficient causation the room's life stands in for. What fills it out as a life — the sense that someone is genuinely at home behind the drinking — is supplied by the player's own \textit{phantasia}; the goblins hold no settled disposition of their own to extend, so nothing in them can close upon the player as a fellow. The acting-together is given to be witnessed and not to be entered — whether by deliberate authorial choice or by the plain limits of the NPC system, the analytic result is one and the same — and the helping tie the fruit exchange seeded is thus called due and left standing: honored with a look, extended by nothing. This is co-disclosedness given as a preview, the disclosure of being-with isolated, deliberately or not, from its enactment.

So the small world holds the player at two depths at once. World-disclosedness is sustained and even widened — the bell that re-scaled the surround now ferries the player through it at will, and the dwelling deepens with each place entered — while co-disclosedness arrives only to its threshold, the co-world and the mutual look disclosed, the acting-together kept just beyond reach. These are not two unions but one incorporation disclosing along its two poles, one in substrate and different in being; and, \textit{Gnomes \& Goblins}, by opening the second pole exactly to the edge of consummation and no further, lets that single incorporation be seen as a consummated collective never could. That the co-world shows at all, even with its enactment withheld, is the work of the same disclosive power the rupture exercised in the breaking — the power of the rendered world to make briefly visible what it silently constitutes.

\subsubsection*{Carry-Over: The Unconcealing and What It Leaves Behind}

The rupture's deepest work is not the change of scale but the change of sight. A world dwelt in withdraws from notice — the fluent body, the inhabited place, the standing from which one acts all do their work precisely by going unremarked, the hammer disappearing into the hammering — and what the miniaturization does is break that transparency for an instant and, in breaking it, unconceal it. Having come to dwell at human scale, the player is handed, without warning, a world in which grass stands like trees; and the gift is not the new scale but the sudden visibility of the old one \textit{as} a scale — one rendering among possible others, contingent where it had seemed simply given. Just as the tool laid bare in its failing discloses the whole equipmental web that had silently carried it, so the breakdown of the player's settled scale lays bare the ordinarily hidden constitution of \textit{being-in-a-world}. The mood in which this arrives is the analysis's confirmation. Anxiety here is not a failure of absorption but its disclosive limit, the affect in which the player's enworldment, having lost its footing, comes into view as enworldment. Anxiety individualizes. It draws the one who undergoes it out of the absorbed, anonymous everyday and sets the bare fact of \textit{being-in-a-world} before the player; the miniaturization works that individualizing stroke in little, and works it as wonder rather than dread, but the structure holds. What had been lived is, for a moment, seen.

Seen — and the seeing changes what can be said of the union. The union in question is rhetorical incorporation itself, the player's absorbed dwelling in the designed world; until the rupture broke its transparency it was only lived, disclosing nothing of its own character, and only now, held up to view, can it be named for what it was. Named precisely, it was a fallen union — a dwelling at once genuine and fallen — and the fallenness admits of statement. In Heidegger's terms the player has fallen into the world it dwells in: ``Dasein has, in the first instance, fallen away [abgefallen] from itself as an authentic potentiality for Being its Self, and has fallen into the `world''' (Heidegger, \textit{Being and Time}, 220). The world's significance comes pre-articulated, settled in advance as the anonymous ``they'' settles the sense of the everyday, so that the player takes on trust a meaning another has authored. The player is drawn ever onward by curiosity — which for Heidegger ``is concerned with the constant possibility of distraction'' and ``has nothing to do with observing entities and marvelling at them'' (Heidegger, \textit{Being and Time}, 216), a fleeing into novelty rather than a thirst to understand, and just what the miniaturization's wonder for a moment breaks; and, beneath both, the player is constituted as a functional object of the programmer's design, moved through a possibility-space foreclosed in advance: every choice, like the chord already lying in the piano before the child's hand finds it, available only because it was built in (Adorno, \textit{Aesthetic Theory}, 32). What the player takes for discovery is the actualizing of what was already laid down — as, for Aristotle, the geometer's constructions are not invented but found, ``the potentially existing relations are discovered by being brought to actuality'' (Aristotle, \textit{Metaphysics}, 1051a21--33); the player draws out only the possibilities the design has drawn in. Incorporation left to itself is this fallen union, and nothing in its felt quality betrays the fact. That is what fallenness means. The rupture is what can convert it. By making the constitution visible, the unconcealing opens the possibility of a dwelling no longer merely undergone but owned — a critical, individualized inhabiting in place of an absorbed one. The productive rupture is thus not the enemy of incorporation but its completion: it converts a union the player had only suffered, blindly, into one the player can see, judge, and take responsibility for.

The disclosure does not end when the headset comes off. After a long duration in a rendered, VR world, the return to the non-virtual world of objective reality is not immediate: the exit brings ``the odd sensation that I remained tethered to the computer,'' depth-perception was briefly unreliable, and the disturbance held. After which, ``it took around 15 minutes for my phenomenological and ontological orientation to recalibrate to my `Being-in-the-world' of objective reality'' (Salvo, \textit{Veri(dis)similitude}, 58 n. 15). For a quarter-hour, that is, the actual world wore the strangeness the virtual one had shed — the eye reaching for a scale no longer there, the body carrying an orientation the room would not answer. The disposition sedimented across hours of dwelling pressed, for that interval, upon the disposition it had left behind. This is the actualization chain closing its longest loop, and the deposit it carried home can be typed. Part of it was the \textit{technē}-fluency of the dwelling itself, the trained reach and read that made the small world navigable. Part of it, seeded at the fruit and honored in the tavern, was the beginning of an ethical bearing toward the world's others. It was the first of these — the bodily attunement — that visibly outlasted its world. Over hours of play the body had learned the small world's scale: how far to reach, how to judge distance, what a step would cover. These trained expectations did not switch off with the headset. For the length of the recalibration they went on operating against the actual room, the eye still gauging distances by a scale no longer there, the hand still set for objects the room did not hold. The completed acts had settled into a standing disposition (\textit{hexis}), and that disposition outlasted the world that formed it, carried back into objective reality as a residue not merely psychological but phenomenological and physiological — a trained way of perceiving and moving, briefly imposed on a world it no longer fit. 

This transient supplanting of the player's dispositional ground---an alteration of settled orientation deep enough to bend how the actual world was met on return---demonstrates the rhetorical ontological nature of human \textit{being-in-the-world} in real time; how we are constantly, inevitably, attuned by the environments within which we participate and come to dwell. Every experience leaves a residue that shapes our future, lays down a dispositional ground from which we encounter the world. This resonance does more than shape the future: in sedimenting, it draws the three tenses of a life together. The past is retained as the disposition from which the world is now encountered, and that retained past, borne as a present readiness, is what casts us forward — conditioning what we will next perceive, wish, and do. This holding-together of the absent with the present is the work of \textit{phantasia} itself, which stretches the soul across past, present, and future; to lay down a \textit{hexis} is, in just this sense, to co-temporalize — to make one's having-been the standing ground of a future not yet come. It must be held within its limit. The limit is the kind-distinction this project has carried from the start. The virtual is distinct in kind from the non-virtual; no depth of carry-over converts the rendered world into an unrendered one, and no recalibration makes the player a new kind of being. What the carry-over alters is the \textit{how} of a dwelling and not the \textit{what} of a substance — an alteration (\textit{alloiōsis}, \gk{ἀλλοίωσις}) so deep it strains toward a remaking of character and yet, the virtual keeping its kind, cannot become one. Fifteen minutes of disturbed orientation will not carry a conversion narrative, and the analysis should not ask it to; what it carries is exactly what the method predicts, a dispositional residue proportioned to the dwelling that produced it, real enough to measure in the body and bounded enough to keep the kinds distinct.

Here, then, the method earns its stake, and the section can close on what the whole demo has demonstrated. If a designed world can drive the actualization of desire to a substantial union with itself and its others; if that union deepens through a wordless pedagogy that teaches by exploration alone; if the world's deepest designed event can break the union's scale without breaking the union, and in breaking unconceal the rendered character of the whole; and if what sediments through all of it carries past the world's edge into the player's standing in the actual — then to design such worlds is to design dispositions, a rhetoric that works not on belief but on the ground from which belief and action set out. The miniaturization is a small and benign instance, a moment of wonder that leaves the one who undergoes it faintly re-pitched toward the world; but the principle it instantiates is not small. If a designed world sediments dispositions that outlast it — and the carry-over shows that it does — then in building such a world one builds, in part, the player who will inhabit it: the perceptions that player will have trained, the desires readied in advance, the bearing carried back out into the actual. That is what it means to call the rendered world a maker of who its inhabitants are becoming; and it is why the design of such worlds cannot be thought as a merely technical craft, but must be thought as an ethical one, answerable for the dispositions it lays down.

\bigskip

\section*{5. \textit{Red Dead Redemption 2} (§1: The Tutorial)}

% ============================================================
% RDR2 SECTIONS 5-6 spliced from RDR2-INTEGRATED-DRAFT-v2.md (2026-07-09).
% Provenance markers (INS-/FIX-/bank/coverage/ledger/gauntlet) stripped.
% OPEN ITEMS carried over from the draft (NOT resolved by this splice):
%   N1  investigative-stance recast of the Sonny section (plan §7b).
%   N2  apparatus purge (\"the frame record / footage / my recording\").
%   ****** newspaper-report YouTube citation (2 footnotes, choke/spare + heist).
%   ****** Sonny revisit-dialogue appendix (footnote).
%   ****** failed-screen image + Sonny-approach video clip (see markers below).
%   bell-toll timing (choke/spare, clip 3a 11:31-11:45) awaits frame verify.
%   FIX-D4 / FIX-D6a/b were applied in the draft; foundation §1-§4 untouched.
%   length: tutorial ~9.3k words; further cuts await user-directed targets.
% Abstract note updated (\"§2 ... pending footage\" -> \"included as Section 6\").
% ============================================================

\subsection*{The Critical Field}

Scholarship on \textit{Red Dead Redemption 2} is young, compact, and unusually varied for its size. The game shipped in October of 2018; within five years it had drawn sustained readings from game studies, rhetoric, history, media studies, ecocriticism, performance studies, and conservation science, and the field's two enduring registers were established in its very first year. In 2019, Nicolas Bagnoli read the game through Bogost's procedural rhetoric and Althusser's apparatuses, arguing that the game ``managed to include a contemporary worldview'' whose delivery is procedural rather than scripted: ``several arguments are expressed through the processes of the game'' (40).
In the same year, Andrew Westerside and Jussi Holopainen set ideology aside for place, coining \textit{gameplace} to name ``the affective relationship between place, experience and play'' and locating the game's power in its lived, inhabitable world (1).
What followed was fan-out rather than consolidation. Where Juho Tuominen pressed the Althusserian line to its strongest form, reading the series whole as ``mainstream but still radically progressive in content'' (60), Esther Wright answered from the archive of Rockstar's own marketing, questioning ``the value of any criticism that merely scratches the surface of deep-rooted, structural violence, while elsewhere profiting from the discourses historically associated with it'' (par.\ 9).
Measuring what the world teaches incidentally, Edward Crowley, Matthew Silk, and Sarah Crowley found that players ``correctly identified more species'' of wildlife than non-players in a quiz of five hundred eighty-six participants (1229, 1234). Hanna Ventomäki concluded that ``there is no single form or role that nature is reduced to'' in the game's world (67). Iain Donald and Andrew Reid distilled the authenticity question to its axis --- ``[a]uthenticity is less about getting the past completely accurate and more about getting the feeling of period and timeline correct'' (18).
John Vanderhoef and Matthew Thomas Payne read the game's deliberate slowness as a cultural politics of time, a design that ``betrays the dominant chrononormativity that governs big-budget video game design practices and neoliberal capitalism more broadly'' (``Once Upon Some Times in the West''). Paolo Ruffino followed the player community's video paratexts as attempts to ``remediate (in the double sense of restoring upon and healing)'' an agency the game withdraws (345); Heather Moser recast the player as a performer who develops empathy through the long inhabitation of a role. Closing the window in 2023, Amalia McEvoy joined reader-response theory to procedural rhetoric to argue that the game ``persuades players to make moral choices while they learn about American history through the game's environment'' (3).
Beneath this variety run two constants. Nearly every strand, whatever its discipline, routes through the honor system and through the frontier myth: the mechanic that scores the player's conduct and the story America tells itself about its own violence.

The field's deepest division, however, is not thematic but methodological: whether the game is to be read from inside the playing or from outside it. One camp reads experientially, from the first person of play. Moser's account is the fullest: ``[w]ith Arthur, players push past the lifeworld boundary, stepping inside the life and world of Arthur while bringing along their own lifeworld to create a doubled structure, their experiences and world become entwined with Arthur's'' (14) --- a phenomenology of felt union that Westerside and Holopainen's \textit{gameplace} shares at the level of world rather than character.
The other camp reads structurally, from outside and against the experience. For Tuominen, the player ``is rendered as a subject, who is interpellated through different institutions of a game'' and made ``to absorb its ideology'' in the act of progressing (2); for Wright, even the game's apparent self-critique is a branded commodity.
The two camps rarely cite one another; across the eleven studies gathered here there is almost no internal citation at all. The disagreement has never been argued out. It stands as the field's open methodological question: if the first camp describes a felt transformation it does not explain mechanically, the second describes a machinery of influence it does not follow into the player's experience.

Where the registers do collide, they collide over agency. On one side, the game's freedom is a fiction maintained and then conspicuously broken. Ruffino argues that \textit{RDR2} ``denies Arthur a specific kind of agency, typically assigned to the white, able-bodied male characters of videogames'' (347). The tuberculosis that kills him arrives as a consequence without a visible cause. The player community's futile search for a cure --- ``there is no cure'' (347) --- becomes the refusal to accept the withdrawal. Vanderhoef and Payne find the same withdrawal distributed through the game's tempo, its waiting, walking, and watching; Tuominen, following Mitchell, holds the open world's freedom to be largely illusory, its choices already made (25).
On the other, the freedom is genuine but fenced. McEvoy maps Salen and Zimmerman's definition of play --- a free space of movement inside a more rigid structure --- onto the open world and its coded social systems, so that the player is at once empowered and disciplined (12), and Moser's performer experiments freely within the bounds a moral world sets. Because the event at their center requires both at once --- a freedom real enough to be exercised, and a design able to seize it --- the analyses that follow cut across this axis rather than choosing a side.

The honor system, the field's shared hub, divides three ways rather than two. McEvoy reads it as moral pedagogy and as the game's rhetorical center --- ``the rhetoric of the game lies within the function of the honor system'' (32--33); Tuominen reads the same reward-loop as an interpellation device, progression itself absorbing the player into an ideology. Bagnoli, borrowing Sicart's verdict on the earlier \textit{Red Dead Redemption}, reads it as arithmetic --- ``a polarized system of values'' that models ethics on a single axis and inscribes the moral imperative ``in the character'' rather than in the player (42; 48).
Notably, the system's two sternest critics also document its designed exceptions: killing members of the Ku Klux Klan costs no honor (Tuominen 53), and in Bagnoli's account it even grants some --- a fact that ``says a lot about Rockstar's opinion on the KKK'' (72). The observation matters more than either study makes of it. The coverage of the moral machine --- what it scores, what it rewards, and what it declines to see --- is itself a designed and therefore expressive variable.

For all this attention, the scenes this study reads closest are nearly untouched. The Valentine bar-fight, the game's first staged brawl, appears in a single study: McEvoy reads its ``dual purpose'' as combat tutorial and as first disclosure of ``the depth of Arthur's dishonorable character'' (22--23). The Sonny event, the game's darkest optional encounter, likewise appears in a single study. Moser treats it in one paragraph, in service of a different argument, as the limit case of what a loaded save can undo --- for Arthur everything, for the player nothing (15--16).
The scene-level ground, in other words, is nearly open, and the field's central impasse is precisely what a method built on the actualization of desire is positioned to mediate. An account that follows a designed world into perception, through \textit{phantasia} and the ratifying opinion of \textit{doxa}, to the completed act of \textit{orexis}, is a structural account \textit{of} the experiential: it explains, in mechanism, the felt transformation the phenomenologists describe, and it follows, into the player, the machinery of influence the ideological critics anatomize. Because the account runs on dispositions --- on \textit{hexis}, which does not differentiate between virtual and non-virtual experience, both being real experiences of different kinds of artifacts --- it can say what neither register can say alone: how the designed becomes the felt, and how the felt walks out of the game.


\bigskip


Where \textit{Gnomes \& Goblins} withholds instruction and lets the player discover, \textit{Red Dead Redemption 2} withholds freedom and scripts the discovery. The two cases invert one another. The virtual-reality world trusts the player to find meaning in an open surround, the candle and the goblin and the fallen fruit each yielding only to a hand that works at them and none of it explained. \textit{RDR2} instead places the player in a prepared corridor and runs a full course of instruction through it, mechanics and ethics alike, before the player has access to the open world.

That argument is built upon a \textit{hexis} (\gk{ἕξις}) — or rather upon two of them, installed together from the first interactive clip. The first is a \textit{technē} \textit{hexis}, the craftsman-outlaw's: horsemanship, gunplay, the tracker's eye, and the heist-man's deliberate hand, each installed across three hours of structured play and then left to run untutored. The second is an ethical \textit{hexis}: the standing disposition to hold the gang's survival as the governing good, to dwell in its camp and be of use there, and to act as the kind of man the honor system is built to measure. It is laid down in Dutch's grief speech before the controller arrives, deepened by camp labor and morning greetings, and put on open display across the Valentine sequence. The claim here is about this artifact rather than about \textit{hexis} as such — a \textit{technē} can certainly be cultivated apart from any ethical disposition, a courage-free marksmanship or a loveless craft. In this tutorial, however, the design never separates them: every skill is taught inside a situation that carries ethical weight, and every ethical situation is reached only by exercising the skills.

Three times across that cultivation the governing desire — the \textit{orekton} (\gk{ὀρεκτόν}) that orients the whole — is raised to higher ground. Blizzard survival names it first: warmth, food, and shelter for the gang, the bare needful goods against the storm. Standing and belonging name it second: being of use inside Dutch's world, earning one's place, mattering to the people who followed Dutch from Blackwater. The outlaw economy names it third: the fence unlock that turns a one-off robbery into a standing livelihood, the gang's prosperity made the shared object of desire. Each horizon carries the settled dispositions forward and adds a layer the next will draw upon, so that the player who begins as a patient in a cutscene — an unnamed body in a blizzard, led by the reins — is the very player who, hours later, acts in silence with a bow drawn on legendary prey and needs no narrator to name what is happening. It begins, before any mechanic, with a name.

The analysis reads six sequences from these opening hours, each chosen because it contributes something the others cannot: the opening cutscene, where the ethical horizon is set before the player can act; the leash and the Colter corridor, where the control scheme and the first moral forks are installed together under a hard constraint; the guided hunt, where Charles narrates the actualization of desire stage by stage and makes the chain legible; the Cornwall train robbery, where the acting is collective and the governing desire shifts from survival to belonging; the camp at Horseshoe Overlook, where repetition rather than scored action does the depositing; and the Valentine trip, where the honor system's calibration is put on open display. A closing section gathers what the whole sequence has built.

\textit{(Clip numbers below are the raw capture indices of the playthrough, not a continuous sequence; the analysis leans on the named beats and cites a clip only where a timestamp anchors a quotation.)}

\bigskip

\subsection*{The Opening Cutscene: Identity Before Embodiment}

The chapter opens not on a player but on a spectator: for nearly six minutes the controller does nothing, the game is a film, and the player watches. This is the tutorial's first lesson, and it names itself as no lesson at all — for before the first mechanic is installed, before a horse is ridden or a rifle raised, the cutscene builds the ethical horizon within which every mechanic will afterward be used, planting a \textit{hexis} of gang loyalty as the orienting good before the player has had any occasion to act upon it. What the section establishes is the flat screen's founding reversal: identity is given before embodiment, the name before the deed.
The one scholarly account written from inside this position corroborates the sequence: Moser, recasting the player as a performer in the theatrical sense, observes that ``The first-time game player, on the other hand, is still party to the perspective of an audience member, discovering and learning information as it presents'' (Moser 25) — the spectator's seat is, on her reading, the first station of a rehearsal, occupied by the one who will later take the role. The cutscene's patiency is that station held deliberately long.

The opening frames set the condition of patiency. The Van der Linde gang flees through the Grizzlies in a blizzard — on horseback, on foot, carrying its wounded — and the player watches without acting, made patient by the medium's own conditions. This is \textit{bia} (\gk{βία}) in its most total and least felt form, not the sudden seizure of an ambush but the structural withholding the cinematic form itself enacts. It arrives, moreover, where the screen supplies least. What the blizzard delivers — white static, wind, indistinct shapes laboring through snow — is very nearly nothing, and where perception is given so little, \textit{phantasia} must constitute the experience rather than merely accompany it: everything the player registers of this world, the felt cold and the gang's desperation and the depth of the ravine they cross, is supplied not by the screen but by the player's own \textit{phantasia} from the backlog of prior perception. The screen provides the cues; the player supplies the world.

Into that primed reservoir Dutch van der Linde speaks. Between 04:11 and 05:01, in a speech the gang does not interrupt, he names the dead — ``Davy… Jenny…'' — grieves them, and rallies the living: ``Stay with me. We ain't done yet'' (clip 1, 05:18). What the speech installs, before the player has moved a thumb, is the ethical \textit{hexis} of gang loyalty: the disposition to hold the gang's survival as the governing good, to count its losses as losses that matter, and to follow Dutch because Dutch is the one still standing and still steering. ``There was nothing more you could have done'' (clip 1, 03:42) — the exculpatory frame opens here and will run the length of the campaign, attributing every loss to circumstance rather than to any agent's failure, to the storm and the Blackwater catastrophe and a world arrayed against the gang. This is the Burkean scene–act ratio running for the first time, the scene term absorbing the moral weight that might otherwise land on the agent; the player is handed the lens before ever being handed its object.

Dutch's charged face, the gathered gang behind him, the grief and the rally together compose the \textit{phantasma} from which the player's opinions about the gang will take their orientation: every \textit{doxa} the player later forms about loyalty — that the gang's survival matters, that its losses are losses, that Dutch is the one to follow — refers back to this image. It is composed almost entirely by \textit{phantasia}. A speech on a flat screen gives image and audio and nothing to do; Dutch's authority, the gang's need, the whole moral economy of survival are constituted by the player from a voice and a face on glass. The master \textit{phantasma} of Dutch as undisputed leader is built this way, and built before the player has done anything at all.

One detail seals the reading. At 05:47 Dutch says ``We've got some work to do'' and the controller becomes active — but Arthur has already been named, already addressed as ``the only one I can rely on to stay strong right now'' (clip 1, 01:13). On the flat screen identity precedes embodiment: the player is handed a name before being handed anything to do, where a head-mounted world would place a body in space and leave it to orient itself from the ground of sensory presence. Arthur exists before the player inhabits him; the ethical horizon exists before the mechanic does. The off-screen Blackwater catastrophe drives the point home — ``What really went down back there on that boat?'' / ``We missed you. That's what happened'' (clip 1, 06:00) — for the originating event, the job that killed Davy and Jenny, is never shown, and \textit{phantasia} must constitute it entire from dialogue and tone and the grief inhabiting Dutch's face. The withholding is a narrative method: the game hands the origin to the player to author, so that the player begins already inside a crisis, already mourning people never met, already indebted to a man the player has only watched speak. What the cutscene has deposited, before a single mechanic arrives, is the beginning of a disposition toward this world and the people in it — the orienting good against which every skill to come will be exercised. The controller arrives second.

\bigskip

\subsection*{The Leash and the Colter Corridor: Instruction Under Fire}

The clip that follows the cutscene runs twenty-nine seconds, and it is the tutorial's most important negative space — important not for what it does but for what it withholds, and for what the withholding installs. The player rides out of Colter in a column with the gang, into a screen that is near-total white. ``This goddamn weather,'' a companion says; ``Been two days or more like this now'' (clip 2, 00:00). The visual field is reduced to a blizzard that is nearly nothing, so that depth and terrain and the cold weight of the air and the gang's silhouettes resolving from the white must all be constituted by \textit{phantasia} from the player's own reservoir of hard mornings. Where a head-mounted display would deliver stereoscopic depth and the vestibular confirmation of a body in space, the flat screen gives white static and wind; the world is carried, almost entire, by \textit{phantasia}.

At 00:11 the player strays from the scripted trail, and the world declines to yield. ``I think you're going the wrong way'' — the sentence falls, companion and NPC and the voice of the designed path at once, and the player is turned back. The rebuff has teeth. A player who keeps riding the wrong way is warned, and a player who persists past the warning is given a game-over screen: ``Failed --- Dutch was abandoned'' (clip 2, 00:27).

% ****** OPEN (image): insert the "Failed --- Dutch was abandoned" screen (leash)

This is \textit{bia} in a hard form, not an ambient nudge: a design that achieves its goals by enforcing a linear, structured sequence of events, with little room for genuine deviation. The environment participates in the enforcement — the storm, the trail, and the companion's voice acting together as what Rickert, describing how a world takes part in the doing, calls ``transhuman powers manifesting themselves as the land's affordances and recalcitrances'' (Rickert, \textit{Ambient Rhetoric}, 243). No visible wall is drawn and no menu explains the rule; the world closes every way but Dutch's way, and ends the game if the player refuses the lesson. The lesson accordingly lands in two registers at once. On its surface it is navigational — stay with the column. Beneath that, the tutorial has drawn its first boundary of possibility: in this world, at this stage, action is what the design says it is, and the player's part is to move well within it.
The slow-time scholarship has read this same ride, and read it as loss: Vanderhoef and Payne hear the announcement in it — ``the game underscores through its initial horseback excursion that nothing will happen in the game `anytime soon.'{}'' (Vanderhoef and Payne, ``Travel and Stillness'') — and find in the whole design ``a game whose design facilitates, and, in fact, demands players slow down and wait'' (Vanderhoef and Payne, ``The Critical Potentiality of Slow Game Time''). Tuominen, following Mitchell and speaking of the open world at large, presses the same diagnosis to its structural end: ``this freedom is largely illusory and in fact, a player often has a little agency to determine a game world's events and outcomes'' (Tuominen 25). Both are right about the mechanism and, on the account developed here, wrong about its yield. What one register reads as agency removed and the other as freedom faked is, run through the chain, a deposit: the rebuff does not subtract movement so much as lay down the \textit{hexis} of being-led on which the whole training will build.

What settles across these twenty-nine seconds is not a skill. No A₄ completes, nothing is finished, the \textit{orekton} is posited but not satisfied. What settles is something more basic and more lasting — the \textit{hexis} of being-led, the disposition of path-adherence, the trained readiness to stay with the column and go where the column goes. The deposit is privative, not a positive capacity but a negative one, and the design's two instruments work at different depths. The blizzard is the soft instrument: by starving the visual field it encourages the player to stay within range of the column, since nothing else is visible and nowhere else is legible. The failed screen is the hard instrument behind it, a limit most players will never touch but which any player who probes the edge will find absolute. Both serve the same later payoff. The player must feel the leash to feel the moment it is loosed, and the player who has tested the hard limit and been turned back will feel the release into the open world more profoundly still.

After the leash comes the corridor, and the corridor makes the chapter's widest claim — that \textit{Red Dead Redemption 2}'s tutorial installs not one \textit{hexis} but two, not in sequence but together, in the same clip and the same thirty-eight minutes of scripted play. The Colter sequence threads the full control scheme — mounted movement, free-look, combat, inventory, melee grapple, horse handling, tracking, Eagle Eye, weapon selection — through a narrative that requires each mechanic to justify itself as something the story needs, so that no menu announces itself and every skill is introduced by the world demanding it. The horse is learned because the gang must move; the gun because the O'Driscolls attack; the satchel because the camp needs supplies. Instruction of this kind is possible only because the sequence is restricted. By limiting the possibilities for action, the design condenses the player's focus to the one mechanic the moment calls for; had the world simply opened and left the player to their own devices, nothing could guarantee that need and instruction arrive together, and the same mechanics would have had to be taught as menus and prompts — announced rather than encountered. The constraint is not the price of the narrative grounding but its enabling condition, and the efficiency with which the \textit{technē} sediments here is that constraint's direct yield: each demand arrives once, at the moment the story makes it intelligible, and is rehearsed immediately under pressure.
Moser catches the same design from the actor's side: the opening hours run ``a tutorial disguised as small missions for him to complete so players can learn how to use the controller and button actions. Just like the actor, players are given techniques and tools to perform'' (Moser 24--25) — instruction costumed as story, the player equipped the way a rehearsal equips a performer. But the instruction runs on two tracks at once, and the second track is ethical: alongside each skill, the world deposits a disposition. The player who is taught to shoot is at the same time being taught whom shooting is for; the player who is taught to loot overhears what the looted man's marriage was worth.

The corridor opens with movement, and with movement's constraint. ``Stay close. We'll do our best to stick to the trail,'' a companion says; ``Careful over this bridge here'' (clip 3, 00:28–00:30). Companion patter scripts every footfall, and navigational agency belongs to the gang and not the player. The leash of clip 2 returns now in narrative form, the only way being Dutch's way, and the \textit{hexis} of being-led deepens through story where before it deepened through rebuff. At 04:17 the camera opens — Dutch places Arthur in the cattle shed and the player finds the first mechanic of the gaze, the free-look, the flat screen's surrogate for the turned head — and where a head-mounted player turns to look with the whole body, here the look is a stick-driven symbolic act, the whole of ``being there'' resting on \textit{phantasia} behind a camera's movement through a flat frame.

The first combat arrives at 07:08 with a recognition and an eruption — ``It's goddamn Dutch Van der Linde!'' / ``Shit! Look out! There's more of the bastards!'' — and the O'Driscoll ambush routes embodiment through \textit{phantasia} rather than through a headset. The involuntary reflex at the on-screen muzzle flash, the body pulling back from a threat on a flat surface, is the flat screen's defining register of embodiment, the felt body of danger substituting for the spatial surround the headset would supply. It fires here for the first time, and the chain it belongs to can be named stage by stage, because the game will run it hundreds of times and this is its first full pass. The world offers its residue and perception takes it up (A₁): a figure, a rifle, a shout. \textit{Phantasia} condenses the uptake into the armed O'Driscoll as a \textit{phantasma} charged with threat (A₂), and the charge carries the discrete event's own \textit{orekton} — survive, and protect the gang the cutscene has just made worth protecting. The committal \textit{doxa} (\gk{δόξα}) — \textit{this is an enemy, to be met with force} — forms and ratifies the image (A₃), and at the A₃→A₄ transition, the event's center, opinion completes itself in movement: the trigger is pulled and the chain closes in the shot (A₄). Run against the \textit{Rhetorica}'s causes of action (I.10, 1368b32--1369a7), the deed is neither chance nor compulsion nor habit — there is no habit yet — but reasoning toward the installed good, with spirit rising under it as the ambush presses. The general causes align the same way: the designed ambush is the efficient occasion, the defense of the gang the end, and the player, newly armed, the agent. So the \textit{technē} of combat seeds itself not by instruction but by an attack the player must survive — which is already the section's lesson in miniature, that the skill and the disposition are installed in one motion: the same shot that teaches aiming teaches whom one aims at, and why.

Inventory arrives under looting's cover. ``Turn the place upside down. Grab as many supplies as you can'' (clip 3, 09:00), and the satchel UI thins the act of rifling a dead man's home — the texture of wood, the smell of blood, the weight of a tin — to a menu of icons and a text prompt, \textit{phantasia} left to supply whatever of the scene the player actually feels.
Even the act's tempo is the design's rather than the player's: as Vanderhoef and Payne observe of the game at large, ``Nearly every in-game action is linked to canned animations that leave the player waiting for the game to catch up to their intentions'' (Vanderhoef and Payne, ``Action and Stasis'') — intention waits on animation, the world metering the pace at which the hand may work. Inside the looting a single line catches and will not be cleared away: ``Poor bastard was married too'' (clip 3, 13:52).

What follows is where the dual claim becomes most precise. At 14:52 an O'Driscoll lunges from the dark of the barn and control is momentarily stripped — the corridor's first \textit{bia} suffered, the agent made patient until counter-input recovers the body — and within minutes Dutch hands Arthur lethal discretion over another O'Driscoll: ``Do what you want with him. I don't care. But bring that horse when you're done'' (clip 3, 16:00). The prisoner begs, and Arthur releases him: ``Get the hell outta here! Go!'' This is the first \textit{prohairesis} (\gk{προαίρεσις}) the player makes in the game, a freely chosen mercy toward a man on his knees asking to live. The honor machine does not score it. No meter moves in either direction; the fork belongs to the mission structure the honor system leaves untouched. But unscored is not unregistered — the world keeps its books by other means, as the beats that follow will show — and the dual installation stands: the \textit{technē} \textit{hexis} and the ethical \textit{hexis} arrive in the same clip, the same half hour, the same breathing sequence of mechanics. The player who learns to shoot at 07:08 is the player who learns to spare at 16:00, and the instruction is dual from its outset.

Sadie Adler's rescue follows at 18:00, where Micah manhandles a hidden woman — ``Get away from me!'' — and Arthur and Dutch intervene: ``Leave her alone. She ain't one of them'' (clip 3, 18:00). Micah's \textit{bia} against Sadie is met by Arthur's agency-restoring force, and the gang's self-understanding arrives in a single line — ``We're bad men. We ain't them'' — so that the ethical \textit{hexis} of in-group distinction, defined not against a positive ideal but against the O'Driscolls, is planted in the space between that line and the deed it names. Horse handling, tracking, Eagle Eye, and the weapon wheel arrive between 17:00 and 29:00 under the same narrative cover, each mechanic something the story needs. Eagle Eye, toggled at 24:00, is the first explicit case of the design lending a perception it cannot cultivate: a vision mode supplying in highlighted tracks and scat what a trained tracker's gaze would otherwise hold as bodily knowledge, the filter standing in for the eye the screen cannot train. At the canyon traversal (29:00–31:00) the companion patter tightens the leash once more — ``It's slippery. Be careful. / Stay close to the wall. / Get up here'' — the sustained scripting of every footfall. Then, at 33:00, the wolves: A₂→A₃→A₄ on repeat under time pressure, the lunging animal as \textit{phantasma}, the snap \textit{doxa} (\textit{threat, kill it}), the shot — and the flinch firing again, sharper now, because the wolves lunge and the screen can place the body nowhere but here. ``Good work, Arthur'' (clip 3, 38:10): the act attributed to the agent's competence, the \textit{technē} of combat under duress named and affirmed.

Against this whole sequence the choke/spare fork stands as the tutorial's plainest disclosure of how its moral bookkeeping actually works. In the interleaved beat (capture 3a) Dutch hands Arthur the same lethal discretion, and again the honor meter is silent: a choke and a spare are mechanically identical, and — like the prisoner release before it — the fork belongs to the mission structure the honor system leaves unscored. But the world does not leave it unmarked. If the player chokes the man, a bell tolls over the act (clip 3a, 11:31–11:45): no attribution, no delta, only a sound that carries the moral weight aesthetically and leaves the sentencing to the player's own \textit{doxa}. The pattern will hold across the tutorial's mandatory events — the train hostages of the coming heist are likewise unscored, their fate reported instead by the local newspaper, survivors or none\footnote{****** Insert citation: playthrough video demonstrating the newspaper's report of the hostage outcome.} — and it is worth stating as a rule from the outset: in the mandatory, narrative events the honor mechanism stays neutral, and registration passes to other channels — a toll, a printed line, a companion's changed regard. The score is one register of a world that is watching in several.

\bigskip

\subsection*{The Hunt: A Desire Made Legible}

Against the corridor's density and pace the hunting tutorial is something rarer — a chain that Charles narrates aloud, stage by stage, so that the actualization of desire can be watched as it happens. The clip is useful not because the chain requires narration to run, for it runs whether or not anyone speaks it, but because Charles verbalizes each stage and so makes the apparatus's operation, everywhere else smuggled under story, unusually legible. That legibility is the section's reason for being.

The chain begins before any animal appears, at the horizon. ``A few cans of food and a rabbit… for 10, 12 people?'' (clip 4, 00:58) — the scarcity is designed, and the design is plain: the camp cannot be fed on what it has, and so someone must go. The \textit{orekton} — food, survival, the gang provisioned against the storm — is installed before the hunt begins, so that the kill will read, from the moment it is made possible, as provision and not as the agent's bloodlust. The motive is pre-assigned, the act licensed in advance, coded as a deed of scene and purpose the honor machine will not need to stir for. Charles hands the bow at 02:24 — ``You take this. I can't use it and you'll have to'' — and the weapon specifies the \textit{technē}: a silent kill, the bow chosen because the gun would scatter the prey. They climb the mountain, Charles reads the terrain aloud, and the scarcity-horizon follows them up into the trees.

At 06:10 the trained perception the analysis calls A₁, the uptake of the world's offered residues, is modeled by the tutor before the player performs it: ``There's deer been here. Recently. How can you tell? / How can you not?'' The Eagle Eye vision mode is toggled, the screen's track highlights supplying what a trained hunter's gaze would otherwise hold as body knowledge without a UI — the design lending outright, in a toggled filter, the perception it cannot cultivate in a single session: the player presses a button and the game paints the ground. ``It's easier in the snow, but once you get your eye in, you'll be able to track them nearly as well in grass and woods'' (clip 4, 06:28): the \textit{hexis} is named as portable, a disposition that will carry from terrain to terrain, the game framing aloud the very thing it is cultivating.

``Good tension on the string before releasing. Just don't overdo it. Now, Arthur'' (clip 4, 07:52). The bow draws, the tension indicator runs, the reticle steadies, the arrow flies. In four seconds the full chain closes, and every element the method names is on the surface. The \textit{orekton} of the discrete event was installed at the horizon an hour before the shot — food, the camp provisioned — so the deer, taken up through the modeled tracking of A₁, is condensed by \textit{phantasia} into a \textit{phantasma} charged with exactly that desire (A₂): not prey in general but provision, the kill already coded as care for the gang. The \textit{doxa} ratifies — \textit{this is the shot, take it} (A₃) — and at the A₃→A₄ transition, the event's center, opinion completes itself in movement: the string is loosed and the chain closes in the kill (A₄). Among the \textit{Rhetorica}'s causes of action the deed runs on reasoning through and through — nothing compels it, no appetite of the agent's own drives it, and the good it serves was handed down by the world an hour earlier. The felt body of the drawn bow — the tension across the back, the steadied shoulder, the held breath — is supplied entirely by \textit{phantasia} behind a button hold and an analog stick. The mechanic abstracts the act, \textit{phantasia} gives it flesh, and four seconds of screen time close what three minutes of narrated preparation had opened. ``Let's try for another'' (clip 4, 08:04): the loop is set to repeat, because repetition is how a \textit{hexis} sediments.

``Easier when they ain't shooting back'' (clip 4, 12:35) is a throwaway line that names the clip's own register. Nothing here charges the player: the deer does not lunge the way the corridor's wolves lunged at 33:00, and the involuntary recoil that fired in that ambush never fires on this mountainside. The embodiment is quiet and accumulated rather than sudden and involuntary — a disposition built by patient repetition where the wolf ambush made its deposit in a single flinch. The carcasses go to Pearson at 19:19, the kill dissolving into the camp economy and the A₄ act depositing its residue into the shared A₀; the \textit{orekton} was always the camp's, merely routed through Arthur's hand. Then the tutor names the \textit{hexis} outright: ``Takes a lifetime of practice to master'' (clip 4, 19:09) — not a skill set but a disposition built by time, not what one can do after thirty minutes with a bow but what one becomes after a lifetime of using one. The bow clip is the moment the tutorial says aloud what, everywhere else, it only does.

\bigskip

\subsection*{The Heist: An Acting-Together}

Against the corridor and the solitary hunt, the Cornwall train robbery is the tutorial's collective event — the first time the full actualization of desire draws not on the player's solitary hand but on the gang's distributed agency, so that what is achieved cannot be credited to Arthur alone. Here Burke's consubstantiation becomes literal — ``a way of life is an acting-together'' (Burke, \textit{Rhetoric}, 21), and the gang's way of life is enacted in the shared act — and here the analysis meets co-disclosedness (\textit{Miterschlossenheit}) in its plainest event-level form: the world of the job is opened to the player \textit{with} others from the outset, its dangers and its spoils disclosed as shared, the ``we'' itself the source of the deed. Dutch opens the sequence by sublimating Colm O'Driscoll's personal revenge into collective provision — ``much less fun to rob him if he never finds out about it'' (clip 5, 06:21) — Burke's act–purpose ratio run deliberately, the act reframed from attitude (revenge) to purpose (plunder for the camp). The approach, the cover positioning, the firefight through the lake camp, and the leap to the moving train (clip 5, 42:42) — ``I'm slipping! Pull me up!'' — are each a phase of a planned collective act in which the gang moves together and Arthur is one part of the moving whole, the vertigo of the teetering car supplied wholly by \textit{phantasia} behind a stick and a button that perform the grab. The flinch fires anyway.

At the private car (clip 5, 51:00) Dutch hands Arthur a hostage choice — ``Kill them, leave them here, take them with you — just make sure they don't send no folk after us'' — and the playthrough takes the non-lethal path, marching the workers down the train and keeping them silent without killing. \textit{Prohairesis} again: the deliberated act that finds the option the ethical \textit{hexis} can live with. The honor meter, true to the rule the corridor established, does not move; the fork is mission structure, and its registration arrives instead in print, the local newspaper reporting the robbery with survivors — or, on the other path, with none. ``Proud of you boys. Not a man down. Outlaws for life'' (clip 5, 25:29): the ethical \textit{hexis} of in-group fidelity is named as a creed, the gang's shared substance offered up as an identity that carries across acts. Horseshoe Overlook arrives at 65:00 with the new horizon carried in four words — ``This place is perfect'' — and the \textit{orekton} shifts with the home base, from survival to prosperity, from Dutch's grief in the blizzard to Dutch's sparkling eyes over a valley he means to hold. ``The camp gets its slice'' (clip 5, 65:40) names the new governing desire — not imposed but solicited, the world beginning to hold out the objects it will invite the player to adopt as their own: provision for the camp, a share in its prosperity, standing among its people. The residue of the heist deposits into that new ground, and the tutorial passes into its second phase.

\bigskip

\subsection*{The Camp: Dwelling Before Desire}

The Horseshoe Overlook clips ask a different question of the method than any combat beat: what deposits a disposition when nothing is at stake? The player grinds no quest here. They chop wood, carry sacks, drink coffee, and are greeted by name — and it is in this unpressured stretch that the \textit{orekton} the heist installed, belonging rather than survival, is worked into a standing \textit{hexis}. What the section establishes is that the deposit is made by repetition, and that its instrument is the re-learning of the world's objects: what a thing in this camp is for turns out to be discoverable only by living with it.

``Morning, Mr. Morgan'' (clip 6, 00:00). The greeting registers not what Arthur does but that the camp \textit{knows} him, knows him by name and acknowledges his standing at the day's first light. It is recognition of standing, not credit for a deed — and standing is exactly what the ethical \textit{hexis} of belonging is, not the act but the deposit of prior acts, the accumulated residue a gang hails when it says the agent's name at dawn. Consider, then, what happens with the coffee. The first time the player notices that coffee can be made and drinks it, the act may be driven by novelty, by curiosity, by the pull of the camp's realism — a sampling of what the world offers, nothing more. But the cup delivers a buff to stamina, and the discovery changes what coffee \textit{is}. In Uexküll's terms its functional tone shifts: the pot at the fire is no longer set dressing to be sampled but a resource, understood from then on through its newly sedimented function, reached for before a long ride the way a tool is reached for. The same arc runs through the camp stew, first tasted out of curiosity and thereafter eaten for what it restores. It runs, most instructively, through the chores. The first log is chopped for the novelty of the animation, or out of the felt realism of camp life; then the player discovers that the labor is scored — the honor system credits it positively, and the gang answers it, members acknowledging the work and warming toward Arthur — and from that discovery forward the woodpile carries a different tone: labor as membership, work done for the camp that the camp visibly returns. Each discovered function recharges the \textit{phantasma} under which the object appears, so that the camp the player perceives after days of dwelling is not the camp they first rode into. The world has taught the player what its things are for, and the teaching is a re-attunement of perception itself.
This is the register Westerside and Holopainen built the term \textit{gameplace} for: on their reading ``The camp is the place of care, maintenance, and familial social interaction. Thus the camp, even though not staying in the same location throughout the game, acts as a home for the player'' (Westerside and Holopainen 9), and the claim beneath it — that ``the gameworld is not just a space for traversal, but a world for being in'' (Westerside and Holopainen 11) — is the phenomenological premise this subsection has been working out. The camp is not a hub with services but a place in the lived sense — and a designed one, its solicitations aimed: the requests for food from gang members, the contribution box, the chore prompts, each an invitation to make the camp's needs the player's own objects of desire.

It is here that Calleja's apparatus earns its place most plainly, for the channel profile of camp dwelling is unlike any combat beat the tutorial has run. The \textit{spatial}, \textit{kinesthetic}, and \textit{affective} channels are dominant — the player moving through a bounded place, doing the things the place asks, and being seen to do them (Calleja, \textit{In-Game}, 36–38) — while the \textit{ludic} channel, the engagement with consequential choice, falls to near silence, no goal pressing and nothing won or lost. This does not mean desire is absent. An \textit{orekton} is always present — a soul without one does not act at all — and the camp's quiet is not desire's absence but its satiation, the governing object (to stay, to be of use, to belong) so nearly at hand that no gap opens for a quest to fill. Nor is the clip itself what cultivates; the game does the cultivating, and it does so through both of its registers of recognition at once. The mechanical register scores: the chores credit honor. The social register answers: gang members and NPCs respond to what Arthur does and has done, their greetings and asides tracking his conduct, so that the world is disclosed \textit{with} others whose regard is part of what is disclosed (\textit{Miterschlossenheit}). ``Good work, Mr. Morgan,'' should the greeting fire mid-chore, is that register speaking — character recognized not from a single act but from the pattern of showing up. What accumulates across the long ambient loops is the \textit{hexis} of membership, deposited quietly and by repetition, since one chore deposits nothing and a hundred chores deposit everything. The blizzard-survival horizon has closed; the standing-and-belonging horizon is now the ground.

Heidegger's second form of boredom bears directly on this stretch, and on that dominant channel-profile — not the boredom of waiting for something but the being-bored-\textit{with} a situation one nonetheless remains in, the being-held in an empty time that yet discloses the world as a whole (Heidegger, \textit{Fundamental Concepts}, 120–124). The looping ``Morning, Mr. Morgan,'' the chore-prompt, the woodpile, the fire: this is exactly that texture, the camp as a world one is bored-with and yet does not leave, because remaining is itself the disposition the tutorial is building. The second form of boredom is at once the analysis's object and its medium — the phenomenal register in which the camp's repetitions do their work, the \textit{ludic} channel near silent, the \textit{spatial} and \textit{affective} channels carrying the whole. What that work yields, however, is not dwelling handed over as a deposit. Being-at-home cannot be laid down in a person the way a skill can; it can only be encouraged, and it settles — when it settles — as a mutual determination of person and world, the camp offering its places and rhythms, the player taking them up until the taking up needs no decision. What the design deposits is everything short of that: the familiarity, the competence, the pattern of showing up. The dwelling itself is the player's answer.
The structural register would dissolve all of this into function: for Tuominen, ``the character development and different side quests, activities and tasks can be considered to exist in these games solely to interpellate a player to progress in the story'' (Tuominen 13). The account is not wrong as design teleology — the chores do keep the player progressing — but it cannot say why the loop \textit{feels} like home rather than like work, and that felt difference is what the disposition account explains: interpellation names what the design wants from the player; the disposition account names what the design cultivates in the player, and why the cultivation can feel like home.

\bigskip

\subsection*{The Town: Honor as Motive-Attribution Machine}

If the camp is where the ethical \textit{hexis} is cultivated quietly, the Valentine trip is where the calibration of the honor machine is put on open display, and where its design-assumptions are brought to the surface. Within thirty minutes the player is handed five consecutive honor-positive events, and in the spaces between them Arthur twice disclaims the very virtue the machine has just credited him — and these two movements, the scoring and the disclaiming, are what the section is built to set against each other.

No mechanic in the game has drawn more scholarly attention than this one, and the readings divide three ways. For McEvoy the machine is the game's persuasive center — ``The honor system in RDR2 mitigates how players can go about experiencing the game and building their experiences'' (McEvoy 32) — and its input logic is plain: ``The system is simple -- be polite to NPCs, choose to spare lives when possible, and perform as little crime as possible, and the game will understand that you are a high-honor individual'' (McEvoy 33). For Tuominen the same loop is Althusserian, an apparatus that hails: ``by giving the most lucrative rewards on completion of its story missions, it calls or hails the player to focus on them'' (Tuominen 45). For Bagnoli it is arithmetic — ``by identifying dialogue actions as positive and violent actions as negative, the game creates a polarized system of values'' (Bagnoli 42) — a morality of one axis, inscribed, he argues after Sicart, in Arthur's character rather than in the player's conscience (48--49). All three readings take the machine's coverage as given; what it declines to see has drawn less notice, though Bagnoli states it flatly: ``it must be said that `honor' is not present in the main missions'' (Bagnoli 42). The playthrough bears him out more fully than his own examples do. The meter is silent not only over the missions' obligatory violence, which goes unscored even against lawmen (43), but over their most charged choices — the Colter prisoner's release, the choke/spare fork, the train hostages — none of them scored, the weightiest of them registered instead through the world's other channels, a tolling bell, a newspaper line. That silence is not a defect. It is design intruding on its own mechanic for the experience's sake, and each face of the intrusion discloses a different design commitment. Were the compulsory killing debited, no high-honor Arthur could exist at all: the meter's neutrality inside the missions protects the story's playability from the story's own body count. Were the mission forks scored, the machine would stand between the player and the choice, converting deliberation into arithmetic; left unscored, the forks are undergone morally rather than transactionally, their weight landing in the player's own \textit{doxa} rather than in a delta. Where the machine does run — in the open world, over greetings and thefts and shot dogs — its pricing is asymmetric, as McEvoy's own observation shows: ``reactions and law enforcement act as a fence or barrier to encourage players to act morally'' (McEvoy 12--13). The fence prices conduct rather than forbidding it; a low-honor run remains possible but substantially harder than an honorable one. In Calleja's terms this is the \textit{ludic} channel worked as rhetoric — the reward structure an underlying, designed encouragement toward a specific playstyle — and it is the machine's scores and silences together, watched act by act, that expose exactly which motives it was built to track. What the Valentine clip adds is the finest-grained view of that calibration the game affords.

The first event opens the sequence: a stranger, feigning lumbago, needs his horse retrieved, and Arthur retrieves it and returns it, the meter bumping. ``You're a gentleman, sir.'' Arthur answers: ``No, not really. I was just trying to impress the women'' (clip 7, 07:45). This is the Burkean scene–agent ratio made explicit by the agent himself, the honorable act credited by the machine to his character and re-credited by the character to an audience — the women the scene-condition without which he would have robbed the man. The machine scores what it sees; Arthur reports what moved him; and neither attribution is false, nor is either complete. Then, a beat later: ``If you lot hadn't been here, I probably would have robbed them'' (clip 7, 08:40). Uncle supplies the scene-attribution in the very form the honor machine elides, for the machine cannot read the subjunctive — the robbery that did not happen because the witnesses were watching — and so scores only the act that did, and scores it to the agent. The dispositions on display are not identical to the scores that track them: this is the clip's first payload, that honor-attribution and moral truth diverge, and that the character names the divergence for the player while the machine files its report.

Tilly's grab (clip 7, 23:00), Karen's beating (clip 7, 24:16), and Jimmy's dangling plea (clip 7, 28:38) are three \textit{bia} events, three de-inversions, three honor-gains, the structure rhythmic — patient to rescuer, each act coded to the agent's standing. But the acts themselves are mixed: intervening for Tilly is protective and uncalculated, a reflex of the scene. Beating the john in the hotel room is violent in excess of the need, Arthur striking him ``a lot harder'' than he hit Karen, as Karen reports afterward (clip 7, 26:50). The honor machine scores both alike, because it cannot read degree. Burke's pentad has a slot for act and a slot for agent, and what falls into each is fixed by the geometry of the scene and not by the player's conscience.

At 29:00 the machine briefly announces itself: an instruction appears in the top-left of the frame explaining that positive decisions will raise Arthur's honor — the system stating its own input logic mid-sequence, inviting the very playstyle its pricing favors. Jimmy Brooks is the clip's keystone. The player has chased him down, lassoed him from horseback, and hauled him to a ledge — all \textit{bia} inflicted, the player the agent of force against a begging man — and then at 30:00 Arthur hauls him up rather than dropping him, and the meter rewards it. The speech that follows is the tutorial's most complete pentad in a single breath: ``I'm not a good man, Jimmy Brooks… I was in Blackwater. I kill people. And maybe I should have killed you… I never saw you. We have an understanding'' (clip 7, 30:00) — scene (Blackwater, the fugitive life), agent (Arthur naming himself a killer), act (spare and intimidate), purpose (the gang's safety), attitude (mercy edged with menace). The honor machine scores the spare as positive, and the player's conscience must supply the weight of the coercion that surrounds it. One terminal is moral credit, the other an asymmetric power relation enforced by the threat of death, and the machine scores the one and is silent on the other. At 32:10 Arthur returns the commandeered horse — ``You really were just borrowing it. Appreciate it'' — and the loop closes. The ethical \textit{hexis} of these thirty minutes is not clean virtue but its approximation under particular scene-conditions, witnesses and women and a fugitive who can name the gang. Strip those conditions away, and Arthur is his own narrator's confession, ``I probably would have robbed them.'' What the machine deposits is a score. What the tutorial cultivates is what the score can only approximate — the disposition to act honorably when the scene calls for it, tracked by a machine that cannot read the gap between the disposition and its measure.

\bigskip

\subsection*{The Disposition Built}

Three hours of tutorial; three raisings of the governing desire, from blizzard survival to standing and belonging to the outlaw economy's shared prosperity; two dispositions, the craftsman's \textit{technē} and the ethical agent's settled orientation, deposited together across combat and camp chores, heist and solitary hunt. What the tutorial has built is neither a skill set nor a personality but a standing readiness to act in certain ways in certain kinds of situation, cultivated by repetition until it runs without instruction. The ethical \textit{hexis} runs across the camp and the town without any explicit lesson in virtue because Dutch's grief speech seeded the orienting good and each subsequent resolution has deepened its furrow. Neither disposition was set by explicit teaching; the tutorial designed a world, and the world cultivated the disposition.

The proof that the \textit{technē} has settled arrives in the one clip this chapter has left unnarrated, because nothing in it is narrated: the legendary bear hunt. Where the guided hunt was legible because Charles spoke every stage aloud, this clip is diagnostic because it is silent — no tutor speaks, no stage is named, and the \textit{hexis} seeded under instruction now runs on its own. The horse is whistled for and comes; the trail is read through Eagle Eye; the bear charges (clip 9, 11:00) and the image lunging within its flat frame still pulls the player's body back, \textit{phantasia} supplying enough of the predator to move a spine through a screen. Even the care that precedes the hunt feeds the performance: the grooming, feeding, and riding that bond horse to rider (clip 9, 05:00–07:30) measurably improve what the animal can do — new gaits and maneuvers unlock, stats rise, the spook recovers faster — unscored acts of care depositing directly into scored capability. Where a body trained over years would hold the tracker's eye, the hunter's felt strangeness near the quarry, and the marksman's composure as its own dispositions, the flat screen lends each as an instrument — Eagle Eye's highlight, the legendary aura's tint, Dead Eye's slowed frames — three eidetic and affective states converted to filters and on open display in a single encounter.
(McEvoy reads Dead Eye as the game's emblem of earned skill — ``If a player masters even this smaller system within the larger game, they can heighten their experience'' (McEvoy 24); Vanderhoef and Payne read this very hunt through tempo: ``While designed as a tutorial, this story-based introduction to legendary animal hunting illustrates the game's systematic deployment of slow game time during optional activities that repeatedly thwart the goals of hurried, determined players'' (Vanderhoef and Payne, ``Action and Stasis'') — the slowness the very medium in which the settled \textit{hexis} can show itself.) The clip ends without commentary: no ``Good work, Arthur,'' no ``takes a lifetime of practice to master.'' The \textit{technē} at ceiling is its own testimony.

This is what the RODA analysis discloses when it is run across the whole tutorial arc: not a series of discrete events but a single compound disposition deposited incrementally, each clip adding a layer the next draws upon, the whole intelligible only from its end looking back. The opening cutscene names an Arthur before the player can act; the legendary hunt shows what that Arthur has become once the player has acted for three hours. Between them lies the tutorial, a guided \textit{paideia}, building the disposition that will run the rest of the game. The two strands are by now one deposit — the \textit{technē} that draws a bow in silence and the ethical \textit{hexis} that spares a begging man — and together they are the prepared agent the open world will inherit.

It was built, moreover, by the route the flat screen makes necessary, and that route is the chapter's contribution to the differential the two cases compose. Where the head-mounted world of \textit{Gnomes \& Goblins} hands the body its surround and asks little of the image, the screen withholds the surround and asks everything — the cold of the blizzard, the heft of the drawn bow, the lurch of the leaping train, the recoil from the charging bear, each supplied by the player's own \textit{phantasia}. Where the eye alone will not suffice, the design lends instruments: Eagle Eye for the tracker's gaze, Dead Eye for the marksman's composure. The disposition deposited is the same kind of thing the virtual-reality world deposited in its player, an incorporation that sediments into a \textit{hexis}; what differs is the road. The screen-based route runs through the image and through what the player makes of it — the design does its work by shaping the player's cognitive orientations, seeding the \textit{phantasmata} under which its world appears, and steering the \textit{doxai} the player forms about that world, reaching the disposition through opinion where the headset reaches it through the body.

One further deposit belongs in the bundle's inventory, and it is a capacity rather than a tendency: the trained readiness to experiment. Moser, importing Stanislavski's magic ``if,'' reads free play as a rehearsal space: ``Playing with scenarios allows players to see not only how Arthur reacts to situations but also how they react when they are confronted with opposition or ethical actions through him'' (Moser 24). The save slot makes the rehearsal literal — a fork can be played, judged, and loaded away, an ethical experiment run under conditions objective reality never grants — and the honor machine gives the experiments a public ledger, its far ends two different endings for the same man. McEvoy's summary of that economy — ``The basic commentary the endings provide is that goodness is rewarded with peace while misdeeds are punished with suffering'' (McEvoy 33) — names what the design counts on: that a player free to try both will come to know the difference from the inside, and will replay to see what the other disposition yields.

The guided phase will end. The world will open, the camp economy will expand, Valentine will become a hub, and the seeds planted in the heist and the town will begin to bear. Everything the tutorial deposited will be drawn upon across the potentially hundreds of hours that follow — and tested, when the gang the loyalty \textit{hexis} was built around begins to come apart, against Micah's corrosive creed and Dutch's slow erosion and the moment in the game's last act when the disposition is all that remains.
Nor will the push-narrative — Calleja's name for the scripted story that advances on the player regardless of the player's own designs — disappear when the guided phase does: the design keeps its right of seizure, and the scholarship has mapped its later uses. Moser, reading the forced beating of the sick farmer Thomas Downes — the debt-collection the player can neither refuse nor soften (17) — concedes the constraint as load-bearing: ``These narrative points must be rigid in their structure for the rest of the story plots to come together, regardless of how players decide to play their Arthur'' (Moser 17). Ruffino follows the same beat downstream into the tuberculosis the player community traces to it (351--352), the arc in which Arthur's hold on his goals, his friends, his resources, and finally his own lungs is withdrawn (353), its cause never shown and its cure nonexistent. These are seizures whose consequences land on Arthur — in stats, in story, in death — and reach the player only through him. The event this study turns to next inverts that landing: a seizure priced at one dollar in Arthur's ledger and paid almost entirely in the player. Those tests belong to the disruption, which the following section reads frame by frame. What the tutorial built was the world that is made to break.

\textit{(One tutorial beat noted above still awaits frame verification before the final pass: the exact timing of the bell toll in the choke/spare beat. The §2 disruption footage has since been captured and frame-analyzed; the section follows.)}

\bigskip

\section*{6. \textit{Red Dead Redemption 2} (§2: The Disruption): The Sonny Event}

Turning to my primary artifact, Rockstar's \textit{RDR2}, I believe one specific event can demonstrate how a virtual world can do violence at the level of a player's being-in-the-world, and how that violence can transform it. I have selected this game because it features one of the most robust, immersive, and realistic virtual ecologies I have ever encountered, and I have selected this event because it is the place where that ecology turns on its inhabitant. Here I would like to pause to mention that this event contains implicit sexual assault and, though nothing is definitively shown, it can still be quite disturbing. Arguably the most controversial event in \textit{Red Dead Redemption 2} is a scene wherein the game's protagonist, Arthur Morgan, is implicitly raped by a man named Sonny who lives near Lakay in the swamp section of the map. The analysis runs the event through the actualization of desire: the chain by which a world is taken in through perception, condensed by \textit{phantasia} into a charged image, ratified by \textit{doxa} as an opinion, and completed by \textit{orexis} in action. As we will see, the event's design demonstrates a command of this chain that goes well beyond anything Calleja's involvement model alone can articulate. The chain itself is not something a design can seize: it is simply how a human being is in its everyday being-in-the-world, and it runs here as it runs everywhere. What the design does is arrange the world so that the chain's ordinary functioning delivers the player, by the player's own completed act, into the one position from which agency can be taken --- and then it takes it.

Consider first where the event lives and who arrives there. Lakay sits near Saint Denis, a city the player, if they are following the main narrative arc, will not reach until the second half of the game, after roughly 25 hours of gameplay if the player solely follows the main questline; given this location, it is likely that the player will have spent enough time playing to have mastered, or internalized, the controls for movement (\textit{kinesthetic involvement}). Having mastered this dimension, the player's conscious attention will be directed to other dimensions of the model, and, as Calleja notes, once a player has mastered the controls it becomes possible to cognitively interpret the actions of the avatar as \textit{their own} (68). Only when the controls and hardware have become what Heidegger calls ``phenomenologically transparent''---a state wherein we do not need to direct conscious attention to them---will the player feel present in the virtual world, as if they exist within the ecological network of mutual affectability the game establishes. But the deeper preparation is dispositional, and it is worth naming precisely. By the second half of the game, hundreds of designed encounters have sedimented a settled \textit{hexis} of seeking: the unmarked cabin holds the treasure map, the unknown stranger opens the memorable quest, the question mark on the horizon is always worth riding toward. The game has spent dozens of hours teaching its player that curiosity is a disposition this world rewards. Their mood (\textit{affective involvement}) may be one of curiosity, leading them to explore this section of the map; it may be one of determination, if they are hunting the hidden items and easter eggs strewn throughout the world (\textit{spatial involvement}). Either way, the player who reaches Sonny's stretch of swamp arrives as a being whose very way of taking in a world---whose \textit{how} of being in this world---has been cultivated by the world itself and the structural, ludic mechanisms that govern it. That cultivation is the trap's true mechanism.
The scholarship describes this arrival from two sides. Moser's account of character attachment makes the cultivation temporal — the bond with the avatar is ``a gradual establishment over time as players get closer to Arthur through their playing'' (Moser 33), until Arthur has become ``less of a computer-generated being and more of a full-fledged person'' (Moser 13) — and her lifeworld frame carries the deeper point: what play deposits does not sort itself by ontology, the dispositions brought into the game and the dispositions built inside it belonging to one continuous life. Westerside and Holopainen supply the complement from the world's side: ``It is a world that appears to exist in spite of us, rather than for us'' (Westerside and Holopainen 10) — and that felt autonomy is precisely what lets the place turn, for a world experienced as existing in spite of the player is a world whose invitations read as discoveries rather than as designs. The two sides are, in the vocabulary this project has carried throughout, world-disclosedness (\textit{Welterschlossenheit}) and co-disclosedness (\textit{Miterschlossenheit}) at full strength: a world opened as one the player skillfully inhabits, and opened \textit{with} others — Arthur nearest of all — whose regard and fate are part of what is opened. The event exploits both at once: the opened world supplies the trust, and the co-disclosed bond supplies the channel along which the harm will travel. My recording opens on exactly this mode of existence: several minutes of unhurried riding through Bayou Nwa's fog, a glance at the map, a passing stranger's cry (``Please... help!'') noted and left behind---a player dwelling in a mastered world, open to what it offers.

% ****** OPEN (clip): insert video clip of the Sonny approach

Consider how much of the work the player does. Nothing about this event is sought---the player rides through the swamp on business of their own---so the world must reach out first. Upon approaching the home the player is hailed by Sonny before the cabin has even fully resolved out of the fog: ``Hello again... don't be shy now, come over to see old Sonny, come on.'' The ``again'' is not a script's slip: the system recognizes a prior visit, and a player who has approached before without entering meets an escalating repertoire of coaxing lines written for exactly that case.\footnote{****** Appendix pending: the full set of dialogue variants Sonny uses across repeated visits when the invitation is declined.} The hail is a designed pull on the chain's opening link --- the world offering, at A₁, precisely the residue it wants taken up, an environmental solicitation addressed to a rider who sought nothing --- and the dialogue that follows is scene painting of a very particular kind. ``Ain't this a fine place... an interesting place, the best of places, huh? Is it land or is it water? Can't make up its mind... I can't make up my mind about things neither.''
Sonny builds himself out of his own landscape: the swamp that is neither land nor water, the man who cannot make up his mind, one substance under two descriptions. The environment's uncanny ambiguity---the very quality that ambiently attunes the approaching player toward unease---is transferred onto the agent who inhabits it, so that the felt strangeness of the place and the felt strangeness of the man become one attunement (\textit{shared involvement} riding on \textit{affective involvement}). Then the lure lands as a determinate image: ``You wanna come in, see my collection of skulls?'' Sonny's repertoire is not single --- on other approaches he baits with food, the invitation I myself first met years ago --- but the skulls are the sharper instrument, and the difference between the two baits makes the design legible in miniature. Food is a need the game's economy met long ago; by this stage of the game nothing about a meal can charge a \textit{phantasma}. A collection of skulls is another matter. It is a curiosity object, macabre and singular, and it is aimed past need entirely --- at the player's cognitive orientations, at the \textit{doxa} the whole game has ratified a hundred times over (\textit{the strange find is worth the detour}), and at the emotion that attends that opinion, curiosity with unease running under it. The invitation pits the player's trained way of being-in-the-world against the plain sense of the words---a lonely man, a swamp, skulls---and, for most players, the training wins. Sonny's doorway is that lesson's trap. The image the invitation plants is the \textit{phantasma} (A₂) on which everything that follows turns: the promised sight, charged with the seeking that hundreds of rewarded detours have made second nature.
McEvoy, mapping Salen and Zimmerman's definition of play onto this game, locates its persuasion exactly here: the open world is the ``free space of movement within a more rigid structure'' (Salen and Zimmerman, qtd. in McEvoy 12), and the structure is the coded part — ``how the characters are coded, the social systems in place, such as law enforcement, and how NPCs will react to the negative or positive things Arthur does'' (McEvoy 12). Sonny's hail is that rigid structure reaching into the free space, a scripted solicitation dressed as an open-world find; in Calleja's terms, the \textit{ludic} channel is worked from the world's side before the player has registered that there is a choice to make.

What the recording shows next is the part of the event that analysis usually cannot see: the formation of the opinion itself. For roughly nineteen seconds the player character stands motionless before the porch, the interaction prompts still live, the doorway a pane of solid black. Nothing forces the next step. The footage even catches the player opening the weapon wheel during the approach---a visible trace of suspicion, the wary opinion weighing against the curious one---and the game's four-option interaction menu (\textit{Aim Weapon / Greet / Antagonize / Stranger}) keeps the armed alternative bright and available through the entire hesitation. This is \textit{ludic involvement} in its fullest sense: a choice under repercussion, made slowly, with the means of refusal on screen. It is also A₃ caught on camera --- the \textit{doxa} forming in real time, the sedimented wariness (a swamp hermit, skulls, unease) weighing visibly against the ratifying pull of trained curiosity.
``A player's choices are wrapped up in the choices of the character'' (Moser 21), Moser writes of this game's forks, and the entanglement is the design's whole interest in the pause: what is being weighed on the porch is not Arthur's risk but a joint stake. The design knows the player is deliberating, because it taunts the delay in Sonny's own voice: ``What you stood there for?'' The line does to the player's hesitation exactly what ``don't be shy'' did to the player's approach---it works on the forming \textit{doxa} from the world's side, reframing caution as a social failing, a discourtesy to a lonely old man, so that the suspicion tipping the scale is made to feel like the thing that needs justifying. The rhetoric is aimed at the ratification itself: no new evidence is offered, only pressure on the weighing. Note what the world is doing: it cannot supply the player's desire, so it manipulates the material situation---the hail, the banter, the lure, the taunt---to solicit the desire the player must supply. When the player finally walks in, the act is theirs, the choice is real, and this matters more than anything else in the sequence. Every element of a free, agent-caused action is present and on the record: an object of desire (the \textit{orekton}), a formed image (the \textit{phantasma}), a visibly hesitating opinion (the \textit{doxa}), and the committing movement---A₄, the step through the doorway, the chain completed by the agent and by no one else. The design requires exactly this, because what it is about to take is the player's agency, and agency can be taken only from someone who has just exercised it: a seizure that arrived mid-ride in an empty field would read as the medium's ordinary interruption, while a seizure built one step past a freely ratified choice converts the choice itself into the trap's final component.

Then, exactly one step past the threshold---at 02:31.67 of my recording, with the avatar's body already inside the doorway---the minimap and prompts vanish, the cinematic bars snap shut, and control is gone. If the player accepts the invitation, a cutscene is triggered that cannot be skipped: the push-narrative, the scripted story asserting its right of way over self-directed play, here pushed to a limit Calleja's category does not anticipate. The frame record sharpens the point that matters most: the seizure does not occur at the threshold but one step \textit{past} it, after the player's own completed act has delivered the avatar beyond the last point of return. The system intercepts the player's self-directed movement and reconfigures it as pure being-moved. This is \textit{bia}, compulsion, the action whose origin lies outside the agent (\textit{Rhetoric} 1368b33--1369a6), and the particular placement of the push-narrative is the driving rhetorical mechanism behind the event's affective force. It does not simply curtail agency in general; it \textit{answers} the player's choice with a sudden re-inscription of that choice into a scripted fate. A cutscene that merely interrupts play suspends agency; this one \textit{uses} it, converting the player's own desire---their curiosity, their cultivated seeking, their freely ratified opinion that the invitation was safe or manageable enough---into the efficient cause of the player's undoing. This is why the event reads as violation rather than as cinema. The threshold thereby becomes a rhetorical site where the world forcibly moves the player in a direction they likely did not intend nor desire, revealing that their affective exposure was always already in play.
The agency scholarship holds one precedent for seizure on this scale, and the comparison is exact enough to be instructive. Ruffino reads the game's tuberculosis arc as its great revocation — agency in games being, as he frames it through Owen, ``the affect of being in control'' (Owen, qtd. in Ruffino 351), and this the game in which ``Arthur loses control of his goals, friends, resources, even his own lungs'' (Ruffino 353). Ruffino also records the engine's small mechanical tell at that arc's origin: before triggering the non-interactive Downes fight, the game removes the mask Arthur wears, stripping a precaution the player had deliberately equipped (352). But the tuberculosis is ``the possibility that consequences could manifest with no preceding actions'' (351): the cause is never shown, and the weight falls on Arthur — his stamina, his story, his death — reaching the player only through him. The Sonny event runs the same seizure in reverse: here the action is fully shown and freely made, and it is the consequence's weight that changes address, almost nothing landing on Arthur — a dollar, three refillable cores — and almost everything landing in the player. The two seizures are inversions of one another, and the inversion states the event's principle: this compulsion was aimed past the avatar from the start. In Calleja's terms, the player's attention is redirected to the \textit{narrative}, \textit{shared}, and \textit{affective} dimensions while the \textit{spatial} and \textit{kinesthetic} dimensions fade once the player realizes they have lost control of the avatar; but ``fade'' is too gentle. The \textit{kinesthetic} dimension is not de-emphasized; it is amputated. The hands still hold a controller that now moves nothing, and the involvement model has no category for that condition, because it is not a redistribution of attention at all. It is an inversion of the player's position in the chain: the one who acts becomes the one who is acted upon, the agent becomes the patient. Aristotle gives us the precise concept the model lacks. Being-affected has two modes --- the \textit{paschein} that preserves and fulfills the perceiver, and the \textit{paschein} that extinguishes (\textit{De Anima} 417b2--5). Ordinary perception, ordinary play, is the first kind. What begins at 02:31.67 is the second, and this distinction is the hinge of the entire analysis: everything the event does from here --- the blackouts, the taunts, the itemized hand-back --- is the destructive mode of being-affected, administered through a medium built for the preservative one.
This is where the event outruns the scholarship's agency debate: against the readings on which the game's freedom is revoked outright (Ruffino 347) or illusory from the start (Tuominen 25), and equally against those on which it is genuine within a fence of coded consequence (McEvoy 12--13), the Sonny event presupposes a freedom real enough to be exercised, precisely so that this exercise can be turned back upon the player---agency not foreclosed, not fenced, but leveraged against them to rhetorically amplify the felt consequences of their own action.

The cutscene's own construction enacts this destruction as a sensory architecture, and the frame record lets us specify it exactly. Sonny approaches---``Now come here''---raises something overhead, and at 02:35.67, at the instant of contact, the image hard-cuts to black. The blow is vaguely depicted; the violence is cut away from at the moment of impact. Nine seconds of darkness follow, filled only with Sonny's voice: ``Don't you hate old Sonny now... don't hate him''---the assailant pre-emptively managing the victim's judgment while the victim's eyes are still closed. At 02:44.71 the black resolves, and the frame record fixes the sequence exactly: the distorted vision is not a fading-out but the coming-to. The player surfaces into a single semi-conscious shot from the floor looking up, the frame vignetted and color-shifted as though perception itself were bleeding, and what it shows is little more than a bare torso and a hand working at a pair of suspenders---the visual counterpart of the belt-buckle sounds---while Sonny leans over the camera: ``Oh, you struggled... and you lost, but it was quite a tussle, I tell you. Quite a tussle, my pet.'' Consider what these lines do. They narrate, in the past tense, a struggle the player never performed and never saw; the player's agency is handed back as something already spent and already beaten, its authorship rewritten while its exercise was hidden. ``My pet'' completes the inversion lexically---from person to owned creature, from agent to patient---and then, at 02:56.42, the screen goes black again and stays black for nearly twenty-four seconds while Sonny delivers the event's thesis: ``See, friendship ain't so tough... and neither is you.'' Somewhere in that darkness a body is dumped. Measured against the recording, the arithmetic is stark: of the roughly seventy seconds during which control is seized, nearly half---some thirty-three seconds---is black screen. Where ordinary play renders a world, this sequence renders the failing of the faculty that renders. That faculty, on the account this project has developed, is \textit{phantasia}: it is \textit{phantasia} that makes the world present across the screen's thin cues, and the vignetted, color-shifted shot from the floor hands it a new and terrible assignment---no longer to constitute a world, but to make present, from the inside, what the collapse of that making-present is like. The camera work that remains seats the player inside the victim's dimming perception and then turns even that off. The game does not show the player a victim. It makes the player one.

This withholding is where \textit{phantasia} must be tied in, because the blackout is not an absence of rhetoric but its concentration. Recall that \textit{phantasia} is the faculty that supplies what is missing when the object cannot be perceived, mediating direct perception with our experiential backlog; recall too that the images it renders are drawn from what each soul has sedimented. When the world withholds depiction, then, it does not withhold the event---it hands the rendering over. It can do so because so little is definitively given: for most of the seized span the player has nothing to work from but lexical and auditory cues---``my pet,'' the metallic clanking, the croon of the voice---and cues of that kind fix almost nothing while suggesting almost everything. The door is left wide open, and \textit{phantasia} walks through it. Into those thirty-three seconds the game pours only sounds and taunts, and the player's \textit{phantasia}, conscripted, builds the scene the screen refuses to build, from materials the player alone supplies. This is why the event's content is genuinely different for differently attuned players, and the mechanism, once visible, lets the comparison be stated exactly. Take Kyle, the seventeen-year-old from Wales, as one point of comparison, and for the other a hypothetical young child with no concept of rape, sexual assault, or perhaps sex itself. (Though I must note I highly discourage any parent from allowing young children to play this game.) Both are in the same situational circumstance and exposed to the same stimuli. For the adult, whose store of cultural knowledge includes the schema of sexual assault, the belt sounds, the ``my pet,'' the drained body, and the bow-legged walk assemble into rape; the metallic sounds a child might hear as rifling through Arthur's pockets assemble instead into robbery, an opinion the game's own narrative frame---a gang story full of theft---renders entirely viable. Here the frame record adds a detail that strengthens the point: the game \textit{funds} the child's reading. The one on-screen fact the aftermath grants is a red debit ticker---exactly one dollar lost, \$1097.35 to \$1096.35---the world's own receipt for the innocent interpretation. The same withheld event yields two different events, each rendered by a different soul from identical prompts. The design's precision here should be stated plainly. No frame of the event depicts an assault, and no line names one; the game's own data underdetermine what happened. Yet for a player whose sedimented experience includes the schema of sexual assault, every cue converges on it. The event's force therefore does not live in its data; it lives in the \textit{phantasmata} the players themselves are conscripted into making---which is why showing less here moves more: each player authors the event, from their own materials, at the design's instruction.

The world, moreover, quietly corroborates the darker rendering for any player who goes looking. Before the event, Sonny's cabin cannot be entered except at his invitation: the back door is sealed, and if the player aims a weapon at him he flees inside and locks a door that no amount of force will open---a barrier that becomes breakable only after the event has run. Inside, the set dressing reads like an evidence room. A child's doll---a girl's, by its bonnet and coloring---sits atop the shelves. Bottles of tooth tonic, a preparation that in this period would likely have run on alcohol and opium, litter the cabin inside and out, and beside them stands a bottle of ``Maryan P. Horsford's QUICK NEUTRALIZING CORDIAL,'' promising ``relief to man or beast'' for ``rheumatism and cholera relief''---a strong sedative, and perhaps a second reason Arthur surfaces dazed. Shackles are mounted on the wall by the bed, at both the head and the foot of the frame. Back at camp, the game supplies its most pointed corroboration in dialogue: Bill Williamson approaches Arthur---``I met an interesting feller in the swamp, real interesting. / Did you? / Sure... he seemed to know all about you. I mean... \textit{all} about you. / Get out of here!''---his tone and diction saying what no line states. (Williamson, dishonorably discharged from the army for attempted murder and deviancy, is elsewhere hinted to be gay; how or why he knows Sonny is never disclosed.) None of this is forced on the player; all of it is left where a seeking \textit{hexis} will find it. The design that refused to depict the event salted the world with the means to conclude it.
The event's single prior scholarly treatment reaches the same conclusion from the other side. Moser — the only critic to have written on this scene — treats it as her limit case of the player/avatar asymmetry: a player who cannot bear what happened can ``load a previous save so it would be as if the event never happened'' (Moser 16), and the game will oblige, but ``the lifeworld of the player will still remember the event, even if it never technically happened, and Arthur does not remember it'' (16). Her observation is made in passing, in service of a lifeworld argument, and it is the present point in embryo: an event that can be deleted from the avatar's world but not from the player's does not principally live in the avatar's world. The erasure option proves the address; what the save-file cannot reach is where the event happened. Nor does the loaded save wipe the slate it claims to wipe: every experience leaves its resonant \textit{kinēsis}, and what has been undergone has already shaped the dispositional ground of the one who underwent it. No file operation reaches that ground. The player who loads the previous save carries the event forward as surely as the player who rides on. A depicted assault would have been Rockstar's image, held at the distance of any spectacle. The undepicted assault is the player's own image, and it is held with the authority of first-person things.

The hand-back completes the design's statement, and it comes in the only language a system speaks: a ledger. The player wakes as Arthur slumped in the marsh grass---``Oh my Lord...'' is the only line the script gives him, an exclamation that registers \textit{that} something happened while refusing to say \textit{what}, preserving in the avatar's mouth the same indeterminacy the blackout preserved on screen. For three and a half seconds after rising (03:31.0--03:34.6 in the recording) the avatar walks bent and shuffling before straightening---the body testifying, in the only register left, to what the screen would not show. Then the interface returns piece by piece, and the sequence of its return reads like a receipt: the money counter first, ticking down the single lost dollar; the sound notifying that all of Arthur's cores have been drained; and, then, of all things, a tutorial toast explaining that resting will refill the player's cores; then the minimap; then the cores themselves---health, stamina, \textit{and} Dead Eye, all drained to red.\footnote{The health core represents the player's total health, the stamina core the ability to sprint, and the Dead Eye core the ability to slow down time in combat. All three are fully depleted after this event; the stamina and health depletions are additionally announced by system messages (``Your Stamina Core is empty''; ``Your Health Core is empty'') in the minutes that follow.} The dollar prices the event at nothing, and thereby announces that the taking was never for money; the cores translate the harm into the machine's own vocabulary. While one would expect health and stamina to be drained after such an event, the Dead Eye drain says something more: Dead Eye is not a bodily quantity but a capacity of mastery---the player's trained power to slow the world and act with superhuman composure, the mechanical correlate of everything hours of play have made of them. Its draining announces that the assault reached past the body's resources into the player's cultivated dispositions: the harm touched the \textit{hexis} itself, the sedimented substrate of mastery.

One more feature of the aftermath must be weighed, and it is a silence. \textit{RDR2} runs a persistent honor system---the moral machine analyzed throughout this project, quick to reward a greeting or a returned wallet, quick to punish a shot dog---and here it says nothing. No honor notification fires at the event. More strikingly, none fires afterward: once the player has recovered the avatar they are free to return and enact vengeance or to ride away (\textit{ludic involvement}), and my own recording carries the whole reprisal---the ride back through the fog, Sonny barricaded behind his door (``Not you again... I don't want you no more!''), the lock that before the event no force could open now giving way, the execution---and no notification appears in either direction. The kill is honor-neutral; I have confirmed this directly.
Such tuning has precedent, and the scholarship has logged it without pressing it. Tuominen notices that ``By not punishing a player for killing the Klansmen, the game's rules (the rules of the world) seem to imply that killing racists is okay'' (Tuominen 53); Bagnoli finds the same exception run further — ``killing members of this organized group will grant you honor, which definitely says a lot about Rockstar's opinion on the KKK'' (Bagnoli 72) — and a second beneficiary in the wandering eugenicist, who can be struck down ``without suffering from any consequences'', a marked exemption from what Bagnoli notes is otherwise a general debit for killed citizens (70). Reward, exemption, silence: the honor machine's coverage is a tuned spectrum, its verdicts legible as much in what it declines to score as in what it scores, and the Sonny kill extends that spectrum to its terminal case. The one system built to attribute motive and pass public verdicts withdraws entirely from the field, and the withdrawal mirrors the visual design exactly: just as the screen went dark and made the player author the scene, the score goes silent and makes the player author the sentence. Was the execution justice, cruelty, or grief? The machine will not say, and so the judgment---like the event---happens \textit{in} the player rather than in front of them. Run the return through the chain and every stage is the player's own. Its cognitive ground is the \textit{doxa} the blackout forced each player to author---for the attuned player, \textit{Arthur was raped}---and the emotion that ratified opinion discloses is anger, or disgust, or grief, each a different attunement of the same wound. The emotion then does what emotion does throughout this method: it discloses some possibilities for action while masking others. Those in a fury may not even conceive of letting the swamp man live; their deliberation runs only across the \textit{how}---some players simply shoot the offender, while others have hog-tied him and fed him to the alligators. The \textit{orekton} of the discrete event is vengeance, Sonny dead, and its cause, in the \textit{Rhetorica}'s catalogue, is nearer \textit{thymos} than reasoning: retaliatory anger for a suffered wrong. \textit{Phantasia} holds the target charged with everything the blackout made of him---the croon, the clanking, ``my pet''---the \textit{doxa} stands already ratified, and at the A₃→A₄ transition the act completes with whatever instrument the player's anger has selected. Nothing in the game's economy motivates any of it---one dollar is an inconsequential amount at this stage---and that is precisely the finding: the vengeance is sourced not by any system of the game but by the wound, by what the event deposited in the player's own dispositional ground. Either way the player must live with a decision the world refuses to score. At every register, the event is designed to happen to the player, not just to the character through which they act.
Ruffino documents the same shape at the game's other wound: ``The desire to save Arthur and the awareness of its impossibility co-exist'' (Ruffino 347), his cure-seekers returning and returning to an impossibility the design will not resolve. The player riding back through the fog to Sonny's door is moved not by anything the world still offers but by what it took.

The reception evidence shows how wide that deposit's range runs, and the mechanism explains the range rather than merely reporting it. As Amelia Tate's article reveals, the felt response (\textit{affective involvement}) spanned a spectrum: at one end players responded with an indifferent ``Gayyy''; at the other, some who had personally experienced sexual assault were forced to relive those experiences, the fear, pain, and anxiety returning with them; and between these extremes, players like Kyle reported a profound disorientation---feeling ``super uncomfortable''---that ultimately transformed their perspective:

\begin{adjustwidth}{1in}{0in}
```Yes, it granted me a new perspective on rape,' says Kyle Rees, 17, from Wales.... `Before experiencing this scene I never really gave much thought about rape. I've always been a bit negligent or closed-minded about it,' he admits. `I've always thought to myself, How can you even allow someone to rape you? Just claw at the dude's eyes or punch him in the nuts.' ''
\end{adjustwidth}

These various attunements, from indifference to discomfiture, demonstrate how attunement expresses the countless modalities of responsiveness to the feedback between game and player---but the mechanism already established, the design's conscription of \textit{phantasia}, its compelling of each player's own faculty to author the scene from that player's sedimented materials, lets us say why the spectrum has the shape it does. Each response is the signature of a different rendering: the resonant ground of the indifferent player supplied a scene that could be shrugged off; the survivor's ground supplied their own past; Kyle's ground supplied enough to build the event but not enough to defend against it. Kyle's case shows what the event was, I believe, built to do. His standing opinion about victims---``just claw at the dude's eyes''---was not amended by counter-argument; it was supplanted by being forcibly situated as the victim of a horrific event, because even someone as strong and combat-seasoned as the game's protagonist is helpless when totally blindsided.

Put back together, the sequence describes a single arc, and it is the arc this project's methodology names rhetorical incorporation---run in reverse. Everything before the doorway is incorporation achieved: a player at home in a mastered world, moving through internalized controls, one with the avatar acted-as and the world acted-in, ``substantially one'' while remaining ``a distinct locus,'' in Burke's formula for consubstantiality.
The union has even a linguistic signature: players, Moser observes, ``they recognize they are not Arthur, but they are also not not Arthur'' (Moser 30), their speech sliding between what Arthur did and what \textit{I} did. The event does not simply break this union, for breaking it would only eject the player into spectatorship; instead it converts the union, holding the player inside the avatar---the floor-level camera guarantees it---while everything that made the union a union of \textit{agency} is amputated, so that oneness itself becomes the medium of helplessness. That conversion is what the players in Tate's article are describing: what marked them was not an image, for there was none to see. This is the exact structure of Sarah Ahmed's discomfort. For, as she argues, comfort is the feeling of inhabiting spaces that ``\textit{extend [our] shape}'', and ``we might only notice comfort as an affect when we lose it'' (158); by removing control after the threshold, the game systematically forces the player to accept their own helplessness, disrupting the sense of fittingness, of belonging to and commanding the game world, into discomfiture. Because incorporation is porous rather than sealed---because what sediments in play crosses back into the player's wider being---the conversion travels. The vengeance so many players enact is its first fruit, sourced by the wound alone; Kyle's changed understanding of victimhood is its second and deeper fruit, carried out of the game into judgment, speech, and conduct in the physically real.
Moser's thesis stakes itself on transfer — ``If this is possible with virtual characters, then it is all the more possible in reality with real people'' (Moser 39) — while conceding that the continuation of her account would require quantitative data on how players portray their Arthur and what the portrayal does to their empathy (40). The Tate testimony is that evidence in qualitative form: transfer observed in the wild, with the mechanism — conscripted \textit{phantasia}, undergone helplessness — visible in a way no instrument could catch.

Rockstar has a history of embedding social commentary in its games' procedures, and this history is one component of the ambient array within which the event is interpreted. As Ian Bogost has observed, \textit{Grand Theft Auto: San Andreas} (2004) requires its protagonist to eat, offers mostly fast food, and locates his headquarters in a low-income area---a procedural critique, Bogost argues, of the affordability of poor nutrition in such communities (``Videogames and Ideological Frames,'' 176).
The \textit{RDR2} scholarship extends the observation to this game directly. Bagnoli announces at his outset that ``The questions of feminism, race, and animal condition will be tackled by the game'' (Bagnoli 7) and reads the game's processes as the delivery mechanism of a contemporary worldview (40); McEvoy goes further and credits the design with prophylaxis — ``Rockstar utilizes the narrative structure of the game to comment on discriminatory ideologies commonly held at that moment in the U.S. and, by adding consequences to discriminatory expressions made by the player, stifles the expansion of those ideologies held today'' (McEvoy 3) — though the claim is asserted rather than evidenced: her reader-response frame gathers no actual reader responses. Whether Rockstar can mount such critique at all is itself contested: Tuominen reads the series as radically progressive under a mainstream skin (60), Wright as a critique undercut by the very frontier branding it profits from (par.\ 9), and the Sonny event is a sharp test for both positions. In light of this, some have stipulated that the rape scene was a developer's attempt to challenge victim-blaming culture and raise awareness in a community with a notoriously misogynistic history; as Kyle puts it, ``I want to believe Rockstar made this scene in order to get people to understand how serious rape cases are.'' Whatever the intention, the procedural form of the commentary is now legible in full. The event teaches what victim-blaming gets wrong not by depicting a blameless victim but by making the player one: entry freely chosen, every warning available, every precaution possible---and none of it mattering, because the seizure was engineered to follow consent's form wherever it led. A player who believes that victims could have fought back is given a protagonist who canonically \textit{did} fight back (``you struggled... quite a tussle'') and lost anyway, off screen, where no skill could be exercised and no deadeye slowed the world. To reiterate Rickert's claim, ``ambience is rhetorical to the extent that we become a furtherance of or site for rhetorical activity'' (120). For Kyle this event catalyzed a new perspective---an effect not isolated to his conduct in the game world but carried over into the physically real, where he may now contest victim-blaming when he encounters it in games and gaming communities. With his dispositional ground affected and an ethical \textit{hexis} sedimented, Kyle has become a site for further rhetorical activity. No single element is responsible for this effect; the agency of the event is ambiently dispersed between hardware, programmers, artists, NPC, cutscene, the structural and systemic properties of the game world, the honor machine's silence, the player's sedimented past, and more, all of which converge in a \textit{kairotic} moment with the power to reattune a person's being-in-the-world. That transformation, ambiently diffused, ripples outward in future rhetorical possibilities that can reach people who will never play the game.

By emplacing the player, as a participant, in a virtual rhetorical ecology which actualizes mutual affectability between player and world, a clever developer can leverage ambient rhetoric to reveal the world in a desired way---to reattune or transform the player's mode of being-in-the-world. What this event demonstrates is the full anatomy of that power in its severest register. A virtual world spent its length cultivating a disposition, used that disposition as a trigger, let the chain of desire run as it always runs and seized its completion one step past a freely made choice, inverted its inhabitant from agent to patient in the destructive mode of being-affected, withheld the event's content so that each player's \textit{phantasia} was conscripted to author it, withheld the verdict so that each player's \textit{doxa} was left to judge it, and returned agency as a drained and limping body. This is what I call ontological violence: harm that, once the player has been rhetorically incorporated and so come to \textit{be-in-the-virtual-world}, does not remain at the level of the avatar but passes through that virtual being-there to reconfigure the player's own \textit{being-in-the-world}. It is rhetorically ontologically transformative for exactly the reasons developed throughout this project, since only an event that reaches the player's \textit{for-the-sake-of-which}, through a relation of mutual affectability rather than of subject to depicted object, can produce a change that walks out of the game with the player once they turn off the program. This is a violence done at the level of the \textit{how} of being: a subversion of the person's actualization of desire in which a world uses a player's own cultivated disposition as the trigger, inverts the player from agent to patient in the destructive mode of being-affected, withholds the event's content so that the player's \textit{phantasia} must author it, withholds the verdict so that the player's \textit{doxa} must judge it, and returns agency only afterwards as a reminder of just how usurpable---how fragile---it truly is, transmuting an experience into the sedimentation of an altered \textit{being-in-the-world}. The phrase is deliberately double, for the violence is ontological in what it targets---the player's mode of \textit{being-in-the-world}---and ontological in what it makes possible---since the same affectability that lets the destruction occur lets the new understanding blossom. An event of this design can thereby reattune a person's \textit{being-in-the-world}, the very dispositional ground from which they will encounter the world thereafter; that is its danger, and its demonstration. A virtual world transformed a seventeen-year-old's grasp of what it means to be helpless not by depicting helplessness, but by compelling him to live it as his own mode of \textit{being-in-the-world}, in which helplessness unconceals itself from itself. This is achieved by designing a virtual world that takes into account and leverages the fundamental nature by which we perceive, understand, and act; in other words, by hacking the human soul.

\end{document}
